
\documentclass{article}
\usepackage{amsmath}

%%%%%%%%%%%%%%%%%%%%%%%%%%%%%%%%%%%%%%%%%%%%%%%%%%%%%%%%%%%%%%%%%%%%%%%%%%%%%%%%%%%%%%%%%%%%%%%%%%%%%%%%%%%%%%%%%%%%%%%%%%%%%%%%%%%%%%%%%%%%%%%%%%%%%%%%%%%%%%%%%%%%%%%%%%%%%%%%%%%%%%%%%%%%%%%%%%%%%%%%%%%%%%%%%%%%%%%%%%%%%%%%%%%%
%TCIDATA{OutputFilter=LATEX.DLL}
%TCIDATA{Version=5.50.0.2960}
%TCIDATA{<META NAME="SaveForMode" CONTENT="1">}
%TCIDATA{BibliographyScheme=Manual}
%TCIDATA{Created=Wednesday, October 21, 2026 11:36:13}
%TCIDATA{LastRevised=Wednesday, October 21, 2026 11:48:35}
%TCIDATA{<META NAME="GraphicsSave" CONTENT="32">}
%TCIDATA{<META NAME="DocumentShell" CONTENT="Scientific Notebook\Blank Document">}
%TCIDATA{CSTFile=Math with theorems suppressed.cst}
%TCIDATA{PageSetup=72,72,72,72,0}
%TCIDATA{AllPages=
%F=36,\PARA{038<p type="texpara" tag="Body Text" >\hfill \thepage}
%}


\newtheorem{theorem}{Theorem}
\newtheorem{acknowledgement}[theorem]{Acknowledgement}
\newtheorem{algorithm}[theorem]{Algorithm}
\newtheorem{axiom}[theorem]{Axiom}
\newtheorem{case}[theorem]{Case}
\newtheorem{claim}[theorem]{Claim}
\newtheorem{conclusion}[theorem]{Conclusion}
\newtheorem{condition}[theorem]{Condition}
\newtheorem{conjecture}[theorem]{Conjecture}
\newtheorem{corollary}[theorem]{Corollary}
\newtheorem{criterion}[theorem]{Criterion}
\newtheorem{definition}[theorem]{Definition}
\newtheorem{example}[theorem]{Example}
\newtheorem{exercise}[theorem]{Exercise}
\newtheorem{lemma}[theorem]{Lemma}
\newtheorem{notation}[theorem]{Notation}
\newtheorem{problem}[theorem]{Problem}
\newtheorem{proposition}[theorem]{Proposition}
\newtheorem{remark}[theorem]{Remark}
\newtheorem{solution}[theorem]{Solution}
\newtheorem{summary}[theorem]{Summary}
\newenvironment{proof}[1][Proof]{\noindent\textbf{#1.} }{\ \rule{0.5em}{0.5em}}
\input{tcilatex}

\begin{document}


\section{Nature of Light and Reflection 3.1.1, 3.1.2}

In the science fiction movies inter-galactic star ships blast each other
with laser cannons. The laser beams streak across the screen. This is
dramatic, but not realistic. For us to see the light path, some of the light
must leave the path and get to our eyes. The light must either travel
directly to our eyes from the source, or it must bounce off of something
along its way. We can make a laser beam visible by providing dust for the
light to bounce off of so it travels to our eyes. But unless that galaxy far
far away is really dusty, Hollywood doesn't understand light waves very
well. Let's start our study of geometric optics with a few comments on what
light is. Then let's study reflection, the bouncing of light off of an
object (like dust, or a mirror).

%TCIMACRO{%
%\TeXButton{Fundamental Concepts}{\vspace{0.10in}\hspace{-0.37in}{\Large Fundamental Concepts\vspace{0.2in}}}}%
%BeginExpansion
\vspace{0.10in}\hspace{-0.37in}{\Large Fundamental Concepts\vspace{0.2in}}%
%EndExpansion

\begin{itemize}
\item Light is a wave in the electromagnetic field

\item Light travels fast, with a speed of $c=3\times 10^{8}\unit{m}/\unit{s}$

\item Law of Reflection
\end{itemize}

\subsection{What is Light?}

Before the 19th century (1800's) light was assumed to be a stream of
particles. Newton was one of the chief proponents of this theory. The
\textquotedblleft corpuscles\textquotedblright\ of light were thought to
bounce like little balls. The theory was able to explain reflection of light
from mirrors and other objects and therefore explain vision. In 1678 Huygens
showed that wave theory could also explain reflection and vision.

In 1801 Thomas Young demonstrated that light had attributes that were best
explained by wave theory. We will study Young's experiment later. The crux
of his experiment was to show that light displayed constructive and
destructive interference--clearly a wave phenomena! The theory of the nature
of light took a dramatic shift.

In 1805 Joseph Smith was born in Sharon, Vermont.

In September of 1832 Joseph Smith received a revelation that said in part :

\begin{quotation}
For the word of the Lord is truth, and whatsoever is truth is light, and
whatsoever is light is Spirit, even the Spirit of Jesus Christ. And the
Spirit giveth light to every man that cometh into the world; and the Spirit
enlighteneth every man through the world, that hearkeneth to the voice of
the Spirit. (D\&C 84:45-46)
\end{quotation}

In December of 1832 Joseph Smith received another revelation that says in
part:

\begin{quotation}
This Comforter is the promise which I give unto you of eternal life, even
the glory of the celestial kingdom; which glory is that of the church of the
Firstborn, even of God, the holiest of all, through Jesus Christ his
Son---He that ascended up on high, as also he descended below all things, in
that he comprehended all things, that he might be in all and through all
things, the light of truth; which truth shineth. This is the light of
Christ. As also he is in the sun, and the light of the sun, and the power
thereof by which it was made. As also he is in the moon, and is the light of
the moon, and the power thereof by which it was made; as also the light of
the stars, and the power thereof by which they were made; and the earth
also, and the power thereof, even the earth upon which you stand. And the
light which shineth, which giveth you light, is through him who enlighteneth
your eyes, which is the same light that quickeneth your understandings;
which light proceedeth forth from the presence of God to fill the immensity
of space---the light which is in all things, which giveth life to all
things, which is the law by which all things are governed, even the power of
God who sitteth upon his throne, who is in the bosom of eternity, who is in
the midst of all things. (D\&C 88:5-12)
\end{quotation}

Light, even non-metaphorical physical light, seems to be of interest to
members of the Restored Church.

In 1847 the saints entered the Salt Lake Valley.

In 1873 James Clark Maxwell published his findings that light is an
electromagnetic wave (something we will try to show in Principles of Physics
III!). That is, the \textquotedblleft group of objects\textquotedblright\
that form a \textquotedblleft medium\textquotedblright\ for light are not
made out of matter. But are part of what we call the electromagnetic field.
This field is something, but not made out of matter.

Planck's work in quantization theory (1900) was used by Einstein In 1905 to
give an explantation of the photoelectric effect that again made light look
like a particle.

Current theory allows light to exhibit the characteristics of a wave in some
situations and like a particle in others. We will study both before the end
of our optics unit.

The results of Einstein's work give us the concept of a \emph{photon} or a
quantized unit of radiant energy. Each \textquotedblleft piece of
light\textquotedblright\ or photon has energy 
\begin{equation}
E=hf
\end{equation}%
where $f$ is the frequency of the light and $h$ is a constant 
\begin{equation}
h=6.63\times 10^{-34}\unit{J}\unit{s}
\end{equation}

The nature of light is fascinating and useful both in physical and religious
areas of thought.

\subsection{Measurements of the Speed of Light}

One of the great foundations of modern physical theory is that the speed of
light is constant in a vacuum. Galileo first tried to measure the speed of
light. He used two towers in town and placed a lantern and an assistant on
each tower. The lanterns had shades. The plan was for one assistant to
remove his shade, and then for the assistant on the other tower to remove
his shade as soon as he saw the light from the first lantern. Back at the
first tower, the first assistant would use a clock to determine the time
difference between when the first lantern was un-shaded, and when they saw
the light from the second tower. The light would have traveled twice the
inter-tower distance. Dividing that distance by the time would give the
speed of light. You can probably guess that this did not work. Light travels
very quickly. The clocks of Galileo's day could not measure such a small
time difference. Ole R\o mer was the first to succeed in measuring the speed
of light.

\subsubsection{R\o mer's Measurement of the speed of light}

\FRAME{dhFU}{1.1761in}{2.0626in}{0pt}{\Qcb{A diagram illustrating R\o mer's
determination of the speed of light. Point A is the Sun, piont B is Jupiter.
Point C is the immersion of Io into Jupiter's shadow at the start of an
eclipse}}{}{Figure}{\special{language "Scientific Word";type
"GRAPHIC";maintain-aspect-ratio TRUE;display "USEDEF";valid_file "T";width
1.1761in;height 2.0626in;depth 0pt;original-width 1.7252in;original-height
3.0468in;cropleft "0";croptop "1";cropright "1";cropbottom "0";tempfilename
'TN9QPD3R.wmf';tempfile-properties "XPR";}}R\o mer performed his measurement
in 1675, 269 years before digital devices existed!. He used the orbital
period of Io, a moon of Jupiter, as Jupiter orbited around the sun. He first
measured the period of Io's rotation about Jupiter, then he predicted an
eclipse of Io three months later. But he found his calculation was off by $%
600\unit{s}.$ After careful thought, he realized that the Earth had moved in
its orbit, and that the light had to travel the extra distance due to the
Earth's new position. Given R\o mer's best estimate for the orbital radius
of the earth and his time difference, R\o mer arrived at a estimate of $%
c=2.3\times 10^{8}\frac{\unit{m}}{\unit{s}}.$ Our current value is $%
c=2.99792458\times 10^{8}\frac{\unit{m}}{\unit{s}}.$ Amazing for 1675!

\subsubsection{Fizeau's Measurement of the speed of light}

\FRAME{dtbpF}{2.5659in}{1.6881in}{0pt}{}{}{Figure}{\special{language
"Scientific Word";type "GRAPHIC";maintain-aspect-ratio TRUE;display
"USEDEF";valid_file "T";width 2.5659in;height 1.6881in;depth
0pt;original-width 2.5253in;original-height 1.6518in;cropleft "0";croptop
"1";cropright "1";cropbottom "0";tempfilename
'TN9QPD3S.wmf';tempfile-properties "XPR";}}Hippolyte Fizeau measured the
speed of light in 1849 using the apparatus drawn in the figure above. He
used a toothed wheel and a mirror and a beam of light. The light passed
through the open space in the wheel's teeth as the wheel rotated. Then was
reflected by the mirror. The speed would be%
\[
v=\frac{\Delta x}{\Delta t}
\]%
We just need $\Delta x$ and $\Delta t.$

It is easy to see that 
\[
\Delta x=2d
\]%
because the light travels twice the distance to the mirror ($d$) and back.
If the wheel turned just at the right angular speed, then the reflected
light would hit the next tooth and be blocked. Think of angular speed $%
\omega $ (not angular frequency this time)%
\[
\omega =\frac{\Delta \theta }{\Delta t}
\]%
So, the time difference for a tooth to move would be 
\[
\Delta t=\frac{\Delta \theta }{\omega }
\]%
We can find $\Delta \theta $ by taking $2\pi $ and dividing by the number of
teeth on the wheel.%
\[
\Delta \theta =\frac{2\pi }{N_{teeth}}
\]

Then the speed of light must be 
\begin{eqnarray*}
c &=&v=\frac{2d}{\frac{\Delta \theta }{\omega }} \\
&=&\frac{2d\omega }{\Delta \theta } \\
&=&\frac{2d\omega N_{teeth}}{2\pi } \\
&=&\frac{d\omega N_{teeth}}{\pi }
\end{eqnarray*}

Then if we have $720$ teeth and $\omega $ is measured to be $d=7500\unit{m}$%
\begin{eqnarray*}
c &=&\frac{\left( 7500\unit{m}\right) \left( 172.\,\allowbreak 79\unit{Hz}%
\right) \left( 720\right) }{\pi } \\
&=&2.\,\allowbreak 97\times 10^{8}\frac{\unit{m}}{\unit{s}}
\end{eqnarray*}%
which is Fizeau's number and it is pretty good!

Modern measurements are performed in very much the same way that Fizeau did
his calculation. The current value is 
\begin{equation}
c=2.99792458\times 10^{8}\frac{\unit{m}}{\unit{s}}
\end{equation}

\subsubsection{Faster than light}

The speed of light in a vacuum is constant, but in matter the average speed
of light changes.

We will study this in detail when we look at refraction. But for now, a
dramatic example is Cherenkov radiation. It is an eerie blue glow around the
core of nuclear reactors. It occurs when electrons are accelerated past the
speed of light in the water surrounding the core. The electrons emit light
and the light waves form a Doppler cone or a light-sonic boom! The result is
the blue glow.\FRAME{dtbpFUw}{2.9457in}{2.2024in}{0pt}{\Qcb{{\protect\small %
\ Cherenkov radiation (United States Department of Energy, image in the
public domain)}}}{}{Figure}{\special{language "Scientific Word";type
"GRAPHIC";maintain-aspect-ratio TRUE;display "USEDEF";valid_file "T";width
2.9457in;height 2.2024in;depth 0pt;original-width 6.2595in;original-height
4.6717in;cropleft "0";croptop "1";cropright "1";cropbottom "0";tempfilename
'TN9QPD3T.wmf';tempfile-properties "XPR";}}This does bring up a problem in
terminology. What does the word \textquotedblleft medium\textquotedblright\
mean? We have used it to mean the substance through which a wave travels.
This substance must have the property of transferring energy between it's
parts, like the coils of a spring can transfer energy to each other, or like
air molecules can transfer energy by collision. For light, the wave medium
is the electromagnetic field. This field can store and transfer energy but
it isn't made of atoms or parts of mater. This may be hard to imagine, but
picture space filled with this non-material substance that can stretch and
hold energy. This is our electromagnetic field. We will study the
electromagnetic field in Principles of Physics III (PH220). So for now, we
just need a mental picture of the field having waves.

But many books on physics call materials like glass a \textquotedblleft
medium\textquotedblright\ through which light travels. The water in our
nuclear reactor example is a \textquotedblleft medium\textquotedblright\ in
this sense. Are glass and water wave mediums for light? The answer is an
emphatic NO! Light does not need any matter to form it's wave. The wave
medium is the electromagnetic field. So we will have to keep this in mind as
we allow light to travel through matter. I will try to say that the light
enters a new \textquotedblleft material\textquotedblright\ to describe
something like light entering a piece of glass. But some books may call the
matter a \textquotedblleft medium.\textquotedblright\ You must remember that
by this they don't mean the wave medium.

Light can't be a wave in some type of matter. We know this because light
travels through empty space. But in Principles of Physics III we will find
that space isn't exactly empty. It has an electromagnetic field in it. And
it turns out that light seems to be a wave in this electromagnetic field. It
will take us a while to fully understand this concept but don't worry.
Physicists knew that light was a wave for almost 80 years before the
electric field was shown to be the medium. We can do a lot just knowing
light behaves like a wave.

This makes the light we see just one small part of a whole class of waves
that are possible in this electromagnetic field medium. Radio waves, and
microwaves, and x-rays are all just different types of electromagnetic
waves. The next figure shows where all of these electromagnetic waves fit
ordered by wavelength (and frequency).\FRAME{dtbpFU}{5.2642in}{0.3991in}{0pt%
}{\Qcb{{\protect\small Image Courtesy NASA}}}{}{Figure}{\special{language
"Scientific Word";type "GRAPHIC";maintain-aspect-ratio TRUE;display
"USEDEF";valid_file "T";width 5.2642in;height 0.3991in;depth
0pt;original-width 7.2032in;original-height 0.518in;cropleft "0";croptop
"1";cropright "1";cropbottom "0";tempfilename
'TN9QPD3U.wmf';tempfile-properties "XPR";}}

There is something very unique about this electromagnetic field medium. The
waves in this medium travel at a constant speed no matter what frame of
reference we are in. This fact lead to the formation of the Special Theory
of Relativity and the famous equation 
\[
E=mc^{2}
\]%
where $c$ is this speed of light%
\[
c=299792458\frac{\unit{m}}{\unit{s}}
\]

Light does slow down when it enters a material, like glass, or even air.
This is not really because it moves slower. what happens is that light is
absorbed by the electrons in the atoms of the material substance. The
electron temporarily takes up all the energy from a bit of the light wave.
But only temporarily. It eventually has to give up the energy and a light
wave is reformed coming from the atoms. There is really a superposition of
what is left of the original light wave, and the new light wave generated by
the atoms. And the resultant wave moves slower. Think of our last lecture on
non-sinusoidal waves. The resultant wave is made of two sinusoidal waves
both traveling at the speed of light. But the combination of the two waves
will also be a wave, and this resultant wave will appear to travel at a
slower speed. How much slower depends on how long the electrons in the atoms
can hang on to the light and the geometry of the situation. Each substance
is different.

We can devise a way to express how much slower light will appear to go in a
substance using the ratio%
\[
\frac{c}{v}
\]%
the ratio of the speed of light, $c,$ to the average speed in the substance, 
$v.$ This ratio is so useful that we give it a name, the \emph{index of
refraction}.%
\[
n=\frac{c}{v}
\]

Let's look at some a simple interaction of light waves with a material, but
let's not look at the wave and atomic re-emition level. Let's treat our
combined wave as just one resultant wave that travels in a specific
direction. We can call this wave a \emph{ray} (as in a ray of light!).

\subsection{Reflection}

Using the ray approximation we wish to find what happens when a bundle of
rays reaches a boundary between materials. If the material boundary is very
smooth, then the rays are reflected (bounce off) in a uniform way. This is
called \emph{specular }reflection\FRAME{dhF}{2.3895in}{2.1344in}{0pt}{}{}{%
Figure}{\special{language "Scientific Word";type
"GRAPHIC";maintain-aspect-ratio TRUE;display "USEDEF";valid_file "T";width
2.3895in;height 2.1344in;depth 0pt;original-width 2.3497in;original-height
2.0954in;cropleft "0";croptop "1";cropright "1";cropbottom "0";tempfilename
'TN9QPD3V.wmf';tempfile-properties "XPR";}}If the surface is not smooth,
then something different happens. The rays are reflected, but they are
reflected randomly

\FRAME{dtbpF}{2.3886in}{1.8498in}{0pt}{}{}{Figure}{\special{language
"Scientific Word";type "GRAPHIC";maintain-aspect-ratio TRUE;display
"USEDEF";valid_file "T";width 2.3886in;height 1.8498in;depth
0pt;original-width 3.1211in;original-height 2.4102in;cropleft "0";croptop
"1";cropright "1";cropbottom "0";tempfilename
'TN9QPD3W.wmf';tempfile-properties "XPR";}}This is called \emph{diffuse}
reflection

This difference can be seen in real life. In the next figure, the surface on
the left is a specular reflector and you can see the reflected light beam.
But the surface on the right photo is diffuse, no reflected beam is seen.%
\FRAME{dtbpF}{4.403in}{1.3553in}{0in}{}{}{Figure}{\special{language
"Scientific Word";type "GRAPHIC";maintain-aspect-ratio TRUE;display
"USEDEF";valid_file "T";width 4.403in;height 1.3553in;depth
0in;original-width 4.4589in;original-height 1.3526in;cropleft "0";croptop
"1";cropright "1";cropbottom "0";tempfilename
'TN9QPD3X.wmf';tempfile-properties "XPR";}}

We said the surface must be smooth for there to be specular reflection. What
does smooth mean? Generally the size of the rough spots must be much smaller
than a wavelength to be considered smooth. So suppose we have a red laser.
How small do the surface variations have to be for the surface to be
considered smooth? The wavelength of a $HeNe$ laser is 
\[
\lambda _{HeNe}=633\unit{nm}
\]%
This is very small. Modern optics for remote sensing are often manufactured
to $1/10$ of a wavelength, which would be $63\unit{nm}.$ Mirrors are smooth
for visible light.

How about a microwave beam of light like your cell phone uses?

\begin{eqnarray*}
c &=&\lambda f \\
\lambda &=&\frac{c}{f}=\frac{3\times 10^{8}\frac{\unit{m}}{\unit{s}}}{1\unit{%
GHz}} \\
&=&\allowbreak 0.3\unit{m}
\end{eqnarray*}%
That is $30\unit{cm}!$ Rough brick walls are smooth for cell phone light! We
can see that we must be careful in our definition of \textquotedblleft
smooth.\textquotedblright

\subsection{Law of reflection}

Experience shows that if we do have a smooth surface, that light bounces
much like a ball. This is why Newton though light was a particle. Suppose we
take a flat surface and we shine a light on it. We have a ray that
approaches at an angle $\theta _{i}$ measured from a line drawn
perpendicular to the flat surface. Then the reflected ray will leave the
surface with an angle $\theta _{r}$ also measured from the same line drawn
perpendicular to the flat surface. In optics, we often measure angles from
lines drawn perpendicular to a surface.

Way back in Principles of Physics I we had forces that were perpendicular to
surfaces. We called them \emph{normal forces} where the word \emph{normal}
meant \textquotedblleft perpendicular to the surface.\textquotedblright\ We
can use the same word, \emph{normal,} to describe our perpendicular lines.
We can say that we measure $\theta _{i}$ and $\theta $ $_{r}$ from a \emph{%
normal line} (or sometimes, we just say that we measure from a \emph{normal}
and leave off the word \textquotedblleft line\textquotedblright ). And we
can write a very simple equation to describe reflection. 
\[
\theta _{r}=\theta _{i}
\]

\FRAME{dhF}{1.8585in}{1.1857in}{0pt}{}{}{Figure}{\special{language
"Scientific Word";type "GRAPHIC";maintain-aspect-ratio TRUE;display
"USEDEF";valid_file "T";width 1.8585in;height 1.1857in;depth
0pt;original-width 1.8213in;original-height 1.1519in;cropleft "0";croptop
"1";cropright "1";cropbottom "0";tempfilename
'TN9QPD3Y.wmf';tempfile-properties "XPR";}}This is called the \emph{law of
reflection}. We will look at how superposition makes this happen later. But
from experiment we can see that this law of reflection works.

\subsubsection{Retroreflection}

Let's take and example\FRAME{dtbpF}{1.881in}{1.7625in}{0in}{}{}{Figure}{%
\special{language "Scientific Word";type "GRAPHIC";maintain-aspect-ratio
TRUE;display "USEDEF";valid_file "T";width 1.881in;height 1.7625in;depth
0in;original-width 1.8438in;original-height 1.7244in;cropleft "0";croptop
"1";cropright "1";cropbottom "0";tempfilename
'TN9QPD3Z.wmf';tempfile-properties "XPR";}}

Let's take our system to be two mirrors set at a right angle. We have a beam
of light incident at angle $\theta _{1}$. By the law of reflection, it must
leave the mirror at $\theta _{2}=\theta _{1}.$ We can see that $\alpha $
must be 
\[
\alpha =90\unit{%
%TCIMACRO{\U{b0}}%
%BeginExpansion
{{}^\circ}%
%EndExpansion
}-\theta _{2}
\]
and it is clear that%
\[
180\unit{%
%TCIMACRO{\U{b0}}%
%BeginExpansion
{{}^\circ}%
%EndExpansion
}=90\unit{%
%TCIMACRO{\U{b0}}%
%BeginExpansion
{{}^\circ}%
%EndExpansion
}+\alpha +\beta
\]%
so then 
\begin{eqnarray*}
\beta &=&90\unit{%
%TCIMACRO{\U{b0}}%
%BeginExpansion
{{}^\circ}%
%EndExpansion
}-\alpha \\
&=&90\unit{%
%TCIMACRO{\U{b0}}%
%BeginExpansion
{{}^\circ}%
%EndExpansion
}-\left( 90\unit{%
%TCIMACRO{\U{b0}}%
%BeginExpansion
{{}^\circ}%
%EndExpansion
}-\theta _{2}\right) \\
&=&\theta _{2}
\end{eqnarray*}%
And then 
\[
\theta _{3}=90\unit{%
%TCIMACRO{\U{b0}}%
%BeginExpansion
{{}^\circ}%
%EndExpansion
}-\beta =90\unit{%
%TCIMACRO{\U{b0}}%
%BeginExpansion
{{}^\circ}%
%EndExpansion
}-\theta _{2}=\alpha
\]
By the law of reflection, $\theta _{3}=\theta _{4}.$ So then 
\[
\theta _{3}=\theta _{4}=\alpha
\]%
Let's find how how much the light turns in this mirror system. That would be

\[
\theta _{1}+\theta _{2}+\theta _{3}+\theta _{4}
\]%
which is 
\begin{eqnarray*}
&&\theta _{1}+\theta _{2}+\theta _{3}+\theta _{4} \\
&=&\theta _{2}+\theta _{2}+\alpha +\alpha \unit{%
%TCIMACRO{\U{b0}}%
%BeginExpansion
{{}^\circ}%
%EndExpansion
}-\theta _{2} \\
&=&\theta _{2}+\theta _{2}+90\unit{%
%TCIMACRO{\U{b0}}%
%BeginExpansion
{{}^\circ}%
%EndExpansion
}-\theta _{2}+90\unit{%
%TCIMACRO{\U{b0}}%
%BeginExpansion
{{}^\circ}%
%EndExpansion
}-\theta _{2} \\
&=&90\unit{%
%TCIMACRO{\U{b0}}%
%BeginExpansion
{{}^\circ}%
%EndExpansion
}+90\unit{%
%TCIMACRO{\U{b0}}%
%BeginExpansion
{{}^\circ}%
%EndExpansion
} \\
&=&180\unit{%
%TCIMACRO{\U{b0}}%
%BeginExpansion
{{}^\circ}%
%EndExpansion
}
\end{eqnarray*}%
or the outgoing ray is sent back toward the source! If we do this in three
dimensions we have a corner cube. Here is a picture of a radar corner cube
array.

\FRAME{dhFU}{2.1926in}{2.7913in}{0pt}{\Qcb{{\protect\small Radar
retroreflector tower located in the center of Yucca Flat dry lake bed. Used
as a radar target by maneuvering aircraft during "inert" contact fusing bomb
drops at Yucca Flat. Sandia National Laboratories conducted the tests on the
lake bed from 1954 to 1956. (Image in the Public Domain in the United States)%
}}}{}{Figure}{\special{language "Scientific Word";type
"GRAPHIC";maintain-aspect-ratio TRUE;display "USEDEF";valid_file "T";width
2.1926in;height 2.7913in;depth 0pt;original-width 2.1543in;original-height
2.7484in;cropleft "0";croptop "1";cropright "1";cropbottom "0";tempfilename
'TN9QPD40.wmf';tempfile-properties "XPR";}}And we left visible corner cube
arrays on the Moon so we can bounce laser beams off of them and monitor the
Earth-Moon distance.

\FRAME{dhFU}{2.2219in}{2.7638in}{0pt}{\Qcb{{\protect\small Apollo
Retroreflector (Images in the Public Domain courtesy NASA)}}}{}{Figure}{%
\special{language "Scientific Word";type "GRAPHIC";maintain-aspect-ratio
TRUE;display "USEDEF";valid_file "T";width 2.2219in;height 2.7638in;depth
0pt;original-width 2.1819in;original-height 2.7207in;cropleft "0";croptop
"1";cropright "1";cropbottom "0";tempfilename
'TN9QPD41.wmf';tempfile-properties "XPR";}}

\subsection{Reflections, Objects, and seeing}

Armed with the law of reflection, we can start to understand how we see
things. Using the ray concept, we can say that a ray of light must leave the
light source. That ray then reflects from something. Suppose you look at the
person sitting next to you in class. We should wonder, how is it that we can
see them? We can only detect (see) light that gets to our eyes. Let's trace
the light from it's source (the light fixture in our room) to our eyes to
see how we see our neighbors.

Light comes from the ceiling lights and goes in all directions. Some of that
light hits our neighbor. And some of that light that hits our neighbor will
reflect. But is the person a specular or diffuse reflector? Once again, we
can only give an answer relative to the wavelength of light. For visible
light, your neighbors do not look like mirrors. They are diffuse reflectors.
Light bounces off of them in every direction. Some of that light that
reflects from your neighbor reflects in the direction of your eyes. That
light can be detected. Your eye is designed to take this diverging set of
rays and condense it into a picture of the person that your brain can
interpret.\FRAME{dhF}{3.173in}{1.3534in}{0pt}{}{}{Figure}{\special{language
"Scientific Word";type "GRAPHIC";maintain-aspect-ratio TRUE;display
"USEDEF";valid_file "T";width 3.173in;height 1.3534in;depth
0pt;original-width 3.128in;original-height 1.3188in;cropleft "0";croptop
"1";cropright "1";cropbottom "0";tempfilename
'TN9QPD42.wmf';tempfile-properties "XPR";}}We tend to not draw the rays that
bounce off the diffuse reflector but that don't get to our eyes, because we
don't see them. So a ray diagram is usually much simpler.\FRAME{dhF}{2.2926in%
}{1.9951in}{0pt}{}{}{Figure}{\special{language "Scientific Word";type
"GRAPHIC";maintain-aspect-ratio TRUE;display "USEDEF";valid_file "T";width
2.2926in;height 1.9951in;depth 0pt;original-width 2.2528in;original-height
1.9562in;cropleft "0";croptop "1";cropright "1";cropbottom "0";tempfilename
'TN9QPD43.wmf';tempfile-properties "XPR";}}This is easy to understand, but
we must keep in mind the wildly fluctuating superposition of waves that is
masked by our macroscopic view.

We can use the idea of a ray diagram to solve problems. Suppose you hold a
mirror half a meter in front of you and look at your reflection. Where would
the reflection appear to be?\FRAME{dhF}{1.9847in}{1.9527in}{0pt}{}{}{Figure}{%
\special{language "Scientific Word";type "GRAPHIC";maintain-aspect-ratio
TRUE;display "USEDEF";valid_file "T";width 1.9847in;height 1.9527in;depth
0pt;original-width 1.9458in;original-height 1.9147in;cropleft "0";croptop
"1";cropright "1";cropbottom "0";tempfilename
'TN9QPD44.wmf';tempfile-properties "XPR";}}Our mind assumes that the light
only travels in straight lines. This is not always true, as we shall see.
But in the light processing center in our brain the algorithm assumes the
light traveled in a straight line. Then, we can use the idea of rays to see
where the light appears to be from. We follow the light backwards from our
eyes to the mirror, but because our mind processing doesn't know about the
reflection we keep tracing the light ray through the mirror. Of course the
light doesn't come from behind the mirror. But our brain processing thinks
it does. The light originally came from us, so it is in a way shaped like
us. Our brain processing tells us there is another one of us about a half a
meter behind the mirror. Our frontal lobe knows there is not a second one of
us, so we interpret the light signal from our brain processing as an image
of us in the mirror.

The image is half a meter behind the mirror. Now suppose we look at an image
of that image in a mirror behind us. \FRAME{dhF}{2.6273in}{1.004in}{0pt}{}{}{%
Figure}{\special{language "Scientific Word";type
"GRAPHIC";maintain-aspect-ratio TRUE;display "USEDEF";valid_file "T";width
2.6273in;height 1.004in;depth 0pt;original-width 2.5858in;original-height
0.9712in;cropleft "0";croptop "1";cropright "1";cropbottom "0";tempfilename
'TN9QPD45.wmf';tempfile-properties "XPR";}}The ray diagram makes it easy to
see that the image will appear to be $2\unit{m}$ behind the big mirror. You
might not feel that this was so easy, but you might find it is not so bad in
a problem in your near future.

%TCIMACRO{%
%\TeXButton{Basic Equations}{\vspace{0.5in}\hspace{-1.3in}{\LARGE Basic Equations\vspace{0.25in}}}}%
%BeginExpansion
\vspace{0.5in}\hspace{-1.3in}{\LARGE Basic Equations\vspace{0.25in}}%
%EndExpansion
\begin{equation}
E=hf
\end{equation}
\begin{equation}
h=6.63\times 10^{-34}\unit{J}\unit{s}
\end{equation}

\[
c=299792458\frac{\unit{m}}{\unit{s}}
\]%
\[
\theta _{r}=\theta _{i}
\]

\section{Refraction 3.1.3, 3.1.4}

We now know about reflection. But you probably know light can go through
some materials. For many of us, this is an everyday experience as light goes
through our eyeglasses and contacts. But there are also stained glass
windows and windshields and many other examples of light going through
materials. We need to include this effect in our study of light.

\FRAME{dhF}{3.83in}{2.8871in}{0pt}{}{}{Figure}{\special{language "Scientific
Word";type "GRAPHIC";maintain-aspect-ratio TRUE;display "USEDEF";valid_file
"T";width 3.83in;height 2.8871in;depth 0pt;original-width
3.781in;original-height 2.8444in;cropleft "0";croptop "1";cropright
"1";cropbottom "0";tempfilename 'TN9QPD46.wmf';tempfile-properties "XPR";}}

%TCIMACRO{%
%\TeXButton{Fundamental Concepts}{\vspace{0.10in}\hspace{-0.37in}{\Large Fundamental Concepts\vspace{0.2in}}}}%
%BeginExpansion
\vspace{0.10in}\hspace{-0.37in}{\Large Fundamental Concepts\vspace{0.2in}}%
%EndExpansion

\begin{itemize}
\item Refraction

\item Total internal reflection
\end{itemize}

In the last lecture we split up electromagnetic waves into different kinds
of light. We categorized the different kinds of light by noting they have
different wavelength ranges. We found visible light, of course, but also
infrared and ultraviolet and x-rays and radio waves. Not all surfaces
reflect all the light. Some, like the lenses shown below, reflect some light
at visible wavelengths, but are transparent so most of the light travels
through them. \FRAME{dhF}{3.2846in}{1.6042in}{0pt}{}{}{Figure}{\special%
{language "Scientific Word";type "GRAPHIC";maintain-aspect-ratio
TRUE;display "USEDEF";valid_file "T";width 3.2846in;height 1.6042in;depth
0pt;original-width 3.2396in;original-height 1.5679in;cropleft "0";croptop
"1";cropright "1";cropbottom "0";tempfilename
'TN9QPD47.wmf';tempfile-properties "XPR";}}We need a way to deal with
transparent materials. This is tricky, because different wavelengths of
light penetrate different materials in different ways (we will see more on
why later). As an example, this is also a lens\FRAME{dhFU}{1.0732in}{1.6285in%
}{0pt}{\Qcb{{\protect\small Infrared lens, but visible mirror (Image
courtesy US Navy, image in the public domain)}}}{}{Figure}{\special{language
"Scientific Word";type "GRAPHIC";maintain-aspect-ratio TRUE;display
"USEDEF";valid_file "T";width 1.0732in;height 1.6285in;depth
0pt;original-width 5.4587in;original-height 8.3195in;cropleft "0";croptop
"1";cropright "1";cropbottom "0";tempfilename
'TN9QPD48.bmp';tempfile-properties "XPR";}}but it clearly is not transparent
at visible wavelengths. It is transparent in the infrared. So what might be
transparent at one wavelength might not be at another.

When light travels into a material, we say it is transmitted. This is
nothing new. Light is a wave, so the part that goes through a junction is
always called the transmitted part just like when we had waves go from one
rope to another. The situation is shown schematically below.

\FRAME{dtbpF}{2.7207in}{2.2719in}{0in}{}{}{Figure}{\special{language
"Scientific Word";type "GRAPHIC";maintain-aspect-ratio TRUE;display
"USEDEF";valid_file "T";width 2.7207in;height 2.2719in;depth
0in;original-width 2.6783in;original-height 2.2321in;cropleft "0";croptop
"1";cropright "1";cropbottom "0";tempfilename
'TN9QPD49.wmf';tempfile-properties "XPR";}}

In the figure we see a ray incident on an air-glass boundary. Some of the
light is reflected just as we saw in our last lecture. But some passes into
the glass. This isn't surprising, light is a wave and at boundaries we
expect reflected and transmitted waves. Notice that the angle between the
normal and the new transmitted ray is \emph{not} equal to the incident angle
for the incoming ray. We say the ray has been bent or \emph{refracted} by
the change in material. Many experiments were performed to find a
relationship between the incident and the refracted angles. It was found
that 
\begin{equation}
\frac{\sin \left( \theta _{2}\right) }{\sin \left( \theta _{1}\right) }=%
\frac{v_{2}}{v_{1}}=\text{constant}
\end{equation}%
Many optics books write this as 
\begin{equation}
\frac{\sin \left( \theta _{t}\right) }{\sin \left( \theta _{i}\right) }=%
\frac{v_{t}}{v_{i}}=\text{constant}
\end{equation}%
where the subscript $i$ stands for \textquotedblleft
incident\textquotedblright\ and the subscript $t$ stands for
\textquotedblleft transmitted.\textquotedblright\ \FRAME{dtbpF}{3.4229in}{%
2.8548in}{0pt}{}{}{Figure}{\special{language "Scientific Word";type
"GRAPHIC";maintain-aspect-ratio TRUE;display "USEDEF";valid_file "T";width
3.4229in;height 2.8548in;depth 0pt;original-width 3.3771in;original-height
2.8115in;cropleft "0";croptop "1";cropright "1";cropbottom "0";tempfilename
'TN9QPD4A.wmf';tempfile-properties "XPR";}}\ Note that we are using the fact
that the speed of light changes in a material. We should probably recall why
this should occur

\subsection{Speed of light in a material}

In a vacuum, light travels as a disturbance in the electromagnetic field
with nothing to encounter. In a material (like glass) the light waves
continually hit atoms. We have not studied antennas, but I\ think many of
you know that an antenna works because the electrons in the metal act like
driven harmonic oscillators. The incoming radio waves drive the electron
motion. In a piece of glass each atom has electrons, and the atoms act like
little antennas, their electrons moving and absorbing the light. But the
atom cannot keep the extra energy\footnote{%
If you get to take Modern Physics (PH279) and Quantum Physics (PH433) you
will learn how this works.}, so it is readmitted as a new photon of light.
It travels to the next atom and the process repeats. Quantum mechanics tells
us that there is a time delay in the re-emission of the light. This
re-emitted light mixes with what is left of the original wave. And this
causes the propagation of the combined light wave to slow down. Thus the
speed of light is slower in a material.\FRAME{dhF}{4.6267in}{1.8135in}{0pt}{%
}{}{Figure}{\special{language "Scientific Word";type
"GRAPHIC";maintain-aspect-ratio TRUE;display "USEDEF";valid_file "T";width
4.6267in;height 1.8135in;depth 0pt;original-width 4.574in;original-height
1.7763in;cropleft "0";croptop "1";cropright "1";cropbottom "0";tempfilename
'TN9QPD4B.wmf';tempfile-properties "XPR";}}But why does the light turn as it
slows down?

As a mechanical analog, consider a rolling barrel. Suppose you are helping a
neighbor move and they have barrels of wheat. And as often happens in a
neighbor move, one of the barrels get's dropped. It lands on the driveway
and rolls toward the grass.\FRAME{dhF}{2.4604in}{1.8005in}{0pt}{}{}{Figure}{%
\special{language "Scientific Word";type "GRAPHIC";maintain-aspect-ratio
TRUE;display "USEDEF";valid_file "T";width 2.4604in;height 1.8005in;depth
0pt;original-width 2.4189in;original-height 1.7634in;cropleft "0";croptop
"1";cropright "1";cropbottom "0";tempfilename
'TN9QPD4C.wmf';tempfile-properties "XPR";}}As the barrel rolls from a flat
low-friction concrete to a higher-friction grass lawn, the friction slows
the barrel. If the barrel hits the lawn parallel to the boundary (so it's
velocity vector is perpendicular to the boundary), then the barrel continues
in the same direction at the slower speed. But if it hits at an angle, the
leading edge is slowed first. \FRAME{dhF}{2.3575in}{1.7236in}{0pt}{}{}{Figure%
}{\special{language "Scientific Word";type "GRAPHIC";maintain-aspect-ratio
TRUE;display "USEDEF";valid_file "T";width 2.3575in;height 1.7236in;depth
0pt;original-width 2.725in;original-height 1.9847in;cropleft "0";croptop
"1";cropright "1";cropbottom "0";tempfilename
'TN9QPD4D.wmf';tempfile-properties "XPR";}}This makes the trailing edge
travel faster than the leading edge, and the barrel turns slightly.\FRAME{dhF%
}{2.4267in}{1.7711in}{0pt}{}{}{Figure}{\special{language "Scientific
Word";type "GRAPHIC";maintain-aspect-ratio TRUE;display "USEDEF";valid_file
"T";width 2.4267in;height 1.7711in;depth 0pt;original-width
3.7948in;original-height 2.7622in;cropleft "0";croptop "1";cropright
"1";cropbottom "0";tempfilename 'TN9QPD4E.wmf';tempfile-properties "XPR";}}

We expect the same behavior from light.\FRAME{dhF}{2.8236in}{2.0237in}{0pt}{%
}{}{Figure}{\special{language "Scientific Word";type
"GRAPHIC";maintain-aspect-ratio TRUE;display "USEDEF";valid_file "T";width
2.8236in;height 2.0237in;depth 0pt;original-width 2.7804in;original-height
1.9847in;cropleft "0";croptop "1";cropright "1";cropbottom "0";tempfilename
'TN9QPD4F.wmf';tempfile-properties "XPR";}}We can see that the left hand
side of the wave hits the slower (green) material first and slows down. The
rest of the wave front moves quicker. The result is the turning of the wave.%
\footnote{%
This is a bit of a simplification, but it will due for now. If you are lucky
enough to take a junior level optics class, you will revisit this. And we
will look at this again in our next lecture.}

\subsection{Change of wavelength}

We have found that when a wave enters a material, its speed may change. But
we remember from wave theory%
\begin{equation}
v=\lambda f
\end{equation}

And it is time to review: does $\lambda $ change, or does $f$ change? If you
will recall, we found that the change in speed at the boundary changes the
wavelength. Recall that if we go from a fast material to a slow material,
the forward part of the wave slows and the rest of the wave catches up to
it. \FRAME{dhF}{2.2727in}{1.2644in}{0pt}{}{}{Figure}{\special{language
"Scientific Word";type "GRAPHIC";maintain-aspect-ratio TRUE;display
"USEDEF";valid_file "T";width 2.2727in;height 1.2644in;depth
0pt;original-width 3.4618in;original-height 1.9147in;cropleft "0";croptop
"1";cropright "1";cropbottom "0";tempfilename
'TN9QPD4G.wmf';tempfile-properties "XPR";}}This will compress pulses, and
lower the wavelength. Now that we know more about light we can also argue
that $f$ cannot change because 
\[
E=hf
\]%
If $f$ changed, then we would either require an input of energy or we would
store energy at the boundary because%
\[
\Delta f=\frac{\Delta E}{h}
\]%
This can't be true. If the wavelength changes, there is no change in energy.

And since 
\[
v_{i}=\lambda _{i}f
\]%
and 
\[
v_{t}=\lambda _{t}f
\]%
then the ratio%
\[
\frac{v_{i}}{v_{t}}=\frac{\lambda _{i}}{\lambda _{t}}
\]%
and we have a solution for the wavelength in the material%
\[
\lambda _{t}=\lambda _{i}\frac{v_{t}}{v_{i}}
\]%
which agrees with our previous analysis.

\subsection{Index of refraction and Snell's Law}

Because the equation%
\begin{equation}
\frac{\sin \left( \theta _{i}\right) }{\sin \left( \theta _{t}\right) }=%
\frac{v_{t}}{v_{i}}=\text{constant}
\end{equation}%
has a constant ratio of velocities, it is convenient to define a term that
represents that ratio. Think, it would be kind of nice to say how much
slower the light is going once it enters the slower glass material. We could
take the current velocity and divide by the speed of light in vacuum to get
the percentage of the speed of light we are going%
\[
\frac{v}{c}=\text{percentage of the speed of light}
\]%
and that makes a lot of sense. But that is not what the optics community
did. They defined a term. 
\begin{equation}
n\equiv \frac{c}{v}
\end{equation}%
which is $1/$percentage of the speed of light. Which is not as intuitive as
the percentage, itself, was. But it caries the same information 
\[
n=\frac{1}{\text{percentage of the speed of light}}
\]%
So we can use this to get an idea of how much slower our transmitted light
is going. This odd ratio of the speed of light to the average speed in the
material has a name. It is called the \emph{index of refraction.} We can
write Snell's law in terms of the index of refraction of air and glass (or
any other two materials). Start with%
\begin{equation}
\frac{\sin \left( \theta _{i}\right) }{\sin \left( \theta _{t}\right) }=%
\frac{v_{i}}{v_{t}}
\end{equation}%
but let's multiply the right hand side by $1,$ only let's write $1$ as 
\[
1=\frac{c}{c}
\]%
which is a very odd way to write $1.$ But then%
\begin{equation}
\frac{\sin \left( \theta _{i}\right) }{\sin \left( \theta _{t}\right) }=%
\frac{v_{i}c}{v_{t}c}
\end{equation}%
which we could write as%
\begin{equation}
\frac{\sin \left( \theta _{i}\right) }{\sin \left( \theta _{t}\right) }%
=\left( \frac{c}{v_{t}}\right) \left( \frac{v_{i}}{c}\right)
\end{equation}%
And with our index of refraction definition, this is just%
\begin{equation}
\frac{\sin \left( \theta _{i}\right) }{\sin \left( \theta _{t}\right) }%
=\left( n_{t}\right) \left( \frac{1}{n_{i}}\right)
\end{equation}%
and for our case $n_{i}=n_{air}\approx 1$ so 
\begin{equation}
\frac{\sin \left( \theta _{i}\right) }{\sin \left( \theta _{t}\right) }%
=n_{glass}
\end{equation}

Suppose we don't have a vacuum (or air that is close to a vacuum). We can
write our formula as 
\begin{equation}
n_{i}\sin \left( \theta _{i}\right) =n_{t}\sin \left( \theta _{t}\right)
\end{equation}%
where we have 
\begin{equation}
n_{i}\equiv \frac{c}{v_{i}}
\end{equation}%
and 
\begin{equation}
n_{t}\equiv \frac{c}{v_{t}}
\end{equation}

This is called \emph{Snell's law of refraction}.

Using the index of refraction we can write our equation relating the ratio
of velocities and wavelengths as%
\begin{equation}
\frac{v_{i}}{v_{t}}=\frac{\lambda _{i}}{\lambda _{t}}=\frac{\frac{c}{n_{i}}}{%
\frac{c}{n_{t}}}=\frac{n_{t}}{n_{i}}
\end{equation}%
which gives 
\begin{equation}
\lambda _{i}n_{i}=\lambda _{t}n_{t}
\end{equation}%
and if we have vacuum and a single material we can find the index of
refraction from 
\begin{equation}
n=\frac{\lambda }{\lambda _{in}}
\end{equation}%
where $\lambda _{in}=\lambda _{t}$ is the wavelength in the material and $%
\lambda =\lambda _{i}$ is the wavelength outside the material.

\subsection{Total Internal Reflection}

Up to now we have assumed that light was coming from a region of low index
of refraction (air) into a region of high index of refraction (water or
glass). We should pause to look at what can happen if we go the other way.%
\FRAME{dtbpF}{3in}{2.5339in}{0pt}{}{}{Figure}{\special{language "Scientific
Word";type "GRAPHIC";maintain-aspect-ratio TRUE;display "USEDEF";valid_file
"T";width 3in;height 2.5339in;depth 0pt;original-width
3.4601in;original-height 2.9187in;cropleft "0";croptop "1";cropright
"1";cropbottom "0";tempfilename 'TN9QPD4H.wmf';tempfile-properties "XPR";}}

We start with Snell's law%
\begin{equation}
n_{i}\sin \theta _{i}=n_{t}\sin \theta _{t}
\end{equation}%
but this time $n_{glass}=n_{i}$ and $n_{t}\approx 1$ so 
\begin{equation}
n_{glass}\sin \theta _{i}=\sin \theta _{t}
\end{equation}%
which gives 
\begin{equation}
\theta _{t}=\sin ^{-1}\left( n_{glass}\sin \theta _{i}\right)
\end{equation}%
If we take $n_{glass}=1.33$ (which is really $n$ for water) we can plot this
expression as a function of $\theta _{i}$

\FRAME{dtbpF}{2.8643in}{1.9804in}{0pt}{}{}{Figure}{\special{language
"Scientific Word";type "GRAPHIC";maintain-aspect-ratio TRUE;display
"USEDEF";valid_file "T";width 2.8643in;height 1.9804in;depth
0pt;original-width 3.6642in;original-height 2.5235in;cropleft "0";croptop
"1";cropright "1";cropbottom "0";tempfilename
'TN9QPD4I.wmf';tempfile-properties "XPR";}}we see that at $\theta
_{i}=\allowbreak 0.850\,91\unit{rad}$ $(48.\,\allowbreak 754\unit{%
%TCIMACRO{\U{b0}}%
%BeginExpansion
{{}^\circ}%
%EndExpansion
})$ the curve becomes infinitely steep. If we use this value in our equation
this gives%
\begin{eqnarray}
\theta _{t} &=&\sin ^{-1}\left( n\sin \left( \allowbreak 0.850\,91\right)
\right) \\
&=&1.\,\allowbreak 570\,8\unit{rad} \\
&=&90\unit{%
%TCIMACRO{\U{b0}}%
%BeginExpansion
{{}^\circ}%
%EndExpansion
}
\end{eqnarray}%
The light skims along the edge of the water! \FRAME{dhF}{2.3756in}{1.9251in}{%
0pt}{}{}{Figure}{\special{language "Scientific Word";type
"GRAPHIC";maintain-aspect-ratio TRUE;display "USEDEF";valid_file "T";width
2.3756in;height 1.9251in;depth 0pt;original-width 2.3359in;original-height
1.887in;cropleft "0";croptop "1";cropright "1";cropbottom "0";tempfilename
'TN9QPD4J.wmf';tempfile-properties "XPR";}}We can find the value of $\theta
_{i}$ that makes this happen without graphing. Set $\theta _{t}=90\unit{%
%TCIMACRO{\U{b0}}%
%BeginExpansion
{{}^\circ}%
%EndExpansion
}$ then%
\begin{equation}
\theta _{i}=\theta _{c}\equiv \sin ^{-1}\left( \frac{1}{n}\right)
\end{equation}%
We give this value of $\theta _{i}$ a special name. It is the \emph{critical
angle} for internal reflection. But what happens if we go farther than this (%
$\theta _{i}>\theta _{c}$). We will no longer have a transmitted ray. But
the wave energy has to go somewhere! The ray will be reflected. This is why
when you dive into a pool and look up, you see a region of the roof of the
pool area (or sky) but off to the side of the pool the surface looks
mirrored. It is also why you sometimes see the sides of a fish tank appear
to be mirrored when you look through the front. \FRAME{dhF}{1.7616in}{%
1.3258in}{0pt}{}{}{Figure}{\special{language "Scientific Word";type
"GRAPHIC";maintain-aspect-ratio TRUE;display "USEDEF";valid_file "T";width
1.7616in;height 1.3258in;depth 0pt;original-width 1.7236in;original-height
1.2912in;cropleft "0";croptop "1";cropright "1";cropbottom "0";tempfilename
'TN9QPD4K.wmf';tempfile-properties "XPR";}}It is also why cut gems (like
diamonds) sparkle. They capture the light with facets that are cut at angles
that create total internal reflection. The light that enters the gem comes
back out the front.

\subsection{Fiber Optics}

Beyond pretty pebbles, this effect is very useful! It is the heart and soul
of fiber optics. \FRAME{dtbpF}{3.1993in}{1.7882in}{0in}{}{}{Figure}{\special%
{language "Scientific Word";type "GRAPHIC";maintain-aspect-ratio
TRUE;display "USEDEF";valid_file "T";width 3.1993in;height 1.7882in;depth
0in;original-width 3.233in;original-height 1.7935in;cropleft "0";croptop
"1";cropright "1";cropbottom "0";tempfilename
'TN9QPD4L.wmf';tempfile-properties "XPR";}}

An interior material with a lower index of refraction is inclosed in a
cladding with a higher index. This creates a light pipe that traps the light
in the fiber. \FRAME{dhF}{3.4662in}{1.535in}{0pt}{}{}{Figure}{\special%
{language "Scientific Word";type "GRAPHIC";maintain-aspect-ratio
TRUE;display "USEDEF";valid_file "T";width 3.4662in;height 1.535in;depth
0pt;original-width 3.4203in;original-height 1.4987in;cropleft "0";croptop
"1";cropright "1";cropbottom "0";tempfilename
'TN9QPD4M.wmf';tempfile-properties "XPR";}}

Modern fibers don't always have a hard boundary. The fibers have a gradual
change in index of refraction that changes the direction of the light
gradually. This keeps the light in the fiber but tends to direct along the
fiber so the beam is not crisscrossing as it goes.

The cutting edge of fiber design today uses hollow fibers or fibers filled
with different index material.

\FRAME{dtbpFU}{1.6319in}{1.5973in}{0pt}{\Qcb{{\protect\small Hollow-Core
Fiber (Courtesy Defense Advanced Research Projects Agency (DARPA), image in
the public domain)}}}{}{Figure}{\special{language "Scientific Word";type
"GRAPHIC";maintain-aspect-ratio TRUE;display "USEDEF";valid_file "T";width
1.6319in;height 1.5973in;depth 0pt;original-width 1.6338in;original-height
1.5992in;cropleft "0";croptop "1";cropright "1";cropbottom "0";tempfilename
'TN9QPD4N.wmf';tempfile-properties "XPR";}}%
%TCIMACRO{%
%\TeXButton{Basic Equations}{\vspace{0.5in}\hspace{-1.3in}{\LARGE Basic Equations\vspace{0.25in}}}}%
%BeginExpansion
\vspace{0.5in}\hspace{-1.3in}{\LARGE Basic Equations\vspace{0.25in}}%
%EndExpansion

\[
n_{i}\sin \theta _{i}=n_{t}\sin \theta _{t}
\]%
\[
\lambda _{i}n_{i}=\lambda _{t}n_{t}
\]%
\[
\theta _{c}\equiv \sin ^{-1}\left( \frac{1}{n}\right)
\]

\section{Dispersion and Huygen's Principle 3.1.5 3.1.6}

We have discovered some important things about how light travels through
materials. But there is something we missed. It turns out the index of
refraction depends on wavelength! This probably shouldn't surprise us. Since
the average speed of our wave in a material depends on the spacing of the
atoms, it seems reasonable that long wavelengths and short wavelengths might
see those atoms differently. This has some profound effects on what light
will do in a material. We will look at this wavelength dependence on
refraction in this lecture.

%TCIMACRO{%
%\TeXButton{Fundamental Concepts}{\vspace{0.10in}\hspace{-0.37in}{\Large Fundamental Concepts\vspace{0.2in}}}}%
%BeginExpansion
\vspace{0.10in}\hspace{-0.37in}{\Large Fundamental Concepts\vspace{0.2in}}%
%EndExpansion

\begin{itemize}
\item The index of refraction is wavelength dependent

\item That different wavelengths bend different amounts when refracted is
called \emph{dispersion}

\item White light is a superposition of many other frequencies
\end{itemize}

\subsection{Dispersion}

\FRAME{dhF}{4.3613in}{2.2745in}{0pt}{}{}{Figure}{\special{language
"Scientific Word";type "GRAPHIC";maintain-aspect-ratio TRUE;display
"USEDEF";valid_file "T";width 4.3613in;height 2.2745in;depth
0pt;original-width 4.3102in;original-height 2.2347in;cropleft "0";croptop
"1";cropright "1";cropbottom "0";tempfilename
'TN9QPD4O.wmf';tempfile-properties "XPR";}}

Who hasn't played with a prism? We immediately recognize a rainbow. But why
does the prism make a rainbow? The secret lies in the nature of the
refractive index. We know that some wavelengths will pass through a
material, say, a piece of glass. But some (like infrared wavelengths) won't.
But there is more to this. Not all the wavelengths that do pass through the
glass will have the same transmitted speed. The speed in the material
depends on the wavelength! Think, blue wavelengths are about the size of
atoms. This means blue photons will hit atoms more frequently. And that will
slow them down more. Red wavelengths are bigger than the atoms in a material
and because of this red photons are not as likely to be absorbed by the
atoms. So red photos will go faster in the material. And this means the
index of refraction depends on the speed.\FRAME{fhFU}{4.3535in}{2.6195in}{0pt%
}{\Qcb{{\protect\small Index of refraction as a function of wavelength (
Ohara optical glass http://www.oharacorp.com/fused-silica-quartz.html data
and Schott optical glass data
http://www.uqgoptics.com/materials\_glasses\_schott.aspx)}}}{\Qlb{Dispersion
of Glass}}{Figure}{\special{language "Scientific Word";type
"GRAPHIC";maintain-aspect-ratio TRUE;display "USEDEF";valid_file "T";width
4.3535in;height 2.6195in;depth 0pt;original-width 5.1436in;original-height
3.084in;cropleft "0";croptop "1";cropright "1";cropbottom "0";tempfilename
'dispersion0.wmf';tempfile-properties "XNPR";}}

We can see light from about $400\unit{nm}$ to about $700\unit{nm}.$ And
notice for the three kinds of glass in the figure, the index of refraction
changes over this region of the electromagnetic spectrum. Think of what this
means for Snell's law. We learned last lecture that 
\[
n_{i}\sin \theta _{i}=n_{t}\sin \theta _{t}
\]%
and if the light goes from air $\left( n_{i}\approx 1\right) $ into some
kind of glass, then 
\[
\sin \theta _{i}=n_{t}\sin \theta _{t}
\]%
and the transmitted angle just depends on $n_{t}.$ This means that as light
enters a material, different wavelengths will be refracted at different
angles. White light is a superposition of many light waves with many
different wavelengths. Thus white light is pulled apart by refraction into a
rainbow. This process is called \emph{dispersion}.

The graph tells us that blue light bends more than red light. \FRAME{dhF}{%
2.4379in}{1.2816in}{0pt}{}{}{Figure}{\special{language "Scientific
Word";type "GRAPHIC";maintain-aspect-ratio TRUE;display "USEDEF";valid_file
"T";width 2.4379in;height 1.2816in;depth 0pt;original-width
6.5207in;original-height 3.4143in;cropleft "0";croptop "1";cropright
"1";cropbottom "0";tempfilename 'TN9QPD4P.wmf';tempfile-properties "XPR";}}

We call the change in direction measured from the original direction of
travel the \emph{angle of deviation} and give it the symbol $\delta $. The
colors we can see are called the visible spectrum. Although this dispersive
element (prism) method of making a spectrum uses the ides of geometric
optics, it is still the wave nature of light that makes the waves of
different wavelengths refract differently.

Let's look at a natural rainbow. The dispersion is caused buy small droplets
of water. The white sunlight enters the drop and is dispersed. Most of the
light refracts back out of the back of the droplet. But a small amount
reflects off the back of the drop and then refracts out of the drop towards
the front, again being dispersed. Red light leaves the drop at about $42%
\unit{%
%TCIMACRO{\U{b0}}%
%BeginExpansion
{{}^\circ}%
%EndExpansion
}$ from its input direction, and violet light leaves at about $40\unit{%
%TCIMACRO{\U{b0}}%
%BeginExpansion
{{}^\circ}%
%EndExpansion
}.$ The other colors leave at angles in between these two angles.\FRAME{dtbpF%
}{3.2681in}{2.3921in}{0pt}{}{}{Figure}{\special{language "Scientific
Word";type "GRAPHIC";maintain-aspect-ratio TRUE;display "USEDEF";valid_file
"T";width 3.2681in;height 2.3921in;depth 0pt;original-width
3.2232in;original-height 2.3514in;cropleft "0";croptop "1";cropright
"1";cropbottom "0";tempfilename 'TN9QPD4Q.wmf';tempfile-properties "XPR";}}%
This effectively separates all the different wavelengths and we see a
rainbow at angles between $40\unit{%
%TCIMACRO{\U{b0}}%
%BeginExpansion
{{}^\circ}%
%EndExpansion
}$ and $42\unit{%
%TCIMACRO{\U{b0}}%
%BeginExpansion
{{}^\circ}%
%EndExpansion
}$ from the incoming light direction.

\subsubsection{Calculation of $n$ using a prism}

Let's do an engineering problem using the idea of refraction in a prism.
Let's find the index of refraction of a the prism material. Suppose we make
a prism as shown. We know the angle $\Phi $ and can measure the exit angle $%
\delta .$ In terms of these two variables, what is $n?$

\FRAME{dtbpF}{2.4595in}{1.9259in}{0in}{}{}{Figure}{\special{language
"Scientific Word";type "GRAPHIC";maintain-aspect-ratio TRUE;display
"USEDEF";valid_file "T";width 2.4595in;height 1.9259in;depth
0in;original-width 2.418in;original-height 1.8887in;cropleft "0";croptop
"1";cropright "1";cropbottom "0";tempfilename
'TN9QPD4R.wmf';tempfile-properties "XPR";}}

Using the notation indicated in the figure, we choose $\theta _{i1}$ such
that the interior ray is horizontal.\footnote{%
WARNING! in an upcoming homework problem you may not be able make the same
assumptions!} This is a refraction problem, so we will want to use Snell's
law. 
\[
n_{i1}\sin \theta _{i1}=n_{t1}\sin \theta _{t1}
\]%
And this has $n_{2}$ in it. That is what we want to find. We know $n_{1}$
because the prism is in air so $n_{1}\approx 1.$ And suppose we know $\Phi $
and $\delta $ from measurement. Our process will be to find $\theta _{1}$
and $\theta _{2}.$ And $n_{t1}$ is $n$ for the material. The we can use 
\[
n_{t1}=\frac{\sin \theta _{i1}}{\sin \theta _{t1}}
\]%
to find $n.$ But we have to find the angles. Let's see one way to do this.
In geometry it is fair game to extend lines and even make some new lines of
our own. I have done this in the figure. You likely will have to do your own
line extensions and additions in the problems you do. But by using the added
lines we can realize that

\begin{equation}
\theta _{i1}=\theta _{t1}+\alpha
\end{equation}%
and that 
\begin{equation}
\delta =180-\beta
\end{equation}%
and it is also true that

\begin{equation}
180=\beta +2\alpha
\end{equation}%
or 
\begin{equation}
180-\beta =2\alpha
\end{equation}

Then

\begin{equation}
\delta =2\alpha
\end{equation}%
and%
\begin{equation}
\alpha =\frac{\delta }{2}
\end{equation}

Now also realize that

\begin{equation}
90=\alpha +\theta _{t1}+\phi
\end{equation}

and

\begin{equation}
180=\Phi +2\alpha +2\phi
\end{equation}%
or%
\begin{equation}
90=\frac{\Phi }{2}+\alpha +\phi
\end{equation}%
Then%
\begin{eqnarray}
\alpha +\theta _{t1}+\phi &=&\frac{\Phi }{2}+\alpha +\phi \\
\theta _{t1} &=&\frac{\Phi }{2}
\end{eqnarray}%
We can put these in our equation for $\theta _{i1}$%
\begin{eqnarray}
\theta _{i1} &=&\theta _{t1}+\alpha \\
&=&\frac{\Phi }{2}+\frac{\delta }{2} \\
&=&\frac{\Phi +\delta }{2}
\end{eqnarray}%
Now we can use Snell's Law%
\begin{eqnarray}
n_{t} &=&\frac{\sin \left( \theta _{i1}\right) }{\sin \left( \theta
_{t1}\right) } \\
&=&\frac{\sin \left( \frac{\Phi +\delta }{2}\right) }{\sin \left( \frac{\Phi 
}{2}\right) }
\end{eqnarray}%
and since we know $\Phi $ and $\delta $, we can find $n_{t}=n$ of the
material.

\subsection{Huygen's Principle}

We have learned the laws of reflection and refraction. And we know the
scattering of the light by atoms in a material cause the bending of light
that we call refraction. But let's look at this again in more detail.

To understand what is going on at a fundamental level we need to remember
that light is a wave. And waves can hit and splash off of objects. Sound
waves could bounce off of things. So can light waves. Only, the light waves
that we can see are very small. So they bounce off of very small things like
atoms. Suppose we send our light into a material. It will hit the atoms and
be absorbed and bounce off. We know this, but there is more. The light hits
the atoms and is absorbed by exciting electrons. When the electrons return
to their original state the light is re-emitted. But this makes a new light
wave coming from the atoms. And that new wave is roughly spherical (think,
most waves start spherical). If we have many atoms, the light will hit many
atoms and make many little spherical \emph{scattered} waves. An early
researcher had an idea of how to find the direction of the light using this
underlying wave nature of light. We could follow the light path by finding
the little spherical scattered waves and look for where their crests line
up. This would be constructive interference, and this is where the ray of
light would end up! \FRAME{dtbpF}{0.5047in}{1.2418in}{0pt}{}{}{Figure}{%
\special{language "Scientific Word";type "GRAPHIC";maintain-aspect-ratio
TRUE;display "USEDEF";valid_file "T";width 0.5047in;height 1.2418in;depth
0pt;original-width 0.33in;original-height 0.8551in;cropleft "0";croptop
"1";cropright "1";cropbottom "0";tempfilename
'TN9QPD4S.wmf';tempfile-properties "XPR";}}This early researcher was named
Huygens, so this method of finding the new wave is called \emph{Huygen's
principle.} He simplified the process by just drawing a line that connected
the crests of each little scattered wave and said that is where the new
combined wave front would be.

We know that light should reflect off a surface. 
\[
\theta _{i}=\theta _{r}
\]

\FRAME{dtbpF}{2.841in}{1.8086in}{0pt}{}{}{Figure}{\special{language
"Scientific Word";type "GRAPHIC";maintain-aspect-ratio TRUE;display
"USEDEF";valid_file "T";width 2.841in;height 1.8086in;depth
0pt;original-width 3.6748in;original-height 2.3283in;cropleft "0";croptop
"1";cropright "1";cropbottom "0";tempfilename
'TN9QPD4T.wmf';tempfile-properties "XPR";}}Let's see how Huygen's principle
gives this law of reflection. Think about the top row (or about $\lambda /2$
thickness) of scatterers. In the next figure we have drawn such a set of
scatters with an incoming plane wave.\FRAME{dtbpF}{3.3987in}{1.9104in}{0pt}{%
}{}{Figure}{\special{language "Scientific Word";type
"GRAPHIC";maintain-aspect-ratio TRUE;display "USEDEF";valid_file "T";width
3.3987in;height 1.9104in;depth 0pt;original-width 4.6778in;original-height
2.6187in;cropleft "0";croptop "1";cropright "1";cropbottom "0";tempfilename
'TN9QPD4U.wmf';tempfile-properties "XPR";}}On the left, the plane wave crest
is striking a scatter and has started a spherical secondary wave. One period
later the situation might look like this \FRAME{dtbpF}{3.5898in}{2.0176in}{%
0pt}{}{}{Figure}{\special{language "Scientific Word";type
"GRAPHIC";maintain-aspect-ratio TRUE;display "USEDEF";valid_file "T";width
3.5898in;height 2.0176in;depth 0pt;original-width 5.1992in;original-height
2.911in;cropleft "0";croptop "1";cropright "1";cropbottom "0";tempfilename
'TN9QPD4V.wmf';tempfile-properties "XPR";}}where now the first crest has
moved on to the second scatterer and the second crest is at the first
scatter. Notice that the first scatter's spherical secondary wave has grown.
This process continues as the wave moves onward. In another period, we would
have three scattered waves \FRAME{dtbpF}{3.6547in}{2.0574in}{0pt}{}{}{Figure%
}{\special{language "Scientific Word";type "GRAPHIC";maintain-aspect-ratio
TRUE;display "USEDEF";valid_file "T";width 3.6547in;height 2.0574in;depth
0pt;original-width 3.6071in;original-height 2.0185in;cropleft "0";croptop
"1";cropright "1";cropbottom "0";tempfilename
'TN9QPD4W.wmf';tempfile-properties "XPR";}}and in another period, four.%
\FRAME{dtbpF}{3.9409in}{2.2182in}{0pt}{}{}{Figure}{\special{language
"Scientific Word";type "GRAPHIC";maintain-aspect-ratio TRUE;display
"USEDEF";valid_file "T";width 3.9409in;height 2.2182in;depth
0pt;original-width 3.8917in;original-height 2.1776in;cropleft "0";croptop
"1";cropright "1";cropbottom "0";tempfilename
'TN9QPD4X.wmf';tempfile-properties "XPR";}}and then five.\FRAME{dtbpF}{%
3.7645in}{2.1292in}{0pt}{}{}{Figure}{\special{language "Scientific
Word";type "GRAPHIC";maintain-aspect-ratio TRUE;display "USEDEF";valid_file
"T";width 3.7645in;height 2.1292in;depth 0pt;original-width
3.717in;original-height 2.0903in;cropleft "0";croptop "1";cropright
"1";cropbottom "0";tempfilename 'TN9QPD4Y.wmf';tempfile-properties "XPR";}}%
By looking for where the wave crests line up and drawing Huygen's line, we
can find the new crest of the reflected wave. By superimposing these
scattered waves, the reflected wave is formed.

We also know \emph{Snell's law of refraction}. But let's see how Huygen's
principle can be used to show this is true. In our discussion on reflection
we just used the top layer of scatters, but we know that there are more
scatterers in a solid. \FRAME{dtbpF}{3.5509in}{2.0064in}{0pt}{}{}{Figure}{%
\special{language "Scientific Word";type "GRAPHIC";maintain-aspect-ratio
TRUE;display "USEDEF";valid_file "T";width 3.5509in;height 2.0064in;depth
0pt;original-width 4.7063in;original-height 2.6463in;cropleft "0";croptop
"1";cropright "1";cropbottom "0";tempfilename
'TN9QPD4Z.wmf';tempfile-properties "XPR";}}So, what happens inside the
material with all the other layers of scatterers? Of course our secondary
waves are kind of spherical shaped, so the situation is something like this%
\FRAME{dtbpF}{4.1286in}{2.335in}{0pt}{}{}{Figure}{\special{language
"Scientific Word";type "GRAPHIC";maintain-aspect-ratio TRUE;display
"USEDEF";valid_file "T";width 4.1286in;height 2.335in;depth
0pt;original-width 4.0785in;original-height 2.2952in;cropleft "0";croptop
"1";cropright "1";cropbottom "0";tempfilename
'TN9QPD50.wmf';tempfile-properties "XPR";}}and we see by drawing another
Huygen's line that we are forming another wave front \emph{inside the
material}. Our scattering model predicts the transmitted wave as well as the
reflected wave. We can see that Snell's law follows from basic geometry just
as the law of reflection did. It depends on the spacing of the atoms and the
size of the wavelengths. This is a more profound view of the laws of
reflection and refraction than our ray model approach.

Huygen's principle is used many times in the study of optics to help us
understand the geometry of an optics problem and from that form an equation
that makes calculations possible!

Snell and others used measurements to find Snell's law. But now that we
understand the geometry we can use it to find Snell's equation. In the next
figure you can see an incoming wave. We have marked a wavefront $\overline{AB%
}$ that is just starting to hit the interface between the two materials. We
can think of the \emph{wavefront} as a particular crest of the wave. And we
have marked another wavefront, $\overline{CD}$ that is the transmitted wave
front that has just become entirely inside the new material.

\FRAME{dtbpF}{5.2183in}{3.2404in}{0pt}{}{}{Figure}{\special{language
"Scientific Word";type "GRAPHIC";maintain-aspect-ratio TRUE;display
"USEDEF";valid_file "T";width 5.2183in;height 3.2404in;depth
0pt;original-width 5.1612in;original-height 3.1955in;cropleft "0";croptop
"1";cropright "1";cropbottom "0";tempfilename
'TN9QPD51.wmf';tempfile-properties "XPR";}}The time it takes the wave front
to go from $B$ to $D$ would be 
\[
\Delta t_{BD}=\frac{\overline{BD}}{v_{i}}
\]%
but because the new transmitted wave is slower, the transmitted wavefront
has only gone the distance $\overline{AC}$ in the same time%
\[
\Delta t_{AC}=\Delta t_{BD}=\Delta t=\frac{\overline{AC}}{v_{t}}
\]%
We have two triangles. And they share a side $\overline{AD}.$ And they are
right triangles. So the two triangles must be similar triangles. Now let's
look at our incident angle $\theta _{i}$ that of course we measure from a
normal line. We are going to want the $\sin \theta _{i}$ for Snell's law and
we can see that 
\[
\sin \theta _{i}=\frac{\overline{BD}}{\overline{AD}}
\]%
and also 
\[
\sin \theta _{t}=\frac{\overline{AC}}{\overline{AD}}
\]%
and we can solve both for $1/\overline{AD}$ 
\begin{eqnarray*}
\frac{\sin \theta _{i}}{\overline{BD}} &=&\frac{1}{\overline{AD}} \\
\frac{\sin \theta _{t}}{\overline{AC}} &=&\frac{1}{\overline{AD}}
\end{eqnarray*}%
and set them equal so that 
\[
\frac{\sin \theta _{i}}{\overline{BD}}=\frac{\sin \theta _{t}}{\overline{AC}}
\]%
but 
\begin{eqnarray*}
\overline{BD} &=&v_{i}\Delta t \\
\overline{AC} &=&v_{t}\Delta t
\end{eqnarray*}%
so 
\[
\frac{\sin \theta _{i}}{v_{i}\Delta t}=\frac{\sin \theta _{t}}{v_{t}\Delta t}
\]%
And the $\Delta t$ cancels from both sides%
\[
\frac{\sin \theta _{i}}{v_{i}}=\frac{\sin \theta _{t}}{v_{t}}
\]%
We can multiply both sides by $c,$ the speed of light 
\[
c\frac{\sin \theta _{i}}{v_{i}}=c\frac{\sin \theta _{t}}{v_{t}}
\]%
to achieve%
\[
n_{i}\sin \theta _{i}=n_{t}\sin \theta _{t}
\]%
which is Snell's law, this time from basic geometry.

\subsection{Filters and other color phenomena}

Of course, we have assumed without statement, that white light is made up of
all the colors of the rainbow. We should ask why a red shirt is red, or why
passing light through a green film makes the light look green as it leaves.

Both of these phenomena are examples of removing wavelengths from white
light.

In the case of the red shirt, the red dye in the cloth absorbs all of the
visible colors except red. The red is reflected, so the shirt looks red. The
filter is much the same. The green pigment in the film causes nearly all
visible colors to be absorbed except green. So only green light is
transmitted. This is why leaves are green. Chlorophyll absorbs red and blue
wavelengths, so the green is reflected or transmitted.

\FRAME{dhFU}{2.4738in}{2.19in}{0pt}{\Qcb{{\protect\small Chlorophyll
Spectrum (Public Domain image courtesy Kurzon)}}}{}{Figure}{\special%
{language "Scientific Word";type "GRAPHIC";maintain-aspect-ratio
TRUE;display "USEDEF";valid_file "T";width 2.4738in;height 2.19in;depth
0pt;original-width 2.4336in;original-height 2.1508in;cropleft "0";croptop
"1";cropright "1";cropbottom "0";tempfilename
'TN9QPD52.wmf';tempfile-properties "XPR";}}

In our next lecture, we will continue to learn about light and we will think
about how light forms images.

\section{Polarization and Image formation 3.1.7, 3.2.1}

Light is a transverse wave. And it turns out this matters! The direction the
electric field is oscillating leads to what we call \emph{polarization} and
this is used in making things like polarized sun glasses. Let's talk about
this in this lecture. Then, let's talk about what we mean in optics when we
use the word \textquotedblleft image.\textquotedblright

%TCIMACRO{%
%\TeXButton{Fundamental Concepts}{\vspace{0.10in}\hspace{-0.37in}{\Large Fundamental Concepts\vspace{0.2in}}}}%
%BeginExpansion
\vspace{0.10in}\hspace{-0.37in}{\Large Fundamental Concepts\vspace{0.2in}}%
%EndExpansion

\begin{itemize}
\item The direction of the electric field in a plane wave is called the
polarization direction.

\item Natural light is usually a superposition of many waves with random
polarization directions. This light is called unpolarized light.

\item Some materials allow light with one polarization to pass through,
while stopping other polarizations. The polaroid is one such material. Two
crossed polaroids will have a final intensity that follows the relationship $%
I=I_{\max }\cos ^{2}\left( \theta \right) $

\item Light reflecting off a surface may be polarized because of the
absorption and re-emission pattern of light interacting with the material
atoms.

\item Scattered light may be polarized because of anisotropies in the
scatterers.

\item Birefringent materials have different wave speeds in different
directions. This affects the polarization of light entering these materials.

\item An image is a diverging set of rays in a recognizable pattern

\item Images can be formed by reflection and refraction
\end{itemize}

\subsection{Polarization of Light Waves}

We said much earlier in our study of light that it was a transverse wave and
that the medium for this wave is the electromagnetic field. Understanding
electric and magnetic fields is a topic for Principles of Physics III
(PH220) so we can't go to deep into electric and magnetic fields here. Let's
assume that all of space (the whole universe) is filled with this stuff we
call electric and magnetic fields. And let's suppose that these fields can
be a medium for waves. And let's identify these waves as light.

The light waves are transverse waves. So when we look at a light wave in
detail we will see the pieces of the electric and magnetic fields
experiencing simple harmonic motion much like the parts of our spring for
waves on springs. It turns out that the electric field medium parts move
perpendicular to the direction the wave goes, and also perpendicular to the
direction the magnetic field medium parts go. This makes them transverse
waves. We will now show some implications of this fact. In a course in
electromagnetic theory, we often draw light as in the figure below.\FRAME{dhF%
}{3.681in}{3.673in}{0pt}{}{}{Figure}{\special{language "Scientific
Word";type "GRAPHIC";maintain-aspect-ratio TRUE;display "USEDEF";valid_file
"T";width 3.681in;height 3.673in;depth 0pt;original-width
3.7218in;original-height 3.7156in;cropleft "0";croptop "1";cropright
"1";cropbottom "0";tempfilename 'TN9QPD53.wmf';tempfile-properties "XPR";}}%
and that looks a bit like our waves on strings or springs only with the
electric field part going up and down and the magnetic field part going back
and forth.

In what follows we will ignore the magnetic field (marked in the figure as $%
B $) because it's amplitude is very much smaller than the electric field
amplitude. So, we will look at the $E$ field an notice that it goes up and
down in the figure. But we could have light in any orientation. If we look
directly at an approaching beam of light we would \textquotedblleft
see\textquotedblright\ many different orientations as shown in the next
figure.\FRAME{dtbpF}{1.5307in}{1.5895in}{0pt}{}{}{Figure}{\special{language
"Scientific Word";type "GRAPHIC";maintain-aspect-ratio TRUE;display
"USEDEF";valid_file "T";width 1.5307in;height 1.5895in;depth
0pt;original-width 1.4944in;original-height 1.5541in;cropleft "0";croptop
"1";cropright "1";cropbottom "0";tempfilename
'TN9QPD54.wmf';tempfile-properties "XPR";}}

When light beams have waves with many orientations, we say they are \emph{%
unpolarized}. But suppose we were able to align all the light so that all
the waves in the beam were transverse waves in the same orientation. Say,
the one in the next figure.\FRAME{dhF}{0.3226in}{1.3508in}{0pt}{}{}{Figure}{%
\special{language "Scientific Word";type "GRAPHIC";maintain-aspect-ratio
TRUE;display "USEDEF";valid_file "T";width 0.3226in;height 1.3508in;depth
0pt;original-width 1.4688in;original-height 6.2923in;cropleft "0";croptop
"1";cropright "1";cropbottom "0";tempfilename
'TN9QPD55.wmf';tempfile-properties "XPR";}}

Then we would describe the light as \emph{linearly polarized}. The plane
that contains the $E$-field is known as the \emph{polarization plane}.

\subsubsection{Polarization by removing all but one wave orientation}

One way to make polarized light is to remove all but one orientation of an
unpolarized beam. A material that does this at visible wavelengths is called
a \emph{polaroid}. It is made of long-chain hydrocarbons that have been
treated with iodine to make them conductive. The molecules are all oriented
in one direction by stretching during the manufacturing process. The
molecules have electrons that can move when light hits them.

Remember that part of the reason light slows down when it enters a material
is because the light get's absorbed by the electrons. But in most materials
the electrons have to stay in the atoms. They just go to a higher shell. But
in polaroid material the electrons can move farther in the long direction of
the molecule. The molecules act like they are little antennas. The
molecules' electrons are driven into harmonic motion along the length of the
molecule. But this means the energy that used to be our light wave gets
turned into kinetic energy of the electrons. This takes energy (and
therefore, light) out of the beam. So the light that is polarized in this
direction stops at the polaroid

But suppose the light is polarized in a direction perpendicular to the long
molecules. Then little electron motion is possible in the short direction of
the molecule. And if the electrons can't move this direction they won't take
up all the light energy. Then the light is passed if it is polarized
perpendicular to the long direction of the molecules. This direction is
called the \emph{transmission axis}.

We can take two pieces of polaroid material to study polarization.\FRAME{dhF%
}{3.7879in}{1.4088in}{0in}{}{}{Figure}{\special{language "Scientific
Word";type "GRAPHIC";maintain-aspect-ratio TRUE;display "USEDEF";valid_file
"T";width 3.7879in;height 1.4088in;depth 0in;original-width
3.7395in;original-height 1.3742in;cropleft "0";croptop "1";cropright
"1";cropbottom "0";tempfilename 'TN9QPD56.wmf';tempfile-properties "XPR";}}

Unpolarized light is initially polarized by the first piece of polaroid
called the \emph{polarizer}. The second piece of polaroid then receives the
light. This piece is called the \emph{analyzer}. If there is an angular
difference in the orientation of the transmission axes of the polarizer and
analyzer, there will be a reduction of light through the system. We expect
that if the transmission axes are separated by $90\unit{%
%TCIMACRO{\U{b0}}%
%BeginExpansion
{{}^\circ}%
%EndExpansion
}$ no light will be seen. If they are separated by $0\unit{%
%TCIMACRO{\U{b0}}%
%BeginExpansion
{{}^\circ}%
%EndExpansion
},$ then there will be a maximum. It is not hard to believe that the
intensity will be given by%
\begin{equation}
I=I_{\max }\cos ^{2}\left( \theta \right)
\end{equation}%
remembering that for waves we have found that the intensity is proportional
to the amplitude squared. This is true for light waves as well (more on this
in Principles of Physics III).

\subsubsection{Polarization by reflection}

If we look at light reflected off of a desk or table through a piece of
polaroid, we can see that at some angles of orientation, the reflection
diminishes or even disappears! Light is often polarized on reflection. Let's
consider a beam of light made of just two polarizations. The actual
polarization might be a combination of these two (like vector components
make up the vector). We will define a plane of incidence. This plane is the
plane of the paper or computer screen. This plane is perpendicular to the
reflective or refractive surface in the figure below.\FRAME{dtbpF}{3.3313in}{%
2.815in}{0pt}{}{}{Figure}{\special{language "Scientific Word";type
"GRAPHIC";maintain-aspect-ratio TRUE;display "USEDEF";valid_file "T";width
3.3313in;height 2.815in;depth 0pt;original-width 3.5544in;original-height
3in;cropleft "0";croptop "1";cropright "1";cropbottom "0";tempfilename
'TN9QPD57.wmf';tempfile-properties "XPR";}}

One of our polarizations is defined as parallel to this plane. This
direction is represented by orange (lighter grey in black and white) arrows
in the figure. The other polarization is perpendicular to the plane of
incidence (the plane of the paper or your device screen). This is
represented by the black dots in the figure. These dots are supposed to look
like arrows coming out of the paper.

When the light reaches the interface between $n_{1}$ and $n_{2}$ it drives
the electrons in the medium into SHM. The perpendicular polarization finds
electrons that are free to move in the perpendicular direction and
re-radiate in that direction. The electron orbitals change shape and
oscillate with the incoming electromagnetic wave.

The parallel ray is also able to excite SHM, but an electromagnetic analysis
tells us that these little \textquotedblleft antennas\textquotedblright\
will not radiate at an angle $90\unit{%
%TCIMACRO{\U{b0}}%
%BeginExpansion
{{}^\circ}%
%EndExpansion
}$ from their excitation direction. And this is something new! Before when
we used Huygens' principle we thought of the scattered waves as little
spheres. But it turns out that the light doesn't go out in all directions
equally. The disturbance for our wave is a moving charged particle like an
electron. If an electron in a wire experiences simple harmonic motion there
will be a wave in the electric field. But think, that wave will go out away
from the electron perpendicular to the electron's motion \FRAME{dtbpF}{%
3.1816in}{1.7279in}{0pt}{}{}{Figure}{\special{language "Scientific
Word";type "GRAPHIC";maintain-aspect-ratio TRUE;display "USEDEF";valid_file
"T";width 3.1816in;height 1.7279in;depth 0pt;original-width
3.2144in;original-height 1.7332in;cropleft "0";croptop "1";cropright
"1";cropbottom "0";tempfilename 'TN9QPD58.wmf';tempfile-properties "XPR";}}%
In the figure we can see the light wave won't travel up or down parallel to
the wire. So, the actual situation with our light striking a surface is a
little more complicated, The pattern of light coming out from the little
atomic oscillators looks more like a doughnut.

\FRAME{dtbpFU}{2.9616in}{1.9842in}{0pt}{\Qcb{Angular dependence of $S$ for a
dipole scatterer.}}{}{Figure}{\special{language "Scientific Word";type
"GRAPHIC";maintain-aspect-ratio TRUE;display "USEDEF";valid_file "T";width
2.9616in;height 1.9842in;depth 0pt;original-width 2.99in;original-height
1.993in;cropleft "0";croptop "1";cropright "1";cropbottom "0";tempfilename
'TN9QPD59.wmf';tempfile-properties "XPR";}}But the point is that it's not
completely spherical and none of the light wave goes up along the direction
of the oscillation.

In the figure you can see that along the antenna axis, the field amplitude
is zero. This means that the wave really does not go that direction. So in
our case, the amount of polarization in the parallel direction decreases
with the angle between the reflected and refracted rays until at $90\unit{%
%TCIMACRO{\U{b0}}%
%BeginExpansion
{{}^\circ}%
%EndExpansion
}$ there is no reflected ray in the parallel direction. \FRAME{dtbpF}{%
4.1035in}{3.4082in}{0pt}{}{}{Figure}{\special{language "Scientific
Word";type "GRAPHIC";maintain-aspect-ratio TRUE;display "USEDEF";valid_file
"T";width 4.1035in;height 3.4082in;depth 0pt;original-width
4.0542in;original-height 3.3624in;cropleft "0";croptop "1";cropright
"1";cropbottom "0";tempfilename 'TN9QPD5A.wmf';tempfile-properties "XPR";}}

The incidence angle that creates an angular difference between the refracted
and reflected rays of $90\unit{%
%TCIMACRO{\U{b0}}%
%BeginExpansion
{{}^\circ}%
%EndExpansion
}$ is called the \emph{Brewster's\ angle} after its discoverer. At this
angle the reflected beam will be completely linearly polarized.

We can predict this angle. Remember Snell's law.%
\[
n_{1}\sin \theta _{1}=n_{2}\sin \theta _{2}
\]%
Let's re-lable the incidence angle $\theta _{1}=$ $\theta _{b}.$ We take $%
n_{1}=1$ and $n_{2}=n$ so%
\[
n=\frac{\sin \theta _{b}}{\sin \theta _{2}}
\]%
Now notice that for Brewster's angle, we have 
\[
\theta _{b}+90\unit{%
%TCIMACRO{\U{b0}}%
%BeginExpansion
{{}^\circ}%
%EndExpansion
}+\theta _{2}=180\unit{%
%TCIMACRO{\U{b0}}%
%BeginExpansion
{{}^\circ}%
%EndExpansion
}
\]%
so%
\[
\theta _{2}=90\unit{%
%TCIMACRO{\U{b0}}%
%BeginExpansion
{{}^\circ}%
%EndExpansion
}-\theta _{b}
\]%
so we have%
\[
n=\frac{\sin \theta _{b}}{\sin \left( 90\unit{%
%TCIMACRO{\U{b0}}%
%BeginExpansion
{{}^\circ}%
%EndExpansion
}-\theta _{b}\right) }
\]%
Ah, but we remember that $\sin \left( 90\unit{%
%TCIMACRO{\U{b0}}%
%BeginExpansion
{{}^\circ}%
%EndExpansion
}-\theta \right) =\cos \left( \theta \right) $ so%
\[
n=\frac{\sin \theta _{b}}{\cos \theta _{b}}
\]%
Again we remember that 
\[
\tan \theta =\frac{\sin \theta }{\cos \theta }
\]%
so%
\begin{equation}
n=\tan \theta _{b}
\end{equation}%
which we can solve for $\theta _{b}.$%
\[
\theta _{b}=\tan ^{-1}\left( n\right)
\]

This phenomena is why we wear polarizing sunglasses to reduce glare.

\subsubsection{Birefringence}

Glass is an amorphic solid--that is--it has no crystal structure to speak
of. But some minerals do have definite order. Sometimes the difference in
the crystal structure creates a difference in the speed of propagation of
light in the crystal. This is not to hard to believe. We said before that
the reason light slows down in a substance is because it encounters atoms
which absorb and re-emit the light. If there are more atoms in one direction
than another in a crystal, it makes sense that there could be a different
speed in each direction.

Calcite crystals exhibit this phenomena. We can describe what happens by
defining two polarizations. One parallel to the plane of the figure below,
and one perpendicular.

\bigskip

\FRAME{dhF}{3.1337in}{2.0117in}{0in}{}{}{Figure}{\special{language
"Scientific Word";type "GRAPHIC";maintain-aspect-ratio TRUE;display
"USEDEF";valid_file "T";width 3.1337in;height 2.0117in;depth
0in;original-width 3.1647in;original-height 2.0214in;cropleft "0";croptop
"1";cropright "1";cropbottom "0";tempfilename
'TN9QPD5B.wmf';tempfile-properties "XPR";}}With a careful setup, we can
arrange things so the perpendicular ray is propagated just as we would
expect for glass. We call this the $O$-ray (for \emph{ordinary}). The second
ray is polarized parallel to the incidence plane. It will have a different
speed, and therefore a different index of refraction. We call it the \emph{%
Extraordinary ray} or $E$-ray.

\FRAME{dhF}{2.1137in}{2.1234in}{0in}{}{}{Figure}{\special{language
"Scientific Word";type "GRAPHIC";maintain-aspect-ratio TRUE;display
"USEDEF";valid_file "T";width 2.1137in;height 2.1234in;depth
0in;original-width 2.1252in;original-height 2.135in;cropleft "0";croptop
"1";cropright "1";cropbottom "0";tempfilename
'TN9QPD5C.wmf';tempfile-properties "XPR";}}

If we were to put a light source in a calcite crystal, we would see the $O$%
-ray send out a sphere of light as shown in the figure above. But the $E$%
-ray would send out an ellipse. The speed for the $E$-ray depends on
orientation. There is one direction where the speeds are equal. This
direction is called the \emph{optic axis} of the crystal.\FRAME{dtbpF}{%
3.2206in}{1.8299in}{0in}{}{}{Figure}{\special{language "Scientific
Word";type "GRAPHIC";maintain-aspect-ratio TRUE;display "USEDEF";valid_file
"T";width 3.2206in;height 1.8299in;depth 0in;original-width
3.1756in;original-height 1.7928in;cropleft "0";croptop "1";cropright
"1";cropbottom "0";tempfilename 'TN9QPD5D.wmf';tempfile-properties "XPR";}}%
If our light entering our calcite crystal is unpolarized, then we will have
two images leaving the other side that are slightly offset because the $O$%
-rays and $E$-rays both form images.

Some materials (notably plastics) become birefringent under stress. A
plastic or other stress birefringent material is molded in the form planned
for a building or other object (usually made to scale). The model is placed
under a stress, and the system is placed between to polaroids. When
unstressed, no light is seen, but under stress, the model changes the
polarization state of the light, and bands of light are seen.\FRAME{dhF}{%
3.1479in}{2.0676in}{0pt}{}{}{Figure}{\special{language "Scientific
Word";type "GRAPHIC";maintain-aspect-ratio TRUE;display "USEDEF";valid_file
"T";width 3.1479in;height 2.0676in;depth 0pt;original-width
3.1798in;original-height 2.0782in;cropleft "0";croptop "1";cropright
"1";cropbottom "0";tempfilename 'TN9QPD5E.wmf';tempfile-properties "XPR";}}

\subsubsection{Polarization due to scattering}

It is important to understand that light is also polarized by scattering. It
really takes a bit of electromagnetic theory to describe this. So for a
moment, lets just comment that blue light is scattered more than red light
because the blue wavelength is about the same size as the molecules. In
fact, the relative intensity of scattered light goes like $1/\lambda ^{4}.$
This has nothing to do with polarization, but it is nice to know.

Now suppose we have long pieces of wire in the air, say, a few microns long.
The pieces of wire would have electrons that could be driven into SHM when
light hits them in a similar manner to our polaroid molecules. If the wires
were all oriented in a common direction, we would expect light to be
absorbed if it was polarized in the long direction of the particles and not
absorbed in a direction perpendicular to the orientation of the particles.
This is exactly what happens when long ice particles in the atmosphere
orient in the wind (think of the moment of inertia). We often get impressive
halo's around the sun due to scattering from ice particles. And these halos
can be strongly polarized.

Rain drops also have a preferential scattering direction because they are
shaped like oblate spheroids (not \textquotedblleft drop
shape\textquotedblright\ like we were told in grade school). This makes
weather radar signals polarized and this lets us measure how much water
there is in a rain storm!

It is also true that small molecules will act like tiny antennas and will
scatter light preferentially in some directions and not in others. This is
called \emph{Rayleigh scattering} and is very like small dipole antennas.

\subsubsection{Optical Activity}

Some substances will rotate the polarization of a beam of light. This is
called being \emph{optically active}. The polarization state of the light
exiting the material depends on the length of the path through the material.
Your calculator display works this way. An electric field changes the
optical activity of the liquid crystal. There are polarizers over the liquid
crystal, so sometimes light passes through the display and sometimes it is
black.

\subsubsection{Laser polarization}

One last comment. Lasers are usually polarized. This is because the laser
light is generated in a \emph{cavity} created by two mirrors. The mirror is
tipped so light approaches it at the Brewster angle. Light with the right
polarization (parallel to the plane of the drawing) is reflected back nearly
completely, but light with the opposite polarization is not reflected at
all. Again, this is because of the little atomic scatterers and
superposition. This reduces the usual loss in reflection from a mirror,
because in one polarization the light must be reflected completely.

\FRAME{dtbpF}{4.2653in}{1.03in}{0pt}{}{}{Figure}{\special{language
"Scientific Word";type "GRAPHIC";maintain-aspect-ratio TRUE;display
"USEDEF";valid_file "T";width 4.2653in;height 1.03in;depth
0pt;original-width 4.2151in;original-height 0.9963in;cropleft "0";croptop
"1";cropright "1";cropbottom "0";tempfilename
'TN9QPD5F.wmf';tempfile-properties "XPR";}}

\subsection{Imaging with Flat Mirrors}

All of us have looked in a mirror at some time. We know what to expect. We
see an image of ourselves. And we just thought about fancy mirrors that we
use to make lasers. But we need to apply math to our study of mirrors so we
can calculate how light is reflected by mirrors.

To study mirrors mathematically, we need to establish a sign convention and
some standard notation.

\FRAME{dtbpF}{1.5463in}{1.6267in}{0pt}{}{}{Figure}{\special{language
"Scientific Word";type "GRAPHIC";maintain-aspect-ratio TRUE;display
"USEDEF";valid_file "T";width 1.5463in;height 1.6267in;depth
0pt;original-width 1.5108in;original-height 1.5904in;cropleft "0";croptop
"1";cropright "1";cropbottom "0";tempfilename
'TN9QPD5G.wmf';tempfile-properties "XPR";}}In the figure above, we have a
person observing an object $O$ in a mirror. The object is located at a
distance $s$ from the mirror. Just like for lenses, we will call this the 
\emph{object distance, }$d_{o}$. The image appears to be located at a point $%
I$ beyond the mirror a distance $d_{i}.$We will call this the \emph{image
distance}. Be careful, $d_{i}$ looks like an incident or initial distance,
but it's not. It is the distance from the mirror to the \emph{image.} We
know we see an image in the mirror, but to understand this we need to define
what we mean by the word \textquotedblleft image.\textquotedblright

\subsubsection{Images}

Let's think about what an image is. Take a piece of paper and a lens, and
hold up the lens in a darkened room that has some bright object in it. Move
the lens or the paper back and forth, and at just the right distance, a
miniature image of the bright object will appear. We should think about what
the word \textquotedblleft image\textquotedblright\ means in this sense.

\FRAME{dhF}{2.6835in}{2.0237in}{0pt}{}{}{Figure}{\special{language
"Scientific Word";type "GRAPHIC";maintain-aspect-ratio TRUE;display
"USEDEF";valid_file "T";width 2.6835in;height 2.0237in;depth
0pt;original-width 2.6411in;original-height 1.9847in;cropleft "0";croptop
"1";cropright "1";cropbottom "0";tempfilename
'TN9QPD5H.wmf';tempfile-properties "XPR";}}We have talked about how we see
objects. Remember the BYU-I guys from a previous lecture.\FRAME{dhF}{3.173in%
}{1.3534in}{0pt}{}{}{Figure}{\special{language "Scientific Word";type
"GRAPHIC";maintain-aspect-ratio TRUE;display "USEDEF";valid_file "T";width
3.173in;height 1.3534in;depth 0pt;original-width 3.128in;original-height
1.3188in;cropleft "0";croptop "1";cropright "1";cropbottom "0";tempfilename
'TN9QPD5I.wmf';tempfile-properties "XPR";}}Our eyes gather rays that are
diverging from the object because light has bounced off of the object. Our
eyes intersect a diverging set of rays that form a definite pattern. That
diverging set of rays forming a pattern is the picture that we see of the
object.

So when we say that the lens has formed a miniature image of our object, we
mean that the lens has somehow formed a diverging set of rays that form a
pattern that looks like the pattern formed by the diverging set of rays
coming from the object, itself. In other words, the object forms a diverging
set of rays. And our lens forms a duplicate set of rays in the same pattern,
so we see the same thing. The lens' version is smaller, upside down, and
backwards, but it is still essentially the same pattern.

But how can a mirror form a diverging set of rays that form a pattern?

Well, the light that we see in a mirror first bounces off of something, say,
a person. This makes a recognizable pattern in the light. Then the light
reflects off of the mirror. Then the light hits our eyes. So when we see
something in a mirror we do see a diverging set of rays that form a pattern.
But what we see makes more makes sense if we know how our brains processes
the light signals that we see. Our eyes intercept rays of light diverging
from an object. Our brain processes those rays as though they traveled only
in straight lines. Our brain really believes light travels just in lines, no
bouncing or reflecting. This means, if we can create a situation that makes
rays diverge in the same way the object did, we will have an image of the
object. But our brain will be wrong in where it thinks that image is because
it does not know about the reflection. In this case, we will see an image as
though it were behind the mirror.

\subsubsection{Virtual Images}

We already know that mirrors often create what we call \emph{virtual} images
because the image \emph{appears} to be created from diverging rays from
behind the mirror. If we look behind the mirror, no rays exist there (if
they did exist, they would not make it through the mirror!).

\FRAME{dtbpF}{2.3869in}{1.4183in}{0in}{}{}{Figure}{\special{language
"Scientific Word";type "GRAPHIC";maintain-aspect-ratio TRUE;display
"USEDEF";valid_file "T";width 2.3869in;height 1.4183in;depth
0in;original-width 2.3471in;original-height 1.3837in;cropleft "0";croptop
"1";cropright "1";cropbottom "0";tempfilename
'TN9QPD5J.wmf';tempfile-properties "XPR";}}

Let's look at a simple image as shown in the figure above. The object is an
arrow. We could trace all the rays that diverge from this object and build a
very nice representation of the arrow (like ray tracing-based computer
graphics--the way movies like \emph{Toy Story} are made) but that would take
time and computation power. We only really need to use two rays, and to
remember what the object looked like to figure out where the image will be.

We pick one ray from the top of the arrow that travels straight to the
mirror. This ray will travel a distance $d_{o}$ and bounce back. Using the
law of reflection we know that it will bounce straight back along the path
it came in on. We pick a second ray from point $P$ that travels the path $%
PR. $ This ray bounces off the mirror at an angle $\theta _{r}=\theta _{i}.$
But our brain optical processing center doesn't believe the light reflected.
It thinks the light traveled straight paths. It mentally traces the rays
backwards along a straight path as though it came from that direction. And
this makes it appear that there is a tip of an arrow is at position $%
P^{\prime }.$ The rays from the virtual tip appear to travel the paths $%
P^{\prime }P$ and $P^{\prime }R.$ The light doesn't really come from point $%
P^{\prime }.$ But it \emph{appears} to come from point $P^{\prime }$. We
could use similar sets of two rays to find where the light from the middle
of the arrow \emph{appears} to come from, and for the bottom of the arrow,
and all the arrow parts. The collection of all these apparent locations is
image location and the light that appears to come from this location is the
image. But, and this is important! light doesn't come from $P^{\prime }.$ It
only appears to come from $P^{\prime }.$ When an image is formed, but the
light doesn't really come from the image location, we call that a \emph{%
virtual image.}

\subsubsection{Mirror reversal}

Look into a mirror. Raise your right hand. \FRAME{dtbpF}{3.173in}{1.9147in}{%
0pt}{}{}{Figure}{\special{language "Scientific Word";type
"GRAPHIC";maintain-aspect-ratio TRUE;display "USEDEF";valid_file "T";width
3.173in;height 1.9147in;depth 0pt;original-width 3.1289in;original-height
1.8766in;cropleft "0";croptop "1";cropright "1";cropbottom "0";tempfilename
'TN9QPD5K.wmf';tempfile-properties "XPR";}}Your image raises what appears to
be a left hand. It looks like a mirror switches the left and right sides of
the image. But it might be better to note that the image raised hand is on
the same side of the room as your real raised hand. Think of lying sideways
on the ground in front of the mirror and raise a hand. Note that your image
feet are on the same side of the room as your real feet. And note that as
you put one hand up, the image hand also goes up. This is not exactly
left-right reversal. What is happening? \FRAME{dtbpF}{4.3716in}{2.6333in}{0pt%
}{}{}{Figure}{\special{language "Scientific Word";type
"GRAPHIC";maintain-aspect-ratio TRUE;display "USEDEF";valid_file "T";width
4.3716in;height 2.6333in;depth 0pt;original-width 4.3206in;original-height
2.5918in;cropleft "0";croptop "1";cropright "1";cropbottom "0";tempfilename
'TN9QPD5L.wmf';tempfile-properties "XPR";}}The light from your raised hand
would hit the mirror, reflect, and then be absorbed by your eye. Your brain
tells you the light traveled in a straight line. So it looks like you have a
raised hand in the mirror. But what has happened is that the light has
traveled from your hand, to the mirror and back to your eye. This is not
really left-right reversal. We need a name to describe what is really
happening. The name is odd. To see why we use this name, notice that from
your perspective (moving left to right across the figure) you first have the
back part of your hand, then the front of your hand and then the front of
the image hand, then the back of the image hand. We call this back-front
reversal. It comes from the reversal in the order of parts of the object,
like the front and back of your hand. So we say that a flat mirror performs
a front-back reversal, not a left-right reversal.

It is nice to know about how images are formed by mirrors. We use them all
the time. But you probably own a phone, and your phone probably has a
camera. Cameras also make images. We should work our way up to understanding
how things like cameras work. We will start that in our next lecture.

\section{Images formed by Curved Mirrors and by Refraction 3.2.2, 3.2.3}

Back in Newton's day, there were no electronic computers, or calculators, or
mechanical adding machines. Early optics researchers did math with a pen and
paper. This is one reason they liked small angle approximations. The
approximation allowed them to do harder problems using easy math. And so
long as the things they built worked, the approximations were good enough.
We are going to use another approximation in this lecture. It is called the
thin mirror approximation. It will make the math that describes mirrors for,
say, telescopes, much easier. We will also introduce a way to draw curved
mirrors and light that will tell you how the mirror works. We will use this
drawing approach to do homework and examination problems. Professional
optical designers use computer applications that do this drawing very
accurately to help the optical engineer understand the system they are
designing and to look for mistakes. So this drawing scheme is very useful.
Because it makes optical systems so much easier to understand, let's start
with the drawings and then work toward a mathematical description of thin
mirrors.

%TCIMACRO{%
%\TeXButton{Fundamental Concepts}{\vspace{0.10in}\hspace{-0.37in}{\Large Fundamental Concepts\vspace{0.2in}}}}%
%BeginExpansion
\vspace{0.10in}\hspace{-0.37in}{\Large Fundamental Concepts\vspace{0.2in}}%
%EndExpansion

\begin{itemize}
\item Curved Mirrors can form images

\item We can write the magnification in terms of the object and image
distances $m=-\frac{d_{i}}{d_{o}}$

\item For \textquotedblleft thin\textquotedblright\ mirrors $\frac{1}{d_{o}}+%
\frac{1}{d_{i}}=\frac{1}{f}$

\item The focal length of a \textquotedblleft thin\textquotedblright\ mirror
is. $f=\frac{R}{2}$

\item We can describe how a mirror operates with just three easy-to-draw
rays.

\item A semi-infinite bump (single sided lens) can be described by the
equation $\frac{1}{d_{o}}+\frac{1}{d_{i}}=\frac{\left( n_{2}-n_{1}\right) }{R%
}$
\end{itemize}

\subsection{Concave Mirrors}

Concave mirrors can form images. I'm sure you know that many telescopes are
made with mirrors. We should see how this works. We recall the law of
reflection\FRAME{dhF}{1.8585in}{1.1857in}{0pt}{}{}{Figure}{\special{language
"Scientific Word";type "GRAPHIC";maintain-aspect-ratio TRUE;display
"USEDEF";valid_file "T";width 1.8585in;height 1.1857in;depth
0pt;original-width 1.8213in;original-height 1.1519in;cropleft "0";croptop
"1";cropright "1";cropbottom "0";tempfilename
'TN9QPD5M.wmf';tempfile-properties "XPR";}}

\[
\theta _{i}=\theta _{r}
\]

Armed with this, we can see what would happen if we curved the mirror
surface. Each ray has a different normal line due to the curvature of the
mirror. \FRAME{dtbpF}{0.8458in}{1.9908in}{0pt}{}{}{Figure}{\special{language
"Scientific Word";type "GRAPHIC";maintain-aspect-ratio TRUE;display
"USEDEF";valid_file "T";width 0.8458in;height 1.9908in;depth
0pt;original-width 0.8138in;original-height 1.9519in;cropleft "0";croptop
"1";cropright "1";cropbottom "0";tempfilename
'TN9QPD5N.wmf';tempfile-properties "XPR";}}And the law of reflection must
works at each one of these normal lines, so we get rays that reflect like
this. \FRAME{dtbpF}{3.0649in}{1.9666in}{0pt}{}{}{Figure}{\special{language
"Scientific Word";type "GRAPHIC";maintain-aspect-ratio TRUE;display
"USEDEF";valid_file "T";width 3.0649in;height 1.9666in;depth
0pt;original-width 2.5806in;original-height 1.6466in;cropleft "0";croptop
"1";cropright "1";cropbottom "0";tempfilename
'TN9QPD5O.wmf';tempfile-properties "XPR";}}The result is that parallel rays
all meet at a spot on the axis.

\FRAME{dhF}{3.243in}{2.079in}{0pt}{}{}{Figure}{\special{language "Scientific
Word";type "GRAPHIC";maintain-aspect-ratio TRUE;display "USEDEF";valid_file
"T";width 3.243in;height 2.079in;depth 0pt;original-width
3.1981in;original-height 2.0401in;cropleft "0";croptop "1";cropright
"1";cropbottom "0";tempfilename 'TN9QPD5P.wmf';tempfile-properties "XPR";}}%
If we place something at this location (like a pice of paper), we would see
it light up. With enough light we could start a fire! We have just
re-invented the solar cooker. Look at the point where all the rays meet.
Let's give this special place a name. We call it the focal point. It is the
place where the rays come together if they originally came into the mirror
parallel.

\subsection{Paraxial Approximation for Mirrors}

The correct shape of a mirror to make all the parallel rays meet at the
focal point is more like a parabola, but parabolas are hard to machine or
build. Spherical shapes are relatively easy to make. So we often see
spherical mirrors used where we should have parabolas. This will work so
long as we allow only rays that strike the mirror not too far from the
principal axis. We can see why this works if we plot a sphere and a parabola
(and a hyperbola). For small deviations from the center, the shape of the
functions all look alike. \FRAME{dhF}{4.3613in}{1.1628in}{0pt}{}{}{Figure}{%
\special{language "Scientific Word";type "GRAPHIC";display
"USEDEF";valid_file "T";width 4.3613in;height 1.1628in;depth
0pt;original-width 4.3102in;original-height 1.3604in;cropleft "0";croptop
"1";cropright "1";cropbottom "0";tempfilename
'TN9QPD5Q.wmf';tempfile-properties "XPR";}}We would expect the reflections
to be similar under these circumstances. Using only the part of a mirror
where the spherical and parabolic shapes are essentially the same is called
the \emph{paraxial approximation}. So, if we meet the criteria for the
paraxial approximation, our spherical mirrors should work. Note that when
you need the entire mirror, say, in a communications antenna, you must do
better than a spherical approximation to the correct shape for your mirror.
Your satellite dish is likely not a spherical section. We studied flat
mirrors last lecture, and found they made images (well, virtual images).
Will curved mirrors make images? We know how to find out. We can trace two
rays from an object and see where the rays go. If the rays eventually
diverge after forming a recognizable pattern, then we have an image!\FRAME{%
dtbpF}{2.88in}{1.7917in}{0in}{}{}{Figure}{\special{language "Scientific
Word";type "GRAPHIC";maintain-aspect-ratio TRUE;display "USEDEF";valid_file
"T";width 2.88in;height 1.7917in;depth 0in;original-width
2.9057in;original-height 1.797in;cropleft "0";croptop "1";cropright
"1";cropbottom "0";tempfilename 'TN9QPD5R.wmf';tempfile-properties "XPR";}}

Like with our flat mirror, we will measure distances from the mirror surface
(from point $V$ in the figure). We can find the image location, $d_{i},$ by
taking two rays. We could use any of billions of rays. But let's try to pick
rays that are easy to draw. One convenient ray is the ray that passes
through the center of curvature, $C.$ This ray will strike the mirror
surface at right angles to the surface and bounce back along the same path.
The incidence angle will be zero, so the reflected angle must be zero by the
law of reflection. That is the orange ray in the figure.

Another convenient ray is the ray from the tip of the object to point $V.$
This ray will bounce back with angle $\theta $. Right at point $V$ the
mirror surface is perpendicular to the optic axis. This makes the bounce of
the blue ray just like the a bounce from a flat mirror! That is easy to
draw. Where these two reflected rays cross and then diverge, we will find
the image of the tip of our tree. Knowing the shape of the tree and that the
bottom of the tree is on the axis, we can fill in the rest of the image.

We can calculate the magnification for this case. We use the gold triangle
to determine that 
\[
\tan \theta =\frac{h_{o}}{d_{o}}
\]%
and the light blue triangle to determine that 
\[
\tan \theta =\frac{-h_{i}}{d_{i}}
\]%
And you might object to the negative sigh. But we want to indicate that the
image is inverted by making it's sign negative. So because the image is
upside down $h_{i}$ is defined to be negative. To make our tangent equation
work, we need to insert another minus sign. This gives us $\tan \theta $
instead of $-\tan \theta $ when $h^{\prime }$ is inverted (when $h^{\prime
}, $ itself, is negative). When we decide to make the sine of a quantity
negative or positive just because we want it to be that way we say it is
negative (or positive) \emph{by convention.} And we have done this here. It
is just a convention. But we will use it. Our magnification would be%
\[
m\equiv \frac{\text{Image height}}{\text{Object height}}=\frac{h_{i}}{h_{o}}
\]%
But we can write 
\begin{eqnarray*}
h_{o} &=&d_{o}\tan \theta \\
h_{i} &=&-d_{i}\tan \theta
\end{eqnarray*}%
so the magnification could be written as 
\[
m=\frac{h_{i}}{h_{o}}=\frac{-d_{i}\tan \theta }{d_{o}\tan \theta }=-\frac{%
d_{i}}{d_{o}}
\]%
which gives us the magnification in terms of how far away the object and
image are.

\subsection{Mirror Equation}

We can further exploit this geometry to get a relationship between $d_{o},$ $%
d_{i},$ and $R.$ We would really like to know where the image is, $d_{i}$,
in terms of the things we know $d_{o}$ and $R.$ Notice that 
\[
\tan \alpha =\frac{h_{o}}{d_{o}-R}
\]%
and that 
\[
\tan \alpha =\frac{-h_{i}}{R-d_{i}}
\]

Then%
\[
\frac{h_{o}}{d_{o}-R}=\frac{-h_{i}}{R-d_{i}}
\]%
or%
\[
\frac{R-d_{i}}{d_{o}-R}=-\frac{h_{i}}{h_{o}}
\]

We can use our magnification definition to replace $h_{i}/h_{o}$%
\[
\frac{R-d_{i}}{d_{o}-R}=\frac{d_{i}}{d_{o}}
\]%
And we have a useful equation, but not in the form you usually see it. Let's
perform some algebra. Starting with our last equation, let's try to make is
so there are no fractions.%
\begin{eqnarray*}
\left( R-d_{i}\right) d_{o} &=&d_{i}\left( d_{o}-R\right) \\
Rd_{o}-d_{i}d_{o} &=&d_{i}d_{o}-d_{i}R
\end{eqnarray*}%
And now let's get all the terms with an $R$ in them on the same side of the
equation. 
\begin{eqnarray*}
Rd_{o}+d_{i}R &=&d_{i}d_{o}+d_{i}d_{o} \\
d_{i}R+Rd_{o} &=&2d_{o}d_{i}
\end{eqnarray*}%
Now, for no apparent reason, I want to divide both sides of the equation by $%
Rd_{o}d_{i}.$ That might not seem reasonable, but let's do it anyway. 
\[
\frac{Rd_{i}}{Rd_{o}d_{i}}+\frac{Rd_{o}}{Rd_{o}d_{i}}=\frac{2d_{o}d_{i}}{%
Rd_{o}d_{i}}
\]%
Some things cancel, and we are left with 
\[
\frac{1}{d_{o}}+\frac{1}{d_{i}}=\frac{2}{R}
\]%
which technically completes our goal. We could solve for $d_{i}.$ We are
close to the usual form for this equation. But let's make one more change to
this equation. And it will involve the focal point.

\subsubsection{Focal Point}

Remember that special place where, if the light comes into the mirror
parallel, the reflected rays all meet? We called that the focal point, $f$.

Now that we know the mirror equation, let's let $d_{o}$ be very large (for
example, let $d_{o}$ be the distance to the Sun). Then the rays would come
very nearly parallel. And then after reflection all the rays would come
together at the focal point. They would diverge from the focal point. We
have a diverging set of rays that are shaped like the sun! We would get an
image of the Sun right at the focal point! Since the sun is approximately
infinitely far away, we could say $d_{o}\approx \infty .$ Then, using our
last equation, 
\begin{eqnarray*}
\frac{1}{\infty }+\frac{1}{d_{i}} &\approx &\frac{2}{R} \\
0+\frac{1}{d_{i}} &\approx &\frac{2}{R} \\
\frac{1}{d_{i}} &\approx &\frac{2}{R}
\end{eqnarray*}%
or 
\[
d_{i}\approx \frac{R}{2}
\]

This is a special image point. But we really already know what to call this
special place where parallel rays come together after reflection. We call it
the \emph{focal point} and the distance from the mirror to the point $f\ $is
called the \emph{focal length}. We see that for a curved mirror the length $%
f $ is given by 
\begin{equation}
f=\frac{R}{2}
\end{equation}%
So we can write the mirror equation as%
\begin{equation}
\frac{1}{d_{o}}+\frac{1}{d_{i}}=\frac{1}{f}
\end{equation}

This equation is very useful. Even when you use powerful optical design
software it can form the beginning of your optical system design.

\subsection{Ray Diagrams for Mirrors}

For our mirror problems we will need to draw a diagram of the situation. And
often the secret to solving the problem is in making the diagram! So, let's
learn how to do this. We could draw an infinite number of rays to find where
the image will be. And, indeed, optical design software does draw many rays
to check our designs. But for us we need to draw just a few easy rays. But
what rays are easy to draw?

Let's think about what we already know. \FRAME{dtbpF}{2.9386in}{1.224in}{0pt%
}{}{}{Figure}{\special{language "Scientific Word";type
"GRAPHIC";maintain-aspect-ratio TRUE;display "USEDEF";valid_file "T";width
2.9386in;height 1.224in;depth 0pt;original-width 2.9661in;original-height
1.2187in;cropleft "0";croptop "1";cropright "1";cropbottom "0";tempfilename
'TN9QPD5S.wmf';tempfile-properties "XPR";}}If parallel rays come into the
mirror, they meet at the focal point. They don't really stop there unless we
put something there.

\FRAME{dtbpF}{3.2907in}{1.3509in}{0pt}{}{}{Figure}{\special{language
"Scientific Word";type "GRAPHIC";maintain-aspect-ratio TRUE;display
"USEDEF";valid_file "T";width 3.2907in;height 1.3509in;depth
0pt;original-width 3.3244in;original-height 1.3473in;cropleft "0";croptop
"1";cropright "1";cropbottom "0";tempfilename
'TN9QPD5T.wmf';tempfile-properties "XPR";}}But often we will just stop our
drawing at the point $f.$ This is easy to draw! Just a parallel line in and
bounce through the focal point, $f$. And we know the focal distance $f=R/2$
so it is even easy to find the point, $f.$ This is one way to make a solar
cooker. \FRAME{dtbpF}{2.866in}{2.1577in}{0pt}{}{}{Figure}{\special{language
"Scientific Word";type "GRAPHIC";maintain-aspect-ratio TRUE;display
"USEDEF";valid_file "T";width 2.866in;height 2.1577in;depth
0pt;original-width 1.8386in;original-height 1.3785in;cropleft "0";croptop
"1";cropright "1";cropbottom "0";tempfilename
'SolarCooker.wmf';tempfile-properties "XNPR";}}so we could call any ray that
comes in parallel and goes through the focal point a \textquotedblleft solar
cooker ray.\textquotedblright\ Note that it is a little weird to have a
point marked $f$ and a distance from a mirror to that point named $f.$ But
we will get used to it.

It turns out we can use this geometry twice. Think, the law of reflection
tells us that every ray that bounced off of the mirror did so with equal
angles $\theta _{i}=\theta _{r}.$ So what would happen if we put a light
source at the focal point? The law of reflection would work the same way,
and the rays would come out parallel! \FRAME{dtbpF}{3.9125in}{1.6356in}{0pt}{%
}{}{Figure}{\special{language "Scientific Word";type
"GRAPHIC";maintain-aspect-ratio TRUE;display "USEDEF";valid_file "T";width
3.9125in;height 1.6356in;depth 0pt;original-width 3.9595in;original-height
1.6383in;cropleft "0";croptop "1";cropright "1";cropbottom "0";tempfilename
'TN9QPD5U.wmf';tempfile-properties "XPR";}}This would be easy to draw! And
this is also useful. Old fashioned flashlights used this to make the light
from a bulb that came out in all directions go in just one direction. Big
versions of this were used to send light out parallel to find enemy aircraft
in World Ward II. These devices were called \emph{spot lights} so we can
call rays that act like this \textquotedblleft spot light\textquotedblright\
rays.

\FRAME{dtbpF}{2.4871in}{2.4002in}{0pt}{}{}{Figure}{\special{language
"Scientific Word";type "GRAPHIC";maintain-aspect-ratio TRUE;display
"USEDEF";valid_file "T";width 2.4871in;height 2.4002in;depth
0pt;original-width 2.5066in;original-height 2.417in;cropleft "0";croptop
"1";cropright "1";cropbottom "0";tempfilename
'TN9QPD5V.wmf';tempfile-properties "XPR";}}Now days, you might only see such
a device used in movies about your favorite super hero (Murci\'{e}lago Boy)
used to call on the super hero for help.

There is one more ray that is easy to draw. We have already encountered it.
It is the ray that goes through the center of curvature for the spherical
mirror so it hits with $\theta _{i}=0$ and bounces off with $\theta _{r}=0.$

Let's put all of these together to dray ray diagrams for a spherical mirror.

\subsubsection{Principal rays for a concave mirror:}

\begin{enumerate}
\item Ray 1 is drawn from the top of the object such that its reflected ray
must pass through $f$. \FRAME{dtbpF}{2.8729in}{1.3739in}{0pt}{}{}{Figure}{%
\special{language "Scientific Word";type "GRAPHIC";maintain-aspect-ratio
TRUE;display "USEDEF";valid_file "T";width 2.8729in;height 1.3739in;depth
0pt;original-width 2.8987in;original-height 1.3713in;cropleft "0";croptop
"1";cropright "1";cropbottom "0";tempfilename
'TN9QPD5W.wmf';tempfile-properties "XPR";}}This is the \textquotedblleft
solar cooker\textquotedblright\ ray.

\item Ray 2 is drawn from the top of the object through the center of
curvature. This ray will be incident on the mirror surface at a right angle
and will be reflected back on itself. \FRAME{dtbpF}{2.9457in}{1.4076in}{0pt}{%
}{}{Figure}{\special{language "Scientific Word";type
"GRAPHIC";maintain-aspect-ratio TRUE;display "USEDEF";valid_file "T";width
2.9457in;height 1.4076in;depth 0pt;original-width 2.9732in;original-height
1.405in;cropleft "0";croptop "1";cropright "1";cropbottom "0";tempfilename
'TN9QPD5X.wmf';tempfile-properties "XPR";}}This is the \textquotedblleft
center of curvature\textquotedblright\ ray.

\item Ray 3 is drawn from the top of the object through the focal point to
reflect parallel to the principal axis.\FRAME{dtbpF}{2.9634in}{1.8103in}{0pt%
}{}{}{Figure}{\special{language "Scientific Word";type
"GRAPHIC";maintain-aspect-ratio TRUE;display "USEDEF";valid_file "T";width
2.9634in;height 1.8103in;depth 0pt;original-width 2.9918in;original-height
1.8157in;cropleft "0";croptop "1";cropright "1";cropbottom "0";tempfilename
'TN9QPD5Y.wmf';tempfile-properties "XPR";}}This is the \textquotedblleft
spot light\textquotedblright\ ray.
\end{enumerate}

We can put these together, and where they meet will be where the image will
form. \FRAME{dtbpF}{3.6765in}{2.2423in}{0in}{}{}{Figure}{\special{language
"Scientific Word";type "GRAPHIC";maintain-aspect-ratio TRUE;display
"USEDEF";valid_file "T";width 3.6765in;height 2.2423in;depth
0in;original-width 3.7182in;original-height 2.2565in;cropleft "0";croptop
"1";cropright "1";cropbottom "0";tempfilename
'TN9QPD5Z.wmf';tempfile-properties "XPR";}} Of course, we did this for the
top of the tree. We could do the same for the middle of the tree and we
would get \FRAME{dtbpF}{3.7404in}{2.2866in}{0in}{}{}{Figure}{\special%
{language "Scientific Word";type "GRAPHIC";maintain-aspect-ratio
TRUE;display "USEDEF";valid_file "T";width 3.7404in;height 2.2866in;depth
0in;original-width 3.783in;original-height 2.3008in;cropleft "0";croptop
"1";cropright "1";cropbottom "0";tempfilename
'TN9QPD60.wmf';tempfile-properties "XPR";}}and we can do this for every
point on the tree and find the corresponding image point. But it turns out
that if we do the top of the object we pretty much know where the rest of
the image will be. So with just three rays we can find the image of our
object!

We can do the same for an object closer than a focal length\FRAME{dtbpF}{%
2.847in}{1.8542in}{0pt}{}{}{Figure}{\special{language "Scientific Word";type
"GRAPHIC";maintain-aspect-ratio TRUE;display "USEDEF";valid_file "T";width
2.847in;height 1.8542in;depth 0pt;original-width 3.3425in;original-height
2.1664in;cropleft "0";croptop "1";cropright "1";cropbottom "0";tempfilename
'TN9QPD61.wmf';tempfile-properties "XPR";}}Only the spotlight ray doesn't
work because light that came from the top of the tree and passed through the
focal point won't hit the mirror. But we could draw a ray that goes along
the line from the focal point to the top of the tree and then hits the
mirror. That ray will hit the mirror and bounce just as though it really
came from the focal point. We usually use dotted lines to show when we are
pretending a ray exists somewhere where it really doesn't. So in the drawing
you see a dotted line come from the focal point, $f,$ and then the line
becomes solid blue where the light really starts at the top of the tree and
goes to the mirror and reflects. In this way we can still use our three ray
system to find the image.

But wait. In this case our three rays never meet! So there seems to be no
image! An it is true, there is no real image. But we learned from flat
mirrors that there can be a virtual image. Virtual images come from our
brain processing thinking that the rays only go in straight lines. So we can
extend the rays that leave the mirror backward as though they didn't bounce
at the mirror. Again we use dotted lines for this. And where the dotted
lines meet we have a virtual image!

\subsection{Principal rays for a convex mirror:}

We also may have a mirror that curves, but curves the other way. We can use
three rays, but this time we need to be careful.

\begin{enumerate}
\item Ray 1 is drawn from the top of the object such that its reflected ray
appears to have come from $f$.

\item Ray 2 is drawn from the top of the object so that it appears to have
come from the center of curvature. This ray will be incident on the mirror
surface at a right angle and will be reflected back on itself.

\item Ray 3 is drawn from the top of the object to reflect parallel to the
principal axis by aiming it at $f$.
\end{enumerate}

\FRAME{dtbpF}{2.7804in}{2.2009in}{0pt}{}{}{Figure}{\special{language
"Scientific Word";type "GRAPHIC";maintain-aspect-ratio TRUE;display
"USEDEF";valid_file "T";width 2.7804in;height 2.2009in;depth
0pt;original-width 2.2719in;original-height 1.7928in;cropleft "0";croptop
"1";cropright "1";cropbottom "0";tempfilename
'TN9QPD62.wmf';tempfile-properties "XPR";}}Note again that these rays won't
meet. So there is no real image. But we can extend the rays that leave the
mirror backward as though they hadn't bounced off the mirror we can see that
they do meet, and that is the location our brain assigns for the virtual
image.

We should tabulate our sign convention for mirrors. So far we know that if $%
h_{i}$ is upside down it will be negative. But we can assign signs for all
the other distances in our ray diagrams. 
\[
\begin{tabular}{|l|l|l|}
\hline
\textbf{Quantity} & \textbf{Positive if} & \textbf{Negative if} \\ \hline
{\small Object location }$\left( d_{o}\right) $ & 
\begin{tabular}{l}
{\small Object is in front of surface} \\ 
{\small (real object)}%
\end{tabular}
& 
\begin{tabular}{l}
{\small Object is in back of surface} \\ 
{\small (virtual object)}%
\end{tabular}%
{\small \ } \\ \hline
{\small Image location }$\left( d_{i}\right) $ & 
\begin{tabular}{l}
{\small Image is in front of surface } \\ 
{\small (real image)}%
\end{tabular}
& 
\begin{tabular}{l}
{\small Image is in back of surface} \\ 
{\small (virtual image)}%
\end{tabular}
\\ \hline
{\small Image height }$\left( h_{i}\right) $ & 
\begin{tabular}{l}
{\small Image is upright}%
\end{tabular}
& 
\begin{tabular}{l}
{\small Image is inverted}%
\end{tabular}
\\ \hline
{\small Radius }$\left( R_{1}\text{ and }R_{2}\right) $ & 
\begin{tabular}{l}
{\small Center of curvature is in} \\ 
{\small front of surface}%
\end{tabular}%
{\small \ } & 
\begin{tabular}{l}
{\small Center of curvature is in} \\ 
{\small back of surface}%
\end{tabular}%
{\small \ } \\ \hline
Focal length $\left( f\right) $ & In {\small front of surface} & In {\small %
back of surface} \\ \hline
\end{tabular}%
\]

\subsection{Aberrations}

It's time to ask what happens when we don't use paraxial rays (when the
paraxial ray approximation is not valid) but our equations are based on
using only small angles or paraxial rays. You might guess that we will have
problems in our imagery. That sounds bad, so let's use a special word for
\textquotedblleft problems\textquotedblright\ in imagery. That word is \emph{%
aberration}.

Our first aberration comes from making the mirror the wrong shape. It should
be a parabola, but instead it is a section of a sphere. We will call the
problem this creates, \emph{spherical aberration.}

Remember that reason the we use spherical shapes is that spherical shapes
are easier to make than parabolas or hyperbole, or other shapes. So optics
manufacturers have been using spherical optics for centuries. Later, when we
do lenses we will find there are ways to correct for spherical aberration.

But what does spherical aberration do? We end up with different images of
our object at different image locations. The center rays (blue dashed lines
in the next figure) make an image at $I_{1}$ but the outside (non-paraxial)
rays focus closer at $I_{2.}$ If we put up a screen to see the image, we
will have a blurry image because some of the light is out of focus at $I_{1}$
but a different bit of light is out of fucus at $I_{2}.$ There is no one
place where all the light is in focus. \FRAME{dtbpF}{3.5976in}{2.1975in}{0pt%
}{}{}{Figure}{\special{language "Scientific Word";type
"GRAPHIC";maintain-aspect-ratio TRUE;display "USEDEF";valid_file "T";width
3.5976in;height 2.1975in;depth 0pt;original-width 3.55in;original-height
2.1577in;cropleft "0";croptop "1";cropright "1";cropbottom "0";tempfilename
'TN9QPD63.wmf';tempfile-properties "XPR";}}Most of the time for this class,
we will design our optical system so the object is near the principal axis,
so we can make the paraxial approximation that fixes this problem. But in
real optical systems we can't do this.

This optical problem was the problem with the Hubble Space Telescope!

\FRAME{dtbpFU}{3.5169in}{1.7722in}{0pt}{\Qcb{{\protect\small This before
(left) WFPC and after (right) WFPC2 image of the core of the galaxy M100
shows the dramatic improvement in Hubble's view of the universe after the
first servicing mission in December 1993. (Images Courtesy NASA)}}}{}{Figure%
}{\special{language "Scientific Word";type "GRAPHIC";maintain-aspect-ratio
TRUE;display "USEDEF";valid_file "T";width 3.5169in;height 1.7722in;depth
0pt;original-width 2.0453in;original-height 1.017in;cropleft "0";croptop
"1";cropright "1";cropbottom "0";tempfilename
'HubbleTrouble.wmf';tempfile-properties "XNPR";}}

And it took some work to correct this problem. We will learn more about this
as we study lenses next.

\subsection{Semi-Infinite Lenses and the Semi-Infinite Lens Image Equation}

We learned how to find an image location graphically for mirrors. But many
optical systems including contacts and glasses, use lenses. Let's look at
lenses next.

Let's start by thinking of a special case for refraction. A curved surface
on a very large piece of glass. We will assume that the piece of glass is
semi-infinate (goes on forever in one direction), but all it has to be is
very large.

I jokingly call this a semi-infinate bump of glass. But I\ have also heard
this called a single sided lens. It may seem like a far fetched situation,
but it turns out that our eyes are essentially a system like this because
our eyes are filled with a jelly like substance, and the light in our eyes
stays in this substance until it is detected in our retina. So let's look at
this situation. \FRAME{dhF}{3.2491in}{1.1865in}{0pt}{}{}{Figure}{\special%
{language "Scientific Word";type "GRAPHIC";maintain-aspect-ratio
TRUE;display "USEDEF";valid_file "T";width 3.2491in;height 1.1865in;depth
0pt;original-width 5.2555in;original-height 1.9009in;cropleft "0";croptop
"1";cropright "1";cropbottom "0";tempfilename
'TN9QPD64.wmf';tempfile-properties "XPR";}}

Take a point object that either glows, or has rays of light reflecting from
it. The rays leave the object and reach the surface of the glass. The rays
will refract at the surface. Each bends toward the normal, but because of
the curvature of the glass, every ray has a different normal. \FRAME{dtbpF}{%
3.351in}{1.2489in}{0in}{}{}{Figure}{\special{language "Scientific Word";type
"GRAPHIC";maintain-aspect-ratio TRUE;display "USEDEF";valid_file "T";width
3.351in;height 1.2489in;depth 0in;original-width 3.3856in;original-height
1.2444in;cropleft "0";croptop "1";cropright "1";cropbottom "0";tempfilename
'TN9QPD65.wmf';tempfile-properties "XPR";}}The result is tat the rays all
converge toward the center. We can identify this convergence point as the
image of the point object. For mirrors we used the law of reflection to know
where the light would go. But for lenses, the light is transmitted. So at
the surface we can find the refracted angles using Snell's law%
\[
n_{1}\sin \theta _{1}=n_{2}\sin \theta _{2}
\]

We will again use the small angle approximation. Then $\theta _{1}$ and $%
\theta _{2}$ are small and none of the rays are very far away from the axis.
We know this as the paraxial approximation. Snell's law becomes 
\[
n_{1}\theta _{1}=n_{2}\theta _{2}
\]

\FRAME{dtbpF}{5.041in}{1.9441in}{0pt}{}{}{Figure}{\special{language
"Scientific Word";type "GRAPHIC";maintain-aspect-ratio TRUE;display
"USEDEF";valid_file "T";width 5.041in;height 1.9441in;depth
0pt;original-width 4.4469in;original-height 1.6985in;cropleft "0";croptop
"1";cropright "1";cropbottom "0";tempfilename
'TN9QPD66.wmf';tempfile-properties "XPR";}}Using the more detailed figure
above with it's exaggerated angles (not so paraxial but easier to see), we
observe triangles $SAC$ and $PAC.$ We can see that for triangle $SAC$ the
top angle labeled $\eta $, plus $\theta _{1}$ must be $180.$ 
\[
\eta +\theta _{1}=180\unit{%
%TCIMACRO{\U{b0}}%
%BeginExpansion
{{}^\circ}%
%EndExpansion
}
\]%
We also know that the sum of interior angles of a triangle must equal $180.$
So%
\[
180\unit{%
%TCIMACRO{\U{b0}}%
%BeginExpansion
{{}^\circ}%
%EndExpansion
}=\alpha +\beta +\eta
\]%
then 
\begin{eqnarray*}
\eta +\theta _{1} &=&\alpha +\beta +\eta \\
\theta _{1} &=&\alpha +\beta
\end{eqnarray*}%
Likewise, from triangle $PAC$,%
\[
\beta +\psi =180\unit{%
%TCIMACRO{\U{b0}}%
%BeginExpansion
{{}^\circ}%
%EndExpansion
}
\]

\[
180\unit{%
%TCIMACRO{\U{b0}}%
%BeginExpansion
{{}^\circ}%
%EndExpansion
}=\psi +\gamma +\theta _{2}
\]%
so then 
\[
\beta =\theta _{2}+\gamma
\]%
then, 
\[
\theta _{2}=\beta -\gamma
\]

and we can write our paraxial Snell's law as%
\begin{eqnarray*}
n_{1}\theta _{1} &=&n_{2}\theta _{2} \\
n_{1}\left( \alpha +\beta \right) &=&n_{2}\left( \beta -\gamma \right) \\
n_{1}\alpha +n_{1}\beta &=&n_{2}\beta -n_{2}\gamma \\
n_{1}\alpha +n_{2}\gamma &=&n_{2}\beta -n_{1}\beta \\
n_{1}\alpha +n_{2}\gamma &=&\beta \left( n_{2}-n_{1}\right)
\end{eqnarray*}

Looking at the figure. We see that $d$ is a leg of three different right
triangles ($SAV,$ $ACV$, and $PAV$). The ray in the figure is clearly not a
paraxial ray. If we use a paraxial ray, then the point $V$ will approach the
air-glass boundary. When this happens, then $SV=d_{o},$ $VC=R$, and $%
VP=d_{i}.$ So we can write (using the small angle approximation)

\begin{eqnarray*}
\tan \alpha &\approx &\alpha \approx \frac{d}{d_{o}} \\
\tan \beta &\approx &\beta \approx \frac{d}{R} \\
\tan \gamma &\approx &\gamma \approx \frac{d}{d_{i}}
\end{eqnarray*}%
So, our Snell's law becomes%
\begin{eqnarray*}
n_{1}\alpha +n_{2}\gamma &=&\beta \left( n_{2}-n_{1}\right) \\
n_{1}\frac{d}{d_{o}}+n_{2}\frac{d}{d_{i}} &=&\frac{d}{R}\left(
n_{2}-n_{1}\right)
\end{eqnarray*}%
We can divide out the $d$'s%
\begin{equation}
\frac{n_{1}}{d_{o}}+\frac{n_{2}}{d_{i}}=\frac{\left( n_{2}-n_{1}\right) }{R}
\label{Bump_Equation}
\end{equation}

We can use this formula to convince ourselves that no matter what the angle
is (providing it is small), the rays will form an image at $P.$

Real images will be in the glass for this special case (we will fix this
with non-infinite lenses soon) so we need to change sign conventions.

\[
\begin{tabular}{|l|l|l|}
\hline
\textbf{Quantity} & \textbf{Positive if} & \textbf{Negative if} \\ \hline
{\small Object location }$\left( d_{o}\right) $ & 
\begin{tabular}{l}
{\small Object is in front of interface} \\ 
{\small \ surface}%
\end{tabular}
& 
\begin{tabular}{l}
{\small Object is in back of interface} \\ 
{\small \ surface (virtual object)}%
\end{tabular}%
{\small \ } \\ \hline
{\small Image location }$\left( d_{i}\right) $ & 
\begin{tabular}{l}
{\small Image is in back of interface} \\ 
{\small \ surface (real image)}%
\end{tabular}
& 
\begin{tabular}{l}
{\small Image is in front of interface} \\ 
{\small \ surface (virtual image)}%
\end{tabular}
\\ \hline
{\small Image height }$\left( h_{i}\right) $ & 
\begin{tabular}{l}
{\small Image is upright}%
\end{tabular}
& 
\begin{tabular}{l}
{\small Image is inverted}%
\end{tabular}
\\ \hline
{\small Radius }$\left( R\right) $ & 
\begin{tabular}{l}
{\small Center of curvature is in} \\ 
{\small back of interface surface}%
\end{tabular}%
{\small \ } & 
\begin{tabular}{l}
{\small Center of curvature is in} \\ 
{\small front of interface surface}%
\end{tabular}%
{\small \ } \\ \hline
\end{tabular}%
\]

We could go through the entire derivation and switch the indices of
refraction. It turns out we get the same equation. The results will be
different, but the equation is the same.

\subsubsection{Flat Refracting surfaces}

\FRAME{dtbpF}{2.6282in}{2.4647in}{0in}{}{}{Figure}{\special{language
"Scientific Word";type "GRAPHIC";maintain-aspect-ratio TRUE;display
"USEDEF";valid_file "T";width 2.6282in;height 2.4647in;depth
0in;original-width 2.5867in;original-height 2.4232in;cropleft "0";croptop
"1";cropright "1";cropbottom "0";tempfilename
'TN9QPD67.wmf';tempfile-properties "XPR";}}

Let's consider a fish tank. The fish tank has an interface, but it is flat.
Can we use our equation (\ref{Bump_Equation}) to describe this?

The answer is yes, if we let $R=\infty .$ This makes sense for a flat
surface. If we have an infinitely large sphere, then our small part of that
spherical surface that makes up the fish tank wall will be very flat.

Then%
\[
\frac{n_{1}}{d_{o}}+\frac{n_{2}}{d_{i}}=\frac{\left( n_{2}-n_{1}\right) }{%
\infty }
\]%
or%
\[
\frac{n_{1}}{d_{o}}+\frac{n_{2}}{d_{i}}=0
\]%
we see that 
\[
d_{i}=-d_{o}\frac{n_{2}}{n_{1}}
\]

We can see from this that, since $n_{2}<n_{1}$ for this case $d_{i}$ is less
than $d_{o}.$ The image is in the water! It is in \emph{front} of the
water/air boundary of the tank. But light is not actually coming from that
location. What has happened is that the light path has bent, but our brain
image processing center thinks light travels only in straight lines. So it
tells us there is a fish there, but the fish appears closer than it is. The
light doesn't come from where we think the fish is. And as we know, this
means the image is virtual, just like for flat mirrors!

\section{Thin Lenses 3.2.4, 3.2.5}

In our last lecture we found an equation for finding the image location from
a single sided lens (or semi-infinite bump) knowing the object distance and
the lens radius and index of refraction. But most lenses have two sides
(think of eye glass lenses). We need to do a little more work to find an
equation for the location of an image with a double sided lens. We will to
this work in this lecture.

%TCIMACRO{%
%\TeXButton{Fundamental Concepts}{\vspace{0.10in}\hspace{-0.37in}{\Large Fundamental Concepts\vspace{0.2in}}}}%
%BeginExpansion
\vspace{0.10in}\hspace{-0.37in}{\Large Fundamental Concepts\vspace{0.2in}}%
%EndExpansion

\begin{itemize}
\item For a thin lens, we can describe where the light goes using $\frac{1}{%
d_{o}}+\frac{1}{d_{i}}=\frac{1}{f}$ where $\frac{1}{f}=\left( n-1\right)
\left( \frac{1}{R_{1}}-\frac{1}{R_{2}}\right) $

\item There is a sign convention for all of these equations.

\item We can think of a complicated lens as a series of interfaces between
materials.

\item We can find the image location of a complicated set of interface
surfaces by going through each surface and finding the image created by that
surface, then using that image as the object of the next surface.

\item The wavelength dependence of the index of refraction gives rise to
another problem, chromatic aberration.
\end{itemize}

\subsection{Images from lenses}

Let's think once again about what an image is. If you haven't done this
before, take a piece of paper and a lens, and hold up the lens in a darkened
room that has some bright object in it. Move the lens or the paper back and
forth, and at just the right distance, a miniature picture of the bright
object will appear. We should think about what the word \textquotedblleft
picture\textquotedblright\ means in this sense.

\FRAME{dhF}{2.6835in}{2.0237in}{0pt}{}{}{Figure}{\special{language
"Scientific Word";type "GRAPHIC";maintain-aspect-ratio TRUE;display
"USEDEF";valid_file "T";width 2.6835in;height 2.0237in;depth
0pt;original-width 2.6411in;original-height 1.9847in;cropleft "0";croptop
"1";cropright "1";cropbottom "0";tempfilename
'TN9QPD68.wmf';tempfile-properties "XPR";}}We have talked about how we see
objects. Remember the BYU-I guys from a few lectures ago.\FRAME{dhF}{3.173in%
}{1.3534in}{0pt}{}{}{Figure}{\special{language "Scientific Word";type
"GRAPHIC";maintain-aspect-ratio TRUE;display "USEDEF";valid_file "T";width
3.173in;height 1.3534in;depth 0pt;original-width 3.128in;original-height
1.3188in;cropleft "0";croptop "1";cropright "1";cropbottom "0";tempfilename
'TN9QPD69.wmf';tempfile-properties "XPR";}}Our eyes gather rays that are
diverging from the object because light has bounced off of the object. Our
eyes intersect a diverging set of rays that form a definite pattern. That
diverging set of rays forming a pattern is the picture of the object.

So when we say that the lens has formed a miniature picture of our object,
we mean that the lens has somehow formed a diverging set of rays that form a
pattern that looks like the pattern formed by the diverging set of rays
coming from the object, itself. In other words, the object forms a diverging
set of rays. And our lens forms a duplicate set of rays in the same pattern,
so we see the same thing. The lens' version may be smaller, upside down, and
backwards, but it is still essentially the same pattern. We saw this with
our semi-infinite lens, but it must be true for less infinite lenses as well.

As a first step to see how this works, consider our fish tank again. It
would be bad on the fish, but think about looking at a fish in air. The room
light would bounce off of the fish, and we would have a diverging set of
rays from every point on the fish (see next figure). \FRAME{dtbpF}{4.1468in}{%
2.0643in}{0pt}{}{}{Figure}{\special{language "Scientific Word";type
"GRAPHIC";maintain-aspect-ratio TRUE;display "USEDEF";valid_file "T";width
4.1468in;height 2.0643in;depth 0pt;original-width 4.0975in;original-height
2.0254in;cropleft "0";croptop "1";cropright "1";cropbottom "0";tempfilename
'TN9QPD6A.wmf';tempfile-properties "XPR";}}We can see that the picture is
made from every point on the fish being \textquotedblleft
imaged\textquotedblright\ to a point on our retina. We collect the rays
leaving every point on the fish, and bring them to corresponding points on
the retina (our biological light detector array) to make the picture.

It will take us a few lectures to see exactly how this is done with an
optical lens included, but as a first step, let's consider the fish tank,
itself. Put the fish back in the tank and look at it (next figure).

\FRAME{dtbpF}{3.9115in}{2.4811in}{0pt}{}{}{Figure}{\special{language
"Scientific Word";type "GRAPHIC";maintain-aspect-ratio TRUE;display
"USEDEF";valid_file "T";width 3.9115in;height 2.4811in;depth
0pt;original-width 4.3552in;original-height 2.7518in;cropleft "0";croptop
"1";cropright "1";cropbottom "0";tempfilename
'TN9QPD6B.wmf';tempfile-properties "XPR";}}Rays still come from the fish.
But we now know that the change from a slow material to a fast material will
bend the light. These bent rays are collected by our eyes, and the picture
of the fish is formed on the retina just as before. But our eyes interpret
the light as though it went in straight lines with no bends (dotted lines in
the last figure). our visual processing center in our brain is designed to
believe light travels in straight lines, so our mind tells us there is a
fish, but that the fish head (and every other part of the fish) is closer
than it really is. We call this apparent fish at the closer location an
image of the fish, because this is where our brain processing center thinks
the diverging set of rays come from that form the fish pattern.

The next figure shows the details of the rays leaving a dot on the fish head
(with the angles exaggerated so it's easier to see them). \FRAME{dtbpF}{%
2.7363in}{2.565in}{0pt}{}{}{Figure}{\special{language "Scientific Word";type
"GRAPHIC";maintain-aspect-ratio TRUE;display "USEDEF";valid_file "T";width
2.7363in;height 2.565in;depth 0pt;original-width 2.6939in;original-height
2.5235in;cropleft "0";croptop "1";cropright "1";cropbottom "0";tempfilename
'TN9QPD6C.wmf';tempfile-properties "XPR";}}A spot on the fish head is our
object for this set of rays. The distance from the fish-head dot and the
edge of the water/air boundary is our object distance,\emph{\ }$d_{o}.$

The distance from the image of the fish-head dot to the edge of the
water/air boundary is our image distance, $d_{i},$ because it appears to be
where the rays come from. We can find where the image is $\left(
d_{i}\right) $ knowing $d_{o}.$ We can see from the figure that%
\begin{eqnarray*}
\ell &=&d_{o}\tan \theta _{1} \\
&=&d_{i}\tan \theta _{2}
\end{eqnarray*}%
so%
\[
d_{o}\tan \theta _{1}=d_{i}\tan \theta _{2}
\]%
or%
\[
d_{o}\frac{\tan \theta _{1}}{\tan \theta _{2}}=d_{i}
\]%
from Snell's law, we know that 
\[
\frac{\sin \theta _{1}}{\sin \theta _{2}}=\frac{n_{2}}{n_{1}}
\]

Usually we can take the small angle approximation. This would limit our
analysis to rays that are near the central axes. We have done this before
with mirrors. We call this central axis the \emph{optics axis} and the rays
that are not too far away from this axis \emph{paraxial rays.} Then for our
small angles we can write 
\[
\tan \theta _{i}\approx \sin \theta _{i}
\]%
so%
\[
d_{o}\frac{\tan \theta _{1}}{\tan \theta _{2}}\approx d_{o}\frac{\sin \theta
_{1}}{\sin \theta _{2}}=d_{o}\frac{n_{2}}{n_{1}}=d_{i}
\]%
and we have the image distance

\[
d_{i}=d_{o}\frac{n_{2}}{n_{1}}
\]%
Note that once again we have a virtual image. The place where we see the
fish does not have rays of light coming from it! But our brain processing
center believes the rays of light didn't bend, so it extends the rays we see
backward into the fish tank and it interprets what we see as a fish being
where those extended rays meet.

This is not so useful unless you have some burning need to know where fish
are in a tank. But we now have the vocabulary to discuss the larger problem
of how a lens works, which we will take up next time.

Before we do a lot of math to describe how lenses work, lets think about our
early childhood experiences. You may have burned things with a magnifying
glass\footnote{%
If you didn't do this as a child, consider trying it now. Be responsible and
safe, but it is valuable to see how this works.}. Using the idea of a ray
diagram, here is what happens.\FRAME{dhF}{3.8017in}{2.2329in}{0pt}{}{}{Figure%
}{\special{language "Scientific Word";type "GRAPHIC";maintain-aspect-ratio
TRUE;display "USEDEF";valid_file "T";width 3.8017in;height 2.2329in;depth
0pt;original-width 3.7533in;original-height 2.1923in;cropleft "0";croptop
"1";cropright "1";cropbottom "0";tempfilename
'TN9QPD6D.wmf';tempfile-properties "XPR";}}We know that the rays from the
Sun come from so far away that they are essentially parallel. We know that
these rays come together to a fine point that can start a fire. The point
where these rays converge is important to us. We found such a point from our
mirrors. We called this point the \emph{focal point}.

Now think about Snell's law. We have 
\[
n_{1}\sin \theta _{1}=n_{2}\sin \theta _{2}
\]%
The light will bend at a surface, but snell's law doesn't care if $%
n_{1}>n_{2}$ of if $n_{2}>n_{1}.$ It works both ways. The light will follow
the same path either direction. We would expect that if we put a light
source at the focal distance, the rays should come out parallel. \FRAME{dhF}{%
4.1926in}{2.4561in}{0pt}{}{}{Figure}{\special{language "Scientific
Word";type "GRAPHIC";maintain-aspect-ratio TRUE;display "USEDEF";valid_file
"T";width 4.1926in;height 2.4561in;depth 0pt;original-width
4.1433in;original-height 2.4146in;cropleft "0";croptop "1";cropright
"1";cropbottom "0";tempfilename 'TN9QPD6E.wmf';tempfile-properties "XPR";}}%
This is one way manufacturers make LED flashlights.

We need one other bit of information that we already have seen from basic
refraction.\FRAME{dhF}{1.996in}{1.6933in}{0pt}{}{}{Figure}{\special{language
"Scientific Word";type "GRAPHIC";maintain-aspect-ratio TRUE;display
"USEDEF";valid_file "T";width 1.996in;height 1.6933in;depth
0pt;original-width 3.6841in;original-height 3.122in;cropleft "0";croptop
"1";cropright "1";cropbottom "0";tempfilename
'TN9QPD6F.wmf';tempfile-properties "XPR";}}A flat block does refract the
light, but when the light leaves the block it is only displaced, it retains
the original direction. With these three ray scenarios, we can describe how
a lens works.

We know that for every point on the object, we get millions of reflected
rays that diverge. In the next figure, rays of light are reflecting off of
the top of the large, black arrow. The lens must collect these rays together
to form the corresponding point on the image.

\FRAME{dhF}{4.1373in}{2.4275in}{0pt}{}{}{Figure}{\special{language
"Scientific Word";type "GRAPHIC";maintain-aspect-ratio TRUE;display
"USEDEF";valid_file "T";width 4.1373in;height 2.4275in;depth
0pt;original-width 4.088in;original-height 2.3869in;cropleft "0";croptop
"1";cropright "1";cropbottom "0";tempfilename
'TN9QPD6G.wmf';tempfile-properties "XPR";}}Of course, this happens for every
point on the image. \FRAME{dhF}{4.1926in}{2.4561in}{0pt}{}{}{Figure}{\special%
{language "Scientific Word";type "GRAPHIC";maintain-aspect-ratio
TRUE;display "USEDEF";valid_file "T";width 4.1926in;height 2.4561in;depth
0pt;original-width 4.1433in;original-height 2.4146in;cropleft "0";croptop
"1";cropright "1";cropbottom "0";tempfilename
'TN9QPD6H.wmf';tempfile-properties "XPR";}}But we usually pick the top of
the object. If we place the bottom of the object on the \emph{optic axis}%
\footnote{%
The line drawn on the figure thorugh the middle of the lens.}, the bottom of
the image will also be on the optic axis. So knowing where the bottom of the
image is, and finding the top of the image gives a pretty good idea of where
the rest of the image must be. So we will draw diagrams for the top of the
object to find the top of the image.

But suppose this is not true? For example, when we use a camera, we do not
align the object with it's bottom on the axis before we shoot. We center the
object on our viewfinder. So how do we find the location for the bottom of
the image?\FRAME{dtbpF}{4.1329in}{2.4967in}{0pt}{}{}{Figure}{\special%
{language "Scientific Word";type "GRAPHIC";maintain-aspect-ratio
TRUE;display "USEDEF";valid_file "T";width 4.1329in;height 2.4967in;depth
0pt;original-width 4.9199in;original-height 2.962in;cropleft "0";croptop
"1";cropright "1";cropbottom "0";tempfilename
'TN9QPD6I.wmf';tempfile-properties "XPR";}}We can, of course, trace two more
rays for the bottom of the image and find the location of the bottom of the
image as well.

We know that light is a superposition of millions of waves and that when
light reflects off of something like your neighbor or a fish or a soccer
player there will be millions of rays. Drawing millions of rays is
impractical, and fortunately, not needed. From our study of mirrors, we can
guess that we can find three simple rays that leave the top of the object
and see where these rays converge to form the top of the image. Let's start
with a ray that travels from the top of the object and travels parallel to
the optic axis. We recognize this ray as being like one of the rays from the
Sun. It comes in parallel, so it will leave the lens and travel through the
focal point.\FRAME{dhF}{4.3889in}{2.5815in}{0pt}{}{}{Figure}{\special%
{language "Scientific Word";type "GRAPHIC";maintain-aspect-ratio
TRUE;display "USEDEF";valid_file "T";width 4.3889in;height 2.5815in;depth
0pt;original-width 4.3379in;original-height 2.5399in;cropleft "0";croptop
"1";cropright "1";cropbottom "0";tempfilename
'TN9QPD6J.wmf';tempfile-properties "XPR";}}For a second easy ray, lets take
the case that is like our flat block. Near the center of the lens, the sides
are nearly flat. So we expect that the ray will leave in about the same
direction as it was going before it struck the lens.\FRAME{dhF}{3.1522in}{%
1.8542in}{0pt}{}{}{Figure}{\special{language "Scientific Word";type
"GRAPHIC";maintain-aspect-ratio TRUE;display "USEDEF";valid_file "T";width
3.1522in;height 1.8542in;depth 0pt;original-width 8.9949in;original-height
5.2736in;cropleft "0";croptop "1";cropright "1";cropbottom "0";tempfilename
'TN9QPD6K.wmf';tempfile-properties "XPR";}}Technically this ray should be
shifted. This is where our thin lens approximation comes in. If the lens is
thin, then the ray going through the center of the lens won't be shifted
much, and we can ignore the shift. We just draw ray 2 as a straight line.

Two rays are really enough to determine where the top of the image will be,
but there is a third ray that is easy to draw, so let's draw it to give us
more confidence in our answer. That ray is one that leaves the top of the
object and passes through the focal point on the object side of the lens.
This situation we also recognize. Since the ray goes through the focal
point, it is as though the light came from that point. This is like our LED
flashlight case. This ray will leave the lens parallel to the optic axis.

\FRAME{dhF}{4.2627in}{2.4976in}{0pt}{}{}{Figure}{\special{language
"Scientific Word";type "GRAPHIC";maintain-aspect-ratio TRUE;display
"USEDEF";valid_file "T";width 4.2627in;height 2.4976in;depth
0pt;original-width 4.2125in;original-height 2.4569in;cropleft "0";croptop
"1";cropright "1";cropbottom "0";tempfilename
'TN9QPD6L.wmf';tempfile-properties "XPR";}}Where all three rays intersect,
we will have the top of the image. \FRAME{dhF}{4.9753in}{2.1067in}{0pt}{}{}{%
Figure}{\special{language "Scientific Word";type
"GRAPHIC";maintain-aspect-ratio TRUE;display "USEDEF";valid_file "T";width
4.9753in;height 2.1067in;depth 0pt;original-width 4.9216in;original-height
2.0678in;cropleft "0";croptop "1";cropright "1";cropbottom "0";tempfilename
'TN9QPD6M.wmf';tempfile-properties "XPR";}}And we can sketch in the rest of
the image. Notice that in this case, the image is upside down. That is
normal. Also notice that it is smaller than the object. We say that the
image is magnified, which may still seem a little bit strange. But in
optics, a magnification of greater than one means that the image is bigger
than the object. This is like a movie projector that makes a large image of
a small film segment. The magnification can be equal to one, meaning the
object and image are the same size. And finally the magnification can be
less than one. This means that the image is smaller than the object. This is
a convenient definition, because then we can use the same equation to
describe all three situations.%
\[
m\equiv \frac{\text{Image height}}{\text{Object height}}=\frac{h_{i}}{h_{o}}
\]%
where $h_{o}$ is the object height, and $h_{i}$ is the image height. This is
the same equation for magnification that we found for mirrors! It's still
true that by convention (meaning physicists got together and voted on this)
we say that an upside down image has a negative magnification. You just have
to memorize this, there is no obvious reason for this except it is
mathematically convenient.

\subsection{Virtual Images with Lenses}

Lets take another case and draw a ray diagram. This time let's place the
object closer than the focal distance. This is the case when we use a lens
as a magnifying glass. The rays will look like this.\FRAME{dhF}{3.6201in}{%
2.3713in}{0pt}{}{}{Figure}{\special{language "Scientific Word";type
"GRAPHIC";maintain-aspect-ratio TRUE;display "USEDEF";valid_file "T";width
3.6201in;height 2.3713in;depth 0pt;original-width 3.5725in;original-height
2.3315in;cropleft "0";croptop "1";cropright "1";cropbottom "0";tempfilename
'TN9QSZ6S.wmf';tempfile-properties "XPR";}}Notice that these rays never
converge! We won't get an image that could project on a paper. But we know
that there is an image, we can look through the lens and see it! And that is
the key. The image does not really exist. There is no place where light
gathers and then diverges into a pattern that we recognize. But we look
through the lens, and our mind interprets the diverging rays coming from the
lens as though they had only traveled in straight lines. If we extend these
rays backwards along straight lines, they appear to come from a common
point. This is the point they would have had to have come from if there were
no lens in or der to make the same recognizable pattern that we see. We
believe we see an image at this location. But no light really comes from
there! Because this image is not really made from light diverging from this
position, we call it a \emph{virtual image} just like we did for this
situation for mirrors The image we formed before that could be projected on
a screen is called a \emph{real image}. For a real image, the light is
actually at the image location and diverges from that position. This is just
like our image terminology for mirrors!

By convention, we say the distance from the lens to the virtual image is a
negative value.

\subsubsection{Diverging Lenses}

So far our lenses have only been the sort that work as magnifying glasses.
We call these \emph{converging lenses.} These lenses are thicker in the
middle and thinner on the edges. Because of this they are sometimes called 
\emph{convex lenses}. By convention, we say the focal distance for this type
of lens is positive. For this reason, they are often called \emph{positive
lenses.}

But what if we make a lens that is thinner in the middle and thicker on the
edges? We can call this sort of lens a \emph{concave lens}, and we will give
it a negative focal length by convention, so we can also call it a \emph{%
negative lens}. But what would this lens do? If we think about our three
rays, ray $1$ won't be bent toward the optic axis for this type of lens. In
fact, if we observe an object through this lens, ray number 1 will appear to
come from the focal point. Ray number 2 will still go through the middle of
the lens, and if the lens is thin enough, ray 2 will pass through
essentially undeviated.

\FRAME{dhF}{3.2984in}{1.7582in}{0pt}{}{}{Figure}{\special{language
"Scientific Word";type "GRAPHIC";maintain-aspect-ratio TRUE;display
"USEDEF";valid_file "T";width 3.2984in;height 1.7582in;depth
0pt;original-width 3.2534in;original-height 1.7201in;cropleft "0";croptop
"1";cropright "1";cropbottom "0";tempfilename
'TN9QSZ6T.wmf';tempfile-properties "XPR";}}Finally ray three leaves the
object in the direction of the far focal point. It will hit the lens and
leave parallel to the optic axis. From the figure we see that these three
rays will never converge. This is like our convex mirror! And from our
mirror experience we expect the rays (with help from our brain) will form a
virtual image. If we extend the rays backward as shown, we see that the
extensions all meet at a point. The rays leaving the lens appear to come
from this point. This is the location of the top of the virtual image of the
object.

You might wonder what good such a lens could do, but we will find that this
type of lens is used to correct vision for nearsighted people. It's likely
that many people in our class have this type of lens with them either on
their eye (contacts) or on their nose (eye glasses).

\subsection{Thin Lenses}

Lets' find an equation for a lens made from a part of a spherical surface
once more. But this time, let's let it be more practical and not make the
\textquotedblleft lens\textquotedblright\ semi-infinite. We will need to
deal with two sides of the lens because (usually) both will be curved.

We found that for refraction%
\[
\frac{n_{1}}{d_{o}}+\frac{n_{2}}{d_{i}}=\frac{\left( n_{2}-n_{1}\right) }{R}
\]

but we did this for a spherical bump on a semi-infinite piece of glass. For
this new problem let's make a few assumptions:

\begin{enumerate}
\item We have two spherical surfaces, with $R_{1}$ and $R_{2}$ as the radii
of curvature

\item We have only paraxial rays

\item The image formed by one refractive surface serves as the object for
the second surface

\item The lens is not very thick (the thickness is much smaller than both $%
R_{1}$ and $R_{2})$
\end{enumerate}

The answer we will get is quite simple

\begin{equation}
\frac{1}{d_{o}}+\frac{1}{d_{i}}=\frac{1}{f}  \label{Thin_lens}
\end{equation}%
where%
\begin{equation}
\frac{1}{f}=\left( n-1\right) \left( \frac{1}{R_{1}}-\frac{1}{R_{2}}\right)
\label{Lens_maker}
\end{equation}%
but to appreciate what it means, lets find out where it comes from.

\subsection{Derivation of the lens equation}

Consider the optical element in the figure below. Notice that our object is
a dot, so our image will also be a dot. This is not as boring as it sounds
if we consider any object can be considered as a collection of dots.%
\footnote{%
In reality, we are a collection of atoms, which are pretty much dots as far
as this analysis is concerned!} And in our past analysis we used the top of
our object as our image, which was really just the dot at the top of the
object. We could use that to figure out where the rest of the image would
be. So let's start with a dot.\FRAME{dtbpF}{4.5714in}{2.8098in}{0pt}{}{}{%
Figure}{\special{language "Scientific Word";type
"GRAPHIC";maintain-aspect-ratio TRUE;display "USEDEF";valid_file "T";width
4.5714in;height 2.8098in;depth 0pt;original-width 4.5186in;original-height
2.7674in;cropleft "0";croptop "1";cropright "1";cropbottom "0";tempfilename
'TN9QSZ6U.wmf';tempfile-properties "XPR";}}Light enters at a spherical
surface on the left hand side. We use a point object located at $S$ (for 
\emph{source}) on the principal axis. We could put our dot object anywhere,
but let's put it on the axis and trace two rays. The ray along the principal
axis crosses each spherical surface at right angles, and therefore traves
straight through the optic (this makes one ray very easy to trace!). The
second ray hits the first spherical surface at point $A.$ It is refracted
and travels to point $B$. It is again refracted and travels toward the
principal axis, crossing at $P.$ The image location is the intersection of
these rays, so we have an image at $P.$

Lets study the surfaces separately

\textbf{Surface 1:}

\FRAME{dtbpF}{4.3448in}{2.6117in}{0pt}{}{}{Figure}{\special{language
"Scientific Word";type "GRAPHIC";maintain-aspect-ratio TRUE;display
"USEDEF";valid_file "T";width 4.3448in;height 2.6117in;depth
0pt;original-width 4.9147in;original-height 2.9438in;cropleft "0";croptop
"1";cropright "1";cropbottom "0";tempfilename
'TN9QSZ6V.wmf';tempfile-properties "XPR";}}

Let's treat surface $1$ as though surface $2$ did not exist. When it's at
surface $1$ The light doesn't know about surface $2.$ It only interacts with
the surface where it is. So ignoring surface $2$ at this point is reasonable.

The light would bend at point $A$ and head off into the lens material. This
is just our single sided lens (semi-infinite bump) problem. And if we
imagine our head to somehow be in the glass looking at the refracted ray,
our brain would think that the ray traveled in a straight line. So our brain
would think the light came from location $P_{1}.$ That would be a virtual
image of our $S$ point, so the distance from the lens to $P_{1}=d_{i1}.$ We
can find this image distance using our single sided lens (semi-infinite
bump) equation.%
\[
\frac{n_{1}}{d_{o1}}+\frac{n_{2}}{d_{i1}}=\frac{\left( n_{2}-n_{1}\right) }{R%
}
\]

We can consider $n_{1}=1$ and $n_{2}=n_{glass}$ for a air-glass interface
and noting that $d_{i1}$ is negative (by convention). Then%
\begin{equation}
\frac{1}{d_{o1}}-\frac{n_{glass}}{d_{i1}}=\frac{\left( n_{glass}-1\right) }{%
R_{1}}  \label{Surface1}
\end{equation}

Note that our rays are \emph{not} converging in the glass this time. We can
find the image formed by this side of our lens by tracing the diverging rays
backward as we did for the fish tank. The image formed from the first side
of the lens is very much virtual.

\textbf{Surface 2}

Now consider the second surface.

\FRAME{dtbpF}{4.4348in}{2.2286in}{0pt}{}{}{Figure}{\special{language
"Scientific Word";type "GRAPHIC";maintain-aspect-ratio TRUE;display
"USEDEF";valid_file "T";width 4.4348in;height 2.2286in;depth
0pt;original-width 4.3837in;original-height 2.188in;cropleft "0";croptop
"1";cropright "1";cropbottom "0";tempfilename
'TN9QSZ6W.wmf';tempfile-properties "XPR";}}The second surface isn't alive
and doesn't know about how light traves. All the second surface sees is
light diverging as though it came from a semi-infinite piece of glass with
the object at $P_{1}.$ It's not true, the light came from $S.$ But the
second surface can't tell the difference. It only knows light is incident in
a particular pattern, and that pattern appears to come from point $P_{1}.$
So we can treat the situation at surface $2$ as though the object is the
virtual image formed by surface $1.$ This is important, so let's say it
again. The object for surface $2$ is the image from surface $1.$ We can find
the object distance for surface $2$ 
\[
d_{o2}=d_{i1}+t
\]

We again use our refractive equation%
\[
\frac{n_{1}}{d_{o2}}+\frac{n_{2}}{d_{i2}}=\frac{\left( n_{2}-n_{1}\right) }{R%
}
\]%
but this time we identify $n_{1}=n_{glass}$ and $n_{2}=1$. We have for
surface $2$%
\begin{equation}
\frac{n_{glass}}{d_{o2}}+\frac{1}{d_{i2}}=\frac{\left( 1-n_{glass}\right) }{%
R_{2}}
\end{equation}%
or%
\begin{equation}
\frac{n_{glass}}{d_{i1}+t}+\frac{1}{d_{i2}}=\frac{\left( 1-n_{glass}\right) 
}{R_{2}}  \label{Surface2}
\end{equation}

Now we take our thin lens approximation. Let $t\rightarrow 0.$ Then
equations (\ref{Surface1}) and (\ref{Surface2}) become%
\[
\frac{1}{d_{o1}}-\frac{n_{glass}}{d_{i1}}=\frac{\left( n_{glass}-1\right) }{%
R_{1}}
\]

\[
\frac{n_{glass}}{d_{i1}}+\frac{1}{d_{i2}}=\frac{\left( 1-n_{glass}\right) }{%
R_{2}}
\]

We would prefer a nice equation that has how far away the object is from the
whole lens and how far away the image is formed from the whole lens like our
thin mirror equation. We can get such an equation by adding these two
equations for the two surfaces%
\[
\frac{1}{d_{o1}}-\frac{n_{glass}}{d_{i1}}+\frac{n_{glass}}{d_{i1}}+\frac{1}{%
d_{i2}}=\frac{\left( n_{glass}-1\right) }{R_{1}}+\frac{\left(
1-n_{glass}\right) }{R_{2}}
\]%
or%
\[
\frac{1}{d_{o1}}+\frac{1}{d_{i2}}=\left( n_{glass}-1\right) \left( \frac{1}{%
R_{1}}-\frac{1}{R_{2}}\right)
\]

This equation is very useful. It reminds us of the mirror equation (well, a
little) . The left hand sides are almost just the same. If we again let $%
d_{o1}=\infty $ (put the object at $\infty $ so the rays enter surface $1$
parallel) we find%
\[
\frac{1}{d_{i2}}=\left( n-1\right) \left( \frac{1}{R_{1}}-\frac{1}{R_{2}}%
\right)
\]%
The spot where the rays gather if the object is infinitely far away is the
focal point, $f.$ So for parallel rays we can identify $d_{i2}=f$ as the
focal length of the optic. That means that 
\[
\frac{1}{f}=\left( n-1\right) \left( \frac{1}{R_{1}}-\frac{1}{R_{2}}\right)
\]%
This equation is known as the \emph{lens makers' equation}. It tells you the
radii of curvature to machine to in or der to make a lens out of a material
with index $n$ in order to have a focal length of $f.$

Then we have a relationship between the object distance in front of the
lens, and the final image in back of the lens: 
\begin{eqnarray*}
\frac{1}{d_{o1}}+\frac{1}{d_{i2}} &=&\left( n-1\right) \left( \frac{1}{R_{1}}%
-\frac{1}{R_{2}}\right) \\
&=&\frac{1}{f}
\end{eqnarray*}%
We can drop the subscripts (which we can do now that we let $t=0$ since the
internal distances for the inside points are not important).%
\[
\frac{1}{d_{o}}+\frac{1}{d_{i}}=\frac{1}{f}
\]%
This is called the \emph{thin lens equation}. The resulting approximate
geometry is shown below.

\FRAME{dtbpF}{5.0548in}{1.5731in}{0pt}{}{}{Figure}{\special{language
"Scientific Word";type "GRAPHIC";maintain-aspect-ratio TRUE;display
"USEDEF";valid_file "T";width 5.0548in;height 1.5731in;depth
0pt;original-width 5.0004in;original-height 1.5376in;cropleft "0";croptop
"1";cropright "1";cropbottom "0";tempfilename
'TN9QSZ6X.wmf';tempfile-properties "XPR";}}And this is fantastic! It is the
same equation as the thin mirror equation! Only we have a different equation
for finding the focal length.

We did our math for just one point of light at $S.$ Of course any real
object is made of lots of points, and not all of the points are on an axis.
But each point will be imaged to a corresponding point on the image. Here is
an example with three points in the object (the $s_{i}$ points) and where
their images are (the $p_{i}$ points).\FRAME{dtbpF}{3.1782in}{2.5685in}{0pt}{%
}{}{Figure}{\special{language "Scientific Word";type
"GRAPHIC";maintain-aspect-ratio TRUE;display "USEDEF";valid_file "T";width
3.1782in;height 2.5685in;depth 0pt;original-width 3.1341in;original-height
2.527in;cropleft "0";croptop "1";cropright "1";cropbottom "0";tempfilename
'TN9QSZ6Y.wmf';tempfile-properties "XPR";}}but it would work for millions of
points making up a whole object. Our simple analysis explains the formation
of actual images and not just point images.

\subsection{Lens Aberrations}

We found that for mirrors if we used a spherical shape rather than a
parabola we had a problem with our images. We called the problem spiracle
aberration. The reason we used a sphere is that spheres are easy to
manufacture. The same problem happens with lenses if we make their sides
from parts of spheres. The correct shape for a lens is a hyperbola, and
hyperbole are hard to manufacture. Spheres are so much easier.\FRAME{dtbpFU}{%
3.3304in}{1.8991in}{0pt}{\Qcb{{\protect\small Example of Spherical
Aberration for Lenses. Optical Ray Trace done in the OSLO\ Optical Design
Package. Rays drawn as green lines are paraxial and rays drawn as red lines
are not.}}}{}{Figure}{\special{language "Scientific Word";type
"GRAPHIC";maintain-aspect-ratio TRUE;display "USEDEF";valid_file "T";width
3.3304in;height 1.8991in;depth 0pt;original-width 3.2846in;original-height
1.8619in;cropleft "0";croptop "1";cropright "1";cropbottom "0";tempfilename
'TN9QSZ6Z.wmf';tempfile-properties "XPR";}}The last figure is a ray trace
done by a professional lens design program. The lens is clearly not thin.
The rays drawn as green lines are the paraxial rays and you can see that
they form a tight focus to the right of the lens. But the rays drawn as red
lines are not paraxial. The non-paraxial and paraxial rays focus at
different spots because we have the wrong shape for our lens. We remember
that spherical aberration was made famous as the main problem with the
Hubble Telescope.

We get another image problem or aberration with lenses. Remember that the
index of refraction is different for different wavelengths of light. If we
send white light into our lens, different colors will bend slightly
differently (think Snell's law and dispersion). So we will get a different
focus point for different colors of light. In the next figure, blue light
(drawn with blue lines) makes an image at a shorter focal distance than red
light (drawn with red lines). This problem is called \emph{chromatic
aberration}. \FRAME{dtbpF}{3.2353in}{2.1862in}{0pt}{}{}{Figure}{\special%
{language "Scientific Word";type "GRAPHIC";maintain-aspect-ratio
TRUE;display "USEDEF";valid_file "T";width 3.2353in;height 2.1862in;depth
0pt;original-width 3.1903in;original-height 2.1465in;cropleft "0";croptop
"1";cropright "1";cropbottom "0";tempfilename
'TN9QSZ70.wmf';tempfile-properties "XPR";}}This will also cause blurry
pictures because different parts of the light will focus at different
distances from the lens. Generally mirrors can't have chromatic aberration
because they don't refract the light. But Lenses always will.

There are many aberrations that come from making lenses that are easy to
manufacture, but that are not perfect in some way. We won't study all these
in this class. If you are curious, we cover more of these aberrations in our
junior level optics class (PH375).

\subsection{Sign Convention}

As we have been working we used our sign convention but we need to add to
our sign convention table a second radius.

\[
\begin{tabular}{|l|l|l|}
\hline
\textbf{Quantity} & \textbf{Positive if} & \textbf{Negative if} \\ \hline
{\small Object location }$\left( d_{o}\right) $ & 
\begin{tabular}{l}
{\small Object is in front of surface} \\ 
\end{tabular}
& 
\begin{tabular}{l}
{\small Object is in back of surface} \\ 
{\small (virtual object)}%
\end{tabular}%
{\small \ } \\ \hline
{\small Image location }$\left( d_{i}\right) $ & 
\begin{tabular}{l}
{\small Image is in back of surface} \\ 
{\small (real image)}%
\end{tabular}
& 
\begin{tabular}{l}
{\small Image is in front of surface } \\ 
{\small (virtual image)}%
\end{tabular}
\\ \hline
{\small Image height }$\left( h^{\prime }\right) $ & 
\begin{tabular}{l}
{\small Image is upright}%
\end{tabular}
& 
\begin{tabular}{l}
{\small Image is inverted}%
\end{tabular}
\\ \hline
{\small Radius }$\left( R_{1}\text{ and }R_{2}\right) $ & 
\begin{tabular}{l}
{\small Center of curvature is in} \\ 
{\small back of surface}%
\end{tabular}%
{\small \ } & 
\begin{tabular}{l}
{\small Center of curvature is in} \\ 
{\small front of surface}%
\end{tabular}%
{\small \ } \\ \hline
Focal length $\left( f\right) $ & Converging lens & Diverging lens \\ \hline
\end{tabular}%
\]

\subsection{Magnification for thin lenses}

We defined the magnification earlier as a comparison of the image height to
the object height. 
\begin{equation}
m=\frac{h_{i}}{h_{o}}
\end{equation}%
And we found we can put this into terms of $d_{o}$ and $d_{i}$ for mirrors.
But will this work for lenses? \FRAME{dtbpF}{3.9773in}{1.9951in}{0pt}{}{}{%
Figure}{\special{language "Scientific Word";type
"GRAPHIC";maintain-aspect-ratio TRUE;display "USEDEF";valid_file "T";width
3.9773in;height 1.9951in;depth 0pt;original-width 5.0764in;original-height
2.5322in;cropleft "0";croptop "1";cropright "1";cropbottom "0";tempfilename
'TN9QSZ71.wmf';tempfile-properties "XPR";}}From the diagram, we can see that 
\begin{eqnarray*}
\tan \theta &=&\frac{h_{o}}{d_{o}} \\
\tan \theta &=&\frac{h_{i}}{d_{i}}
\end{eqnarray*}%
then 
\[
\frac{h_{o}}{d_{o}}=\frac{h_{i}}{d_{i}}
\]%
or 
\[
\frac{h_{i}}{h_{o}}=\frac{d_{i}}{d_{o}}
\]%
And once again we can write the magnification for a converging lens as minus
the ratio of the image distance over the object distance. 
\[
m=-\frac{d_{i}}{d_{o}}
\]%
Just like before the minus sign comes from our sign convention because the
image will be upside down.

We found the thin lens formula using converging lenses, but it works for
diverging lenses as well, so long as the thin lens approximation is valid.

\section{Eyes and Cameras 3.2.6, 3.2.7}

in 1900 George Eastman introduced the Brownie Camera. This event has changed
society dramatically. The camera is very like an artificial eye. We will
study both eyes and cameras in this lecture.

%TCIMACRO{%
%\TeXButton{Fundamental Concepts}{\vspace{0.10in}\hspace{-0.37in}{\Large Fundamental Concepts\vspace{0.2in}}}}%
%BeginExpansion
\vspace{0.10in}\hspace{-0.37in}{\Large Fundamental Concepts\vspace{0.2in}}%
%EndExpansion

\begin{itemize}
\item The eye is like a semi-nfinate double lens with the cornea and the
crystalline lens as the two lenses

\item The crystalline lens does focusing but in biology it is called
accommodation

\item We have three color sensors for red green and blue light

\item Myopia and hyperopia are problems with the shape of the eye and are
corrected with additional lenses

\item Optical lenses are measured in optical power which is $1/f\left(
m\right) $ with the unit if a diopter.

\item Medical people measure the power of a lens in diopters. A diopter is
one over the focal length with the focal length measured in meters.

\item Cameras and other imaging systems use a strange term called and
f-number to tell what the intensity of the image will be%
\[
I\propto \frac{1}{\left( f/\#\right) ^{2}}
\]%
where $f/\#$ is the symbol for f-number and the f-number is given by%
\[
f/\#\equiv \frac{f}{D}
\]
\end{itemize}

\subsection{The Eye}

We in our last lecture we learned how images are formed by lenses. Our eyes
use what we learned to form images that are then converted to nerve signals
that go to our brain for processing. Here is the basic design.

\FRAME{dhF}{3.48in}{1.3673in}{0pt}{}{}{Figure}{\special{language "Scientific
Word";type "GRAPHIC";maintain-aspect-ratio TRUE;display "USEDEF";valid_file
"T";width 3.48in;height 1.3673in;depth 0pt;original-width
3.4342in;original-height 1.3327in;cropleft "0";croptop "1";cropright
"1";cropbottom "0";tempfilename 'TN9QSZ6N.wmf';tempfile-properties "XPR";}}

The figure below shows the parts of the eye. The parts that concern us are
the cornea, crystalline lens, pupil, and the retina. \FRAME{dtbpF}{2.6913in}{%
1.8213in}{0pt}{}{}{Figure}{\special{language "Scientific Word";type
"GRAPHIC";maintain-aspect-ratio TRUE;display "USEDEF";valid_file "T";width
2.6913in;height 1.8213in;depth 0pt;original-width 2.6481in;original-height
1.7841in;cropleft "0";croptop "1";cropright "1";cropbottom "0";tempfilename
'TN9QSZ6O.wmf';tempfile-properties "XPR";}}The Cornea-lens system refracts
the light onto the retina, which detects the light. The lens is focused by a
set of mussels that flatten the lens to change it's focal length. The
focusing process is different from a standard camera (which we will discuss
next). The camera moves the lens to achieve a different image distance. Our
eye can't change the distance between the lens system and the retina. So our
eye changes the shape of the lens, changing it's focal length. \FRAME{dhFU}{%
3.2387in}{1.446in}{0pt}{\Qcb{The crystaline lens becomes thicker, and
therefore more curved when the cilliary musscle flexes. Austin Flint,
\textquotedblleft The Eye as an Optical Instrument,\textquotedblright\ \emph{%
Popular Science Monthly,} Vol. 45, p203, 1894 (Image in the public domain)}}{%
}{Figure}{\special{language "Scientific Word";type
"GRAPHIC";maintain-aspect-ratio TRUE;display "USEDEF";valid_file "T";width
3.2387in;height 1.446in;depth 0pt;original-width 4.1987in;original-height
1.8602in;cropleft "0";croptop "1";cropright "1";cropbottom "0";tempfilename
'TN9QSZ6P.wmf';tempfile-properties "XPR";}}

The focusing process is called \emph{accommodation}. This system becomes
less effective at about the time you reach an age of 40 years because the
lens becomes less flexible. The closest point that can be focused by
accommodation is called the near point. It is about $25\unit{cm}$ on
average. There is, of course, no such thing as an average person, all of us
are a little bit different. Young students probably have a much shorter near
point than $25\unit{cm}.$ For those of us that are a little older, $25\unit{%
cm}$ or more is more likely.

The farthest point that can be focused should be a long way away. It is
called the far point. Both the near and far points degrade with years
leading to bifocal glasses (and much irritation because you can't see as
well, but I'm really not complaining because at least I can see).

The iris changes the area of the pupil (the hole or \emph{aperture} of the
eye). The pupil is, on average, about $7\unit{mm}$ in diameter.

\subsubsection{Nearsightedness}

In some people the cornea-lens system focuses in front of the retina.
Usually this is because the shape of their eyes is not the normal spherical
shape. In this case it is elongated along the optic axis of the eye. This is
called nearsightedness or myopia.

\FRAME{dhF}{3.5475in}{2.4189in}{0pt}{}{}{Figure}{\special{language
"Scientific Word";type "GRAPHIC";maintain-aspect-ratio TRUE;display
"USEDEF";valid_file "T";width 3.5475in;height 2.4189in;depth
0pt;original-width 4.7547in;original-height 3.2335in;cropleft "0";croptop
"1";cropright "1";cropbottom "0";tempfilename
'TN9QSZ72.wmf';tempfile-properties "XPR";}}So their perfectly good
cornea-lens system makes a great image in the vitreous humor (the jellylike
stuff that fills the eye) and not on the retina where it would be detected.
From your experience with lenses you know this would produce a blurry image.
And that is what nearsighted people see much of the time. We can help
nearsighted people by effectively changing the cornea-lens system of the eye
by adding another lens. We want a diverging lens that makes the light more
spread apart so that the lens system of the eye can make it focus where the
retina actually is. Alternately we could flatten the cornea, itself, with
laser ablation.

\subsubsection{Farsightedness}

Sometimes the cornea-lens system focuses in back of the retina. This is
usually because the eye grew too flat or \textquotedblleft
oblate.\textquotedblright\ This is called farsightedness or hyperopia. Once
again we can fix the problem by adding an additional lens. This time a
converging lens.

\FRAME{dhF}{3.7593in}{2.693in}{0pt}{}{}{Figure}{\special{language
"Scientific Word";type "GRAPHIC";maintain-aspect-ratio TRUE;display
"USEDEF";valid_file "T";width 3.7593in;height 2.693in;depth
0pt;original-width 3.7118in;original-height 2.6507in;cropleft "0";croptop
"1";cropright "1";cropbottom "0";tempfilename
'TN9QSZ73.wmf';tempfile-properties "XPR";}}The converging lens will make the
effective focal length of the system shorter, so that it can form an image
where the retina actually is.

It would be convenient and very understandable if medical professionals
prescribed eye glasses by telling us what focal length we needed to correct
our vision. But that would not be the biology way! Eye glasses use a
different unit of measure to describe how they bend light. The unit is the 
\emph{diopter} and it is equal one over the focal length, but with the focal
length measured in meters. 
\begin{equation}
\text{diopter}=\frac{1}{f\left( \unit{m}\right) }
\end{equation}%
This measurement is called the \emph{power }of the lens. It provides the
same information as giving the focal length.

\subsubsection{Color Perception}

The eye detects different colors. The receptors called cones can detect red,
green, and blue light.\FRAME{dhF}{2.4587in}{1.0326in}{0pt}{}{}{Figure}{%
\special{language "Scientific Word";type "GRAPHIC";maintain-aspect-ratio
TRUE;display "USEDEF";valid_file "T";width 2.4587in;height 1.0326in;depth
0pt;original-width 4.3656in;original-height 1.8178in;cropleft "0";croptop
"1";cropright "1";cropbottom "0";tempfilename
'TN9QSZ74.wmf';tempfile-properties "XPR";}}The eye combines the red, green,
and blue response to allow us to perceive many different colors.

Most digital cameras also have red, green, and, blue pixels to provide color
to images. The detectors in digital cameras are often have much narrower
frequency bands than the eye. Likewise, television displays and monitors
have red, green, and blue pixels. By targeting the eye receptors, power need
not be wasted in creating light that is not detected well by the eye. But
the difference in bandwidth can cause problems in color mixing. Yellow
school busses (perceived as different amounts of green and red light) may
seem reddish or green if the detector bandwidths are chosen poorly.

The science of human visual perception of imagery is called \emph{image
science}. There are many applications for this field, from forensics to
intelligence gathering.

\subsection{The Camera}

The idea behind a camera is very simple. It is very much like an artificial
eye!\FRAME{dtbpF}{3.0926in}{1.8542in}{0pt}{}{}{Figure}{\special{language
"Scientific Word";type "GRAPHIC";maintain-aspect-ratio TRUE;display
"USEDEF";valid_file "T";width 3.0926in;height 1.8542in;depth
0pt;original-width 3.0485in;original-height 1.817in;cropleft "0";croptop
"1";cropright "1";cropbottom "0";tempfilename
'TN9QSZ75.wmf';tempfile-properties "XPR";}}

The camera has a lens, often a compound lens made of several individual
lenses, and a screen for projecting a real image created by the lens. We
haven't learned how to deal with multiple lenses yet, but it's not too bad
of an approximation to treat the whole camera lens system as a single
equivalent lens with an equivalent focal length. Photographers speak as
though this were true and it's good enough for them to do their work, but
not good enough to do optical engineering. But we can gain a lot of insight
into how cameras work with this simple model. So let's stick with this
approximation in this lecture.

Let's take an example camera. Say we are at a family reunion and wish to
take a picture of Aunt Agatha. Aunt Agatha is about $1.\,\allowbreak 5\unit{m%
}$ tall. She is standing about $5\unit{m}$ away. And suppose we have a $58%
\unit{mm}$ lens. 
\[
\begin{tabular}{l}
$h_{o}=1.5\unit{m}$ \\ 
$h_{i}=0.035\unit{m}$ \\ 
$d_{o}=5\unit{m}$ \\ 
$f=0.058\unit{m}$%
\end{tabular}%
\]%
Then to fit the image of Aunt Agatha on our $35\unit{mm}$ detector, we must
have

\begin{eqnarray*}
m &=&\frac{h_{i}}{h_{o}}=\frac{-0.035\unit{m}}{1.5\unit{m}} \\
&=&-2.\,\allowbreak 333\,3\times 10^{-2}
\end{eqnarray*}%
so our image is small and inverted.

The small size we wanted. But now we know from the negative sign that the
image in our cameras is upside down. A digital camera uses it's built-in
computer to turn the image right side up for us on the display on the back
of the camera.

Now let's find $d_{i}.$ 
\begin{eqnarray*}
d_{i} &=&\frac{fd_{o}}{d_{o}-f} \\
&=&5.\,\allowbreak 868\,1\times 10^{-2}\unit{m} \\
&=&58.681\unit{mm}
\end{eqnarray*}%
Our detector must be $58.681\unit{mm}$ from the lens.

Now suppose we want to photograph a $1000\unit{m}$ tower from $2\unit{km}$
away. Then%
\begin{eqnarray*}
m &=&-\frac{0.035\unit{m}}{1000\unit{m}} \\
&=&-3.\,\allowbreak 5\times 10^{-5}
\end{eqnarray*}%
and 
\begin{eqnarray*}
d_{i} &=&\frac{\left( 0.058\unit{m}\right) \left( 2000\unit{m}\right) }{2000%
\unit{m}-\left( 0.058\unit{m}\right) } \\
&=&5.\,\allowbreak 800\,2\times 10^{-2}\unit{m} \\
&=&58.002\unit{mm}
\end{eqnarray*}%
Notice that the image distance changed, but not by very much. This is why
you need a focus adjustment on the lens of a good camera. Objects far away
require a different $d_{i}$ value than objects that are close. Usually you
twist the lens housing to make this adjustment in manual focus mode. The
lens housing has a threaded screw system that increases or decreases $d_{i}$
as you twist. Cell phones and consumer cameras often have an motor that
makes this adjustment for you. In some cameras you may see the lens move
inward and outward as someone takes a picture.

There are several things that govern whether a picture will be good. When
you buy a quality manual lens, it will be marked in $f/\#$'s. The
specification of an automatic lens will also be given in terms of $f/\#$. To
help us buy such lenses, we should understand what the terminology means.

Most things we want to take a snapshot of are much farther than $58\unit{mm}$
from the camera. For such objects we can revisit the magnification.%
\[
m=-\frac{d_{i}}{d_{o}}
\]%
but from the thin lens formula%
\[
\frac{1}{d_{o}}+\frac{1}{d_{i}}=\frac{1}{f}
\]%
If $d_{o}\gg f$ then we can say that $1/d_{o}\approx 0$ and so $d_{i}\approx
f.$ Then%
\[
m\approx -\frac{f}{d_{o}}
\]%
and we see that the size of the image 
\[
h_{i}=mh_{o}
\]%
is directly proportional to the focal distance. If we change the focal
distance, we can change the size of the image. 
\[
h_{i}=\left( -\frac{f}{d_{o}}\right) h_{o}
\]%
This is how a zoom lens works. A zoom lens is a compound lens, and the focal
length is changed by increasing the distance between the component lenses
(more on compound lenses soon). You may have a camera with a telephoto
button and this button zooms in and out by moving one part of the complex
lens. Cell phones usually do not do this to zoom in. They just make a photo
out of the center of their photosensor array and because they just use the
middle, it looks like they zoomed in. This is a \emph{digital zoom}. But
let's go back to our standard non-zoom camera lens.\footnote{%
If you are a photographer, we would call this a prime lens.}

Remember we studied intensity%
\[
I=\frac{P}{A}
\]

Photographic film and digital focal plane arrays detect the intensity of
light that hits them. We can see that the area of our image depends on our
magnification, which depends on $d_{i}$ and for our distant objects it is
proportional to $f.$ The image area is proportional to $d_{i}^{2}\approx
f^{2}.$ So we can say that the area is proportional to $f^{2}.$ Then 
\[
I\propto \frac{P}{f^{2}}
\]%
The power entering the camera is proportional to the size of the aperture
(hole the light goes through). A bigger aperture lets in more light. A
smaller aperture lets in less light. If the camera has a circular opening,
this area is proportional to the square of the diameter of the opening, $%
D^{2}$ so%
\[
I\propto \frac{D^{2}}{f^{2}}
\]%
This ratio is useful because it gives us a relative measure of how much
intensity we get in terms of things we can easily know about our camera.
Good cameras have changeable aperture sizes, and many lenses have changeable
focal lengths. By using the combination of these two terms, we can ensure we
will get enough light (but not too much) when we take the picture.

It would be good to give this ratio a special name. But instead, we named
the ratio 
\begin{equation}
f/\#\equiv \frac{f}{D}
\end{equation}%
It is called the $f/\#$ (pronounced f-number, or f-stop) so our intensity is
proportional to one over the f-number. 
\begin{equation}
I\propto \frac{1}{\left( f/\#\right) ^{2}}
\end{equation}%
Good cameras have adjustable lens systems marked in $f/\#^{\prime }s$.
Typical values are $f/2.8,$ $f/4,$ $f/5.6,$ $f/8,$ $f/11,$ and $f/16.$
Notice that since $I$ is proportional to $\left( 1/f/\#\right) ^{2}$ these
common $f/\#$ values give an increase of intensity of a factor of $2$ each
time you change the $f/\#$ by one marked stop. That is because, say, 
\[
\frac{f/4}{f/2.8}=\frac{4}{2.8}=1.\,\allowbreak 428\,6\allowbreak =\sqrt{2}
\]
and the intensity goes as $1/\left( f/\#\right) ^{2}$ so the change in
intensity from $f/2.8$ to $f/4$ would be proportional to $\left( \sqrt{2}%
\right) ^{2}=2.$ It is the same between any two of our standard $f/\#$
values.%
\[
\begin{tabular}{|c|c|}
\hline
$f/\#$ change & Intensity change factor \\ \hline
$f/11$ to $f/8$ & $2$ \\ \hline
$f/8$ to $f/5.6$ & $2$ \\ \hline
$f/5.6$ to $f/4$ & $2$ \\ \hline
$f/4$ to $f/2.8$ & $2$ \\ \hline
\end{tabular}%
\]

This terminology is used for telescope design as well. The Hubble telescope
is an $f/24$ Ritchey-Chretien Cassegrainian system with a $2.4\unit{m}$
diameter aperture. The effective focal length is $57.6\unit{m}$.

\FRAME{dhF}{3.7593in}{2.3298in}{0pt}{}{}{Figure}{\special{language
"Scientific Word";type "GRAPHIC";maintain-aspect-ratio TRUE;display
"USEDEF";valid_file "T";width 3.7593in;height 2.3298in;depth
0pt;original-width 3.7118in;original-height 2.29in;cropleft "0";croptop
"1";cropright "1";cropbottom "0";tempfilename
'TN9QSZ76.wmf';tempfile-properties "XPR";}}

It is important to realize that electronic (and biological) sensors don't
react instantly to what we see. The intensity is 
\[
I=\frac{P}{A}=\frac{\Delta E}{\Delta tA}
\]%
so there is a time involved. The time it takes to collect enough light to
form an image on the sensor is called the \emph{exposure time}. 
\[
\Delta t=\frac{\Delta E}{IA}
\]%
So changing our $f/\#$ changes the needed exposure time by changing the
intensity. 
\[
\Delta t\propto \frac{\Delta E}{\frac{1}{\left( f/\#\right) ^{2}}A}=\frac{%
\Delta E}{A}\left( f/\#\right) ^{2}
\]%
How sensitive our camera sensor is also affects the exposure time. Modern
sensors have adjustable sensitivity. The photography world gives the three
letters ISO for the name for this detector sensitivity. There isn't a
standard for exactly what ISO setting gives what exposure. Different
manufacturers use slightly different numbers. But a change in ISO settings
usually are equivalent to one $f/\#$ change in exposure.

This is part of what a good photographer does in taking a picture. The
photographer will adjust the $f/\#$ and the exposure time and ISO to get a
photograph that is not over exposed (too light) or underexposed (to dark).

Now that we know a little about cameras, let's consider some details, like
how to get an equivalent focal distance for a combination of lenses. And
then let's look at some real optical systems.

\section{Optical Systems 3.2.8, 3.3 1}

So far we have looked at single lenses and how they work. But we hinted that
actual camera lenses and even our eyes are have lens systems that are made
of more than one lens. We should look at how lens systems work. First,
though, let's go back to the simple magnifier to define a more meaningful
magnification measure for using optical systems.

%TCIMACRO{%
%\TeXButton{Fundamental Concepts}{\vspace{0.10in}\hspace{-0.37in}{\Large Fundamental Concepts\vspace{0.2in}}}}%
%BeginExpansion
\vspace{0.10in}\hspace{-0.37in}{\Large Fundamental Concepts\vspace{0.2in}}%
%EndExpansion

\begin{itemize}
\item Angular magnification $M=\frac{\theta }{\theta _{o}}$ describes how
much bigger the image looks with an optical system compared to with just
your eye.

\item The magnification of a two-lens system is just the product of the
magnifications of the individual lenses $m_{\text{combined}}=m_{1}m_{2}$

\item If the two lenses in a two-lens system are placed so the distance
between them is essentially zero, then the focal length of the two-lens
system is given by 
\[
\frac{1}{f}=\frac{1}{f_{1}}+\frac{1}{f_{2}}
\]%
where $f_{1}$ and $f_{2}$ are the focal lengths of the individual lenses

\item If the two lenses in a two-lens system are not placed so the distance
between them is essentially zero, the previous equation is not valid!

\item Telescopes and Microscopes are double lens systems
\end{itemize}

\subsection{Angular Magnification}

We already encountered the simple magnifier when we studied ray diagrams.
But by this point in our study of optics you are probably wondering about
our definition of magnification. If you are an Idahoan and are out hunting,
when you look through your binoculars or scope you don't want to know how
big the image is compared the actual size of the dear. No, you want to know
if you can see the deer better with the scope than you could with just your
eyes. The magnification on your scope doesn't compare the image size to the
object size. It is a comparison between what you see with two different
optical systems. One is just your eyes, the other is your eyes and the scope
working together.

Let's use the simple magnifier that we already know about (rather than a
complex scope) to define a new kind of magnification that does this
comparison between two optical systems. We can use what we know about rays
that are easy to draw, and about how our eyes work to describe both optical
systems.\FRAME{dtbpF}{3.3278in}{0.8717in}{0pt}{}{}{Figure}{\special{language
"Scientific Word";type "GRAPHIC";maintain-aspect-ratio TRUE;display
"USEDEF";valid_file "T";width 3.3278in;height 0.8717in;depth
0pt;original-width 5.2174in;original-height 1.3456in;cropleft "0";croptop
"1";cropright "1";cropbottom "0";tempfilename
'TN9QSZ6Q.wmf';tempfile-properties "XPR";}}

We usually use three principle rays, but for this analysis, let's just use
one for each side of an object. Since the image needs to be real and on the
retina, we can see where these singular rays strike the retina and
understand the size of the image on our eye light detection system. Let's
choose to draw the rays that go straight through the middle of the lens of
the eye because they are the easiest ones to draw. They don't bend, they
just go straight through. If we pick a ray from the top of our object that
goes through the center of the lens, that ray won't seem to change direction
at all. It will hit the retina to form the top of the image of the object.

We can do the same for the bottom of the object. Then we can see from the
next figure that the angle $\theta _{o}$ subtends\footnote{%
The angle that \textquotedblleft subtends\textquotedblright\ and object is
the one who's bounding lines just hit either side of the object.} both the
object and the image of the object. If you think about the previous figure,
you will see that if the angle were to increase, so would the size of the
image on the retina. And that would mean that we would see a bigger version
of the image and would be able to see more detail. We can increase this
angle by moving the object closer to our eyes. In the next figure the top is
just a repeat of the last figure, but the bottom is with the object moved
closer.\FRAME{dtbpF}{3.6089in}{1.9484in}{0pt}{}{}{Figure}{\special{language
"Scientific Word";type "GRAPHIC";maintain-aspect-ratio TRUE;display
"USEDEF";valid_file "T";width 3.6089in;height 1.9484in;depth
0pt;original-width 3.5613in;original-height 1.9104in;cropleft "0";croptop
"1";cropright "1";cropbottom "0";tempfilename
'TN9QSZ6R.wmf';tempfile-properties "XPR";}}

You can see that the image covers more of the retina. And so we see a bigger
version of the object. If we move our object too close we will reach the
limit of the eye for focusing. \FRAME{dtbpF}{3.8882in}{1.0758in}{0pt}{}{}{%
Figure}{\special{language "Scientific Word";type
"GRAPHIC";maintain-aspect-ratio TRUE;display "USEDEF";valid_file "T";width
3.8882in;height 1.0758in;depth 0pt;original-width 4.8585in;original-height
1.324in;cropleft "0";croptop "1";cropright "1";cropbottom "0";tempfilename
'TN9QSZ77.wmf';tempfile-properties "XPR";}}If we move the object any closer
than this limit, it will appear blurry. Remember, we call this position, the
closest point where we can place an object and still bring it into focus
with our eye, the \emph{near point.} Thus the maximum value of $\theta _{o}$
where the image won't be blurry will be at the near point.

But suppose we want to see this object in more detail. That is, we want to
spread the image of the object over more of our retina so our retinal pixels
do a better job of showing the details. We can use a magnifying glass to
spread the light over more of our retina. Think, from our study of
magnifying glasses we know that the magnifying glass spreads out the light. 
\FRAME{dtbpF}{3.2348in}{2.1092in}{0pt}{}{}{Figure}{\special{language
"Scientific Word";type "GRAPHIC";maintain-aspect-ratio TRUE;display
"USEDEF";valid_file "T";width 3.2348in;height 2.1092in;depth
0pt;original-width 4.4198in;original-height 2.872in;cropleft "0";croptop
"1";cropright "1";cropbottom "0";tempfilename
'TN9QSZ78.wmf';tempfile-properties "XPR";}}It doesn't form a real image. So
if we look from the left hand side toward the magnifying glass in the figure
we see teh light from the lamp get spread out. And that is what we want! We
want to spread out the light so it covers more of our retina once it enters
our eye.

This worked for our lamp diagram. But let's go back to a normal ray diagram
and draw two principle rays for the top of the bug object. \FRAME{dtbpF}{%
5.4846in}{2.565in}{0pt}{}{}{Figure}{\special{language "Scientific Word";type
"GRAPHIC";maintain-aspect-ratio TRUE;display "USEDEF";valid_file "T";width
5.4846in;height 2.565in;depth 0pt;original-width 5.4267in;original-height
2.5235in;cropleft "0";croptop "1";cropright "1";cropbottom "0";tempfilename
'TN9QSZ79.wmf';tempfile-properties "XPR";}} Sure enough, the two rays spread
out and never meet to form a real image, so we have a virtual image of the
bug. To achieve this we placed the object closer to the magnifying glass
than the focal distance $(d_{o}<f)$, then we have a virtual image with
magnification 
\begin{equation}
m=\frac{-d_{i}}{d_{o}}
\end{equation}%
which is larger than one and positive (because $d_{i}$ is negative by sign
convention). This is what we have studied so far, but you probably guessed
there was more to our virtual images than just this. After all, we do see
the image with our eye. So, we need to think about how our eye sees this
virtual image.

In the next figure I\ have departed from our three principle rays. We know
they were just representative rays, there are millions of actual rays. I
drew one ray that is one of the principle rays (the one that comes in
parallel to the axis) but I also drew a non-principle ray. Still, we see
that these two rays are spreading after the magnifying glass and we can
trace them backwards to find our virtual image. They never meet to form a
real image on their own. But now we introduce another optical system, our
eye! \FRAME{dtbpF}{4.8153in}{2.4448in}{0pt}{}{}{Figure}{\special{language
"Scientific Word";type "GRAPHIC";maintain-aspect-ratio TRUE;display
"USEDEF";valid_file "T";width 4.8153in;height 2.4448in;depth
0pt;original-width 5.1552in;original-height 2.6031in;cropleft "0";croptop
"1";cropright "1";cropbottom "0";tempfilename
'TN9QSZ7A.wmf';tempfile-properties "XPR";}}The lens of our eye is another
positive lens and it can refract the rays back so they meet. Our eye
refracts the diverging rays to form a real image on our retina. So in the
end our virtual image from the lens became a real image on our retina, but
it took an additional lens to make this happen. In the last figure I just
drew rays for the top of the bug, but we would need another set of rays for
the bottom of the bug to see how big the real image at the retina became.
But you can see from the figure that the rays that do enter our eye are
diverging before they enter and they appear to have come from the virtual
image. Then the lens of our eye brings the rays back to gether to form a
real image on the retina. This is how we see virtual images.

Back to magnification. What we really want to know is how much bigger the
image looks with the lens than it did without the lens. By looking at what
happens to the rays when they enter our eye, we now can see why the image
looks bigger. Let's go back to just a ray for the top of the bug and a ray
for the bottom of the bug and align our eye so these rays go through the
cornea as the undeviated rays. Then our diagram is simpler and we can see
why the magnifying lens makes a bigger image on our retina. \FRAME{dtbpF}{%
3.8337in}{3.6028in}{0in}{}{}{Figure}{\special{language "Scientific
Word";type "GRAPHIC";maintain-aspect-ratio TRUE;display "USEDEF";valid_file
"T";width 3.8337in;height 3.6028in;depth 0in;original-width
3.7853in;original-height 3.5544in;cropleft "0";croptop "1";cropright
"1";cropbottom "0";tempfilename 'TN9QSZ7B.wmf';tempfile-properties "XPR";}}

The magnifying glass has bent the light, and the bent rays make a bigger
angle, $\theta _{i}>\theta _{o}$ so the image on our retina is bigger. Since
the image fills more of our retina, we perceive the image as being bigger.
We need a mathematical formula that will tell us how much bigger the image
on our retina will be. We need to be careful with two parts of our optical
system. We have two images, one virtual and one real. Let's call the virtual
image $I_{1}$ with height $h_{i1}$ and the real image $I_{2}$ with height $%
h_{i2}.$

Notice from the figure that if we look inside our eye for the case of no
extra lens 
\[
\tan \theta _{o}\approx \frac{h_{i,eye}}{d}=\frac{h_{i}}{d}
\]%
and for the case with the lens%
\[
\tan \theta _{i}\approx \frac{h_{i,lens+eye}}{d}=\frac{h_{i2}}{d}
\]%
so 
\begin{eqnarray*}
h_{i} &=&h_{i,eye}\approx d\tan \theta _{o} \\
h_{i2} &=&h_{i,lens+eye}\approx d\tan \theta _{i}
\end{eqnarray*}%
Then if we compare the new, larger image on the retina formed with the
lens-eye system to the one formed with just the eye, we get 
\[
\frac{h_{i2}}{h_{i}}\approx \frac{d\tan \theta _{i}}{d\tan \theta _{o}}=%
\frac{\tan \theta _{i}}{\tan \theta _{o}}
\]%
and if we once again use the small angle approximation 
\[
\frac{h_{i2}}{h_{i}}\approx \frac{\theta _{i}}{\theta _{o}}
\]%
This would tell us how much bigger our object looks when viewed with the
magnifying glass compared to how it looked without the magnifying glass.
This is just what we want! Let's give this a new symbol 
\begin{equation}
M=\frac{\theta _{i}}{\theta _{o}}
\end{equation}

Remember, this is really different than the magnification we have studied
before. The magnification we have been using compared the size of the image
with the size of the object. The magnification $M$ compares the size of an
image made with just our eye to the size of the image made by our eye and
another lens system. We call $M$ the \emph{angular magnification} so we can
keep these two types of magnification straight.

Then, the angular magnification compares how big the object seems to be with
and without an additional lens or lens system. It is really a comparison
between the size of the area the diverging set of rays make on the retina
formed with just our eye, and the one formed with the magnifier.

We now know what angular magnification is, lets find an equation for the
angular magnification for our simple magnifier. \FRAME{dtbpF}{4.2929in}{%
2.9049in}{0pt}{}{}{Figure}{\special{language "Scientific Word";type
"GRAPHIC";maintain-aspect-ratio TRUE;display "USEDEF";valid_file "T";width
4.2929in;height 2.9049in;depth 0pt;original-width 5.0021in;original-height
3.3754in;cropleft "0";croptop "1";cropright "1";cropbottom "0";tempfilename
'TN9QSZ7C.wmf';tempfile-properties "XPR";}}Let's say the image is not right
at the near point, and even make it harder by saying the lens is a distance $%
\ell $ away from the eye. This is really how we use a simple magnifier. Then%
\[
L=\ell +d_{i1}
\]%
is how far the virtual image is from the eye. And we can go back to
trigonometry. 
\[
\tan \theta _{i}\approx \frac{h_{i1}}{L}\approx \theta _{i}
\]%
And it is still true that without the lens to get the biggest view we need
to put the object at our near point. Let's choose this for our case without
the magnifier. 
\[
\tan \theta _{o}=\frac{h_{o}}{25\unit{cm}}\approx \theta _{o}
\]%
Our angular magnification is 
\[
M=\frac{\theta _{i}}{\theta _{o}}\approx \frac{\frac{h_{i1}}{L}}{\frac{h_{o}%
}{25\unit{cm}}}=\frac{h_{i1}}{L}\frac{25\unit{cm}}{h_{o}}
\]%
We can use our regular magnification 
\[
m=\frac{h_{i1}}{h_{o}}=-\frac{d_{i1}}{d_{o1}}
\]%
to write $M$ as 
\[
M\approx -\frac{d_{i1}}{d_{o1}}\frac{25\unit{cm}}{L}
\]%
and we can eliminate $d_{o1}$ by solving for it and then substituting%
\[
\frac{1}{d_{o1}}+\frac{1}{d_{i1}}=\frac{1}{f_{1}}
\]%
so 
\[
\frac{1}{d_{o1}}=\frac{1}{f_{1}}-\frac{1}{d_{i1}}
\]%
and we actually want $1/d_{o1}$ so let's substitute now%
\[
M\approx -d_{i1}\left( \frac{1}{f_{1}}-\frac{1}{d_{i1}}\right) \frac{25\unit{%
cm}}{L}
\]%
Let's bring the $-d_{i1}$ inside%
\begin{eqnarray*}
M &\approx &\left( \frac{-d_{i1}}{f_{1}}-\frac{-d_{i1}}{d_{i1}}\right) \frac{%
25\unit{cm}}{L} \\
&=&\left( \frac{-d_{i1}}{f_{1}}+1\right) \frac{25\unit{cm}}{L} \\
&=&\left( 1-\frac{d_{i1}}{f_{1}}\right) \frac{25\unit{cm}}{L}
\end{eqnarray*}

Then recall that $-d_{i1}=L-\ell $ (the minus sign is because the image is
virtual so by sign convention $d_{i1}$ is negative) so 
\[
M\approx \left( 1+\frac{L-\ell }{f_{1}}\right) \frac{25\unit{cm}}{L}
\]%
This is our angular magnification for a simple magnifier.

From this we can get the maximum value for our angular magnification for a
simple magnifier. We just let $L=25\unit{cm}$ and we make $\ell \approx 0$
then the virtual image is as close as we can get it and have it not be
blurry. This gives 
\begin{eqnarray*}
M_{\max } &\approx &\left( 1+\frac{25\unit{cm}-0}{f_{1}}\right) \frac{25%
\unit{cm}}{25\unit{cm}} \\
&\approx &\left( 1+\frac{25\unit{cm}}{f_{1}}\right)
\end{eqnarray*}%
This is our maximum magnification (try it, put the magnifying glass right
next to your eye and see how big the image looks. then try other locations
for the lens).

We could also get a minimum angular magnification for our simple magnifier.
We can think of what we know about ray diagrams. If we put the object right
at the focal distance so $d_{o1}=f_{1}$ then the image would be at infinity
because the rays from the object would be bent into parallel rays by the
lens. And looking at our diagram that would mean $L=\infty .$ Of course that
is too far, we wouldn't actually see a virtual image. But suppose we only
let $L\rightarrow \infty $ (so $d_{o1}\rightarrow f_{1}$ but doesn't quite
equal $f_{1}$). Then we should get the minimum magnification 
\begin{eqnarray*}
M_{\min } &=&\underset{L\rightarrow \infty }{\lim }\left( 1+\frac{L-\ell }{%
f_{1}}\right) \frac{25\unit{cm}}{L} \\
&=&\underset{L\rightarrow \infty }{\lim }\left( 25\unit{cm}\left( \frac{1}{L}%
+\frac{L-\ell }{f_{1}L}\right) \right)
\end{eqnarray*}%
and as $L\rightarrow \infty $ little $\ell $ becomes unimportant and $1/L$
becomes zero so 
\begin{eqnarray*}
M_{\min } &=&\underset{L\rightarrow \infty }{\lim }\left( \left( 1+\frac{%
L-\ell }{f_{1}}\right) \frac{25\unit{cm}}{L}\right) \\
&=&\underset{L\rightarrow \infty }{\lim }\left( 25\unit{cm}\left( \frac{1}{L}%
+\frac{L}{Lf_{1}}\right) \right) \\
&=&25\unit{cm}\left( 0+\frac{1}{f_{1}}\right) \\
&=&\frac{25\unit{cm}}{f_{1}}
\end{eqnarray*}%
There were a lot of approximations, but still, we can see that we can have a
magnification approximately in the range 
\[
\frac{25\unit{cm}}{f_{1}}<M<1+\frac{25\unit{cm}}{f_{1}}
\]

We use the idea of a simple magnifier in combination with other lenses to
make the magnification happen in telescopes, microscopes, and other
instruments that magnify. So all of these systems can be described using the
idea of an angular magnification.

Angular magnification is important. Many optical systems that we buy are
really designed to make things seem bigger (or smaller) than they seem with
just our eyes. Things like telescopes and microscopes and binoculars.

\subsection{Combinations of Lenses}

Our eyes and cameras have more than one lens. in using a magnifying glass
technically we have more than one lens because our eye is involved. We found
that to correct chromatic aberration we used two lenses. Together they are
called an \textquotedblleft achromat\textquotedblright\ and all good cameras
use achromats to fix chromatic aberration. We can do something similar to
correct for spherical aberration. But in each case, we are combining two
lenses (or even more!). Microscopes and telescopes are also systems that
have a minimum of two lenses. How do we predict what the lens system will
do? It turns out that we have all we need to know to combine lenses already.

To combine lenses, we do the same thing we did for the two surfaces of a
thin lens or for our simple magnifier case. We form the image from the first
lens as though the second lens is not there. Then we use the image from the
first lens as the object for the second lens. Suppose we take two lenses of
focal lengths $f_{1}$ and $f_{2}$ and place them a distance $d$ apart. 
\FRAME{dtbpF}{4.286in}{1.7876in}{0pt}{}{}{Figure}{\special{language
"Scientific Word";type "GRAPHIC";maintain-aspect-ratio TRUE;display
"USEDEF";valid_file "T";width 4.286in;height 1.7876in;depth
0pt;original-width 4.235in;original-height 1.7504in;cropleft "0";croptop
"1";cropright "1";cropbottom "0";tempfilename
'TN9QSZ7D.wmf';tempfile-properties "XPR";}}

Because this system would use a magnified image of our object as the object
for lens $2,$ the final magnification is the product of the two lens
magnifications

\begin{equation}
m_{\text{combined}}=m_{1}m_{2}
\end{equation}

For the first lens we have%
\begin{equation}
\frac{1}{d_{o1}}+\frac{1}{d_{i1}}=\frac{1}{f_{1}}
\end{equation}%
where $d_{i1}$ is our first lens image distance. We can solve for $d_{i1}$%
\begin{equation}
d_{i1}=\frac{d_{o1}f_{1}}{d_{o1}-f_{1}}  \label{q}
\end{equation}%
This first image from the first lens will be the object for the second lens

We take as the second object distance as 
\[
d_{o2}=d-d_{i1}
\]%
and we use the lens formula again.%
\[
\frac{1}{d_{o2}}+\frac{1}{d_{i2}}=\frac{1}{f_{2}}
\]%
and again find the image distance%
\[
d_{i2}=\frac{d_{o2}f_{2}}{d_{o2}-f_{2}}
\]%
but we can use our value of $d_{o2}=d-d_{i1}$ to give us%
\begin{eqnarray*}
d_{i2} &=&\frac{\left( d-d_{i1}\right) f_{2}}{\left( d-d_{i1}\right) -f_{2}}
\\
&=&\frac{\left( d-d_{i1}\right) f_{2}}{d-d_{i1}-f_{2}}
\end{eqnarray*}%
We have and expression relating the image distances, $d,$ $d_{i1,}$ $d_{i2}$
and $f_{2}.$ But we would really like to have an expression that relates $%
d_{o1}$ and $d_{i2}.$ That is, we want a relationship that by knowing how
far away the object is from the lens system we get out how far the final
image is from the lens system. Then we can treat the lens system like a
single lens (most of the time).

Lets use our expression for 
\[
d_{i1}=\frac{d_{o1}f_{1}}{d_{o1}-f_{1}}
\]%
and substitute it into our expression for $d_{i2}$%
\[
d_{i2}=\frac{\left( d-\frac{d_{o1}f_{1}}{d_{o1}-f_{1}}\right) f_{2}}{d-\frac{%
d_{o1}f_{1}}{d_{o1}-f_{1}}-f_{2}}
\]%
This looks messy, but it does the job. Knowing $d_{o1}$ we can find $d_{i2}$%
!. We can do some simplification%
\begin{equation}
d_{i2}=\frac{df_{2}-\frac{d_{o1}f_{1}f_{2}}{d_{o1}-f_{1}}}{d-f_{2}-\frac{%
d_{o1}f_{1}}{d_{o1}-f_{1}}}  \label{TwoLensesSeparated}
\end{equation}%
Well, it is still messy, but we have achieved our goal! We have $d_{i2}$ in
terms of the focal lengths $f_{1}$ and $f_{2}$, the distances $d,$ and $%
d_{o1}.$

Let's think of our achromat lens system where the lenses are glued together.
They are right next to each other. The distance between our lenses is $d.$
Suppose we let $d\rightarrow 0.$ That would mean the two lenses are right
next to each other, like our achromat. Then%
\[
d_{i2}=\frac{-\frac{d_{o1}f_{1}f_{2}}{d_{o1}-f_{1}}}{-f_{2}-\frac{d_{o1}f_{1}%
}{d_{o1}-f_{1}}}
\]%
which hardly looks better. But let's rearrange some with a common
denominator in the bottom of our expression 
\begin{eqnarray*}
d_{i2} &=&\frac{\frac{d_{o1}f_{1}f_{2}}{d_{o1}-f_{1}}}{\frac{f_{2}\left(
d_{o1}-f_{1}\right) }{d_{o1}-f_{1}}+\frac{d_{o1}f_{1}}{d_{o1}-f_{1}}} \\
&=&\frac{d_{o1}f_{1}f_{2}}{f_{2}d_{o1}-f_{2}f_{1}+d_{o1}f_{1}} \\
&=&\frac{d_{o1}f_{1}f_{2}}{d_{o1}\left( f_{2}+f_{1}\right) -f_{2}f_{1}}
\end{eqnarray*}%
So%
\[
d_{i2}=\frac{d_{o1}f_{1}f_{2}}{d_{o1}\left( f_{2}+f_{1}\right) -f_{2}f_{1}}
\]%
Lets undo the math that brought us $d_{i2}$ in the first place. I want to
make this look like the thin lens formula, but with $d_{i2}$ as the image
distance and $d_{o1}$ as the object distance. 
\begin{eqnarray*}
\frac{1}{d_{i2}} &=&\frac{d_{o1}\left( f_{2}+f_{1}\right) -f_{2}f_{1}}{%
d_{o1}f_{1}f_{2}} \\
&=&\frac{d_{o1}\left( f_{2}+f_{1}\right) }{d_{o1}f_{1}f_{2}}-\frac{f_{2}f_{1}%
}{d_{o1}f_{1}f_{2}} \\
&=&\frac{\left( f_{2}+f_{1}\right) }{f_{1}f_{2}}-\frac{1}{d_{o1}}
\end{eqnarray*}%
or%
\[
\frac{1}{d_{i2}}+\frac{1}{d_{o1}}=\frac{\left( f_{2}+f_{1}\right) }{%
f_{1}f_{2}}
\]%
Which looks very like the lens formula with 
\[
\frac{1}{f}=\frac{\left( f_{2}+f_{1}\right) }{f_{1}f_{2}}
\]%
If we unwind this expression for $f$, we find%
\begin{equation}
\frac{1}{f}=\frac{f_{2}}{f_{1}f_{2}}+\frac{f_{1}}{f_{1}f_{2}}  \nonumber
\end{equation}%
\begin{equation}
\frac{1}{f}=\frac{1}{f_{1}}+\frac{1}{f_{2}}
\end{equation}

We see that the two lenses are equivalent to a single lens with focal length 
$f$ as long as they are close together.

Of course, we had to place our lenses right next to each other for this to
work. This is not the case for a telescope or microscope. We should look at
such a case. There is no need for more math. We can go back to equation (\ref%
{TwoLensesSeparated}) for the case of two lenses separated by a distance $d$%
. 
\[
d_{i2}=\frac{df_{2}-\frac{d_{o1}f_{1}f_{2}}{d_{o1}-f_{1}}}{d-f_{2}-\frac{%
d_{o1}f_{1}}{d_{o1}-f_{1}}}
\]

Let's look at a case using ray diagrams. For this case, let's take two
lenses, and let's have the first lens make a real image. Once again, let's
have that image be the object for the second lens. But this time, let's move
the second lens so that the image from the first lens (object for the second
lens) is closer to the second lens than $f_{2}.$ If that is the case, the
second lens works like a magnifier. The final image is enlarged.

\FRAME{dtbpF}{3.9409in}{2.5746in}{0pt}{}{}{Figure}{\special{language
"Scientific Word";type "GRAPHIC";maintain-aspect-ratio TRUE;display
"USEDEF";valid_file "T";width 3.9409in;height 2.5746in;depth
0pt;original-width 3.8917in;original-height 2.5339in;cropleft "0";croptop
"1";cropright "1";cropbottom "0";tempfilename
'TN9QSZ7E.wmf';tempfile-properties "XPR";}}This is the basis of both
microscopes and telescopes!

\subsection{The Microscope}

Microscopes are designed to allow us to see very small things. They go
beyond a simple magnifier. But it turns out we don't have to go very far
beyond a magnifying glass. To see things that are very small, we add another
lens to our simple magnifier. We will place this lens near the object and
call this lens the \emph{objective} because it is next to the object. We
will keep a simple magnifier and place it near the eye. This lens will be
called the \emph{eyepiece} because it is near your eye.

The objective will have a very short focal length. The eyepiece will have a
longer focal length (a few centimeters).\FRAME{dtbpF}{3.9409in}{2.8383in}{0in%
}{}{}{Figure}{\special{language "Scientific Word";type
"GRAPHIC";maintain-aspect-ratio TRUE;display "USEDEF";valid_file "T";width
3.9409in;height 2.8383in;depth 0in;original-width 3.8917in;original-height
2.7951in;cropleft "0";croptop "1";cropright "1";cropbottom "0";tempfilename
'TN9QSZ7F.wmf';tempfile-properties "XPR";}}

We separate the lenses by a distance $L$ where 
\begin{eqnarray*}
L &>&f_{o} \\
L &>&f_{e}
\end{eqnarray*}

We place the object \emph{just outside} the focal point of the objective
lens. The image formed by the objective lens is then real and inverted. Then
(and here is the trick we just learned!) we use this real image formed by
the objective lens as the new object for the eyepiece. For a real image,
there is really light there. So a real image can act like light bouncing off
of an object. We can use the real image as an object for the eyepiece.

The image formed by the eyepiece is upright and virtual, but it looks upside
down because the object for the eyepiece (the image made by the objective)
is upside down.

To say it once more, what we are doing is forming a real image using the
objective lens and then using the eyepiece lens as a simple magnifier to
make a larger virtual image of the real image made by the objective lens. We
could find $d_{ie}$ using equation (\ref{TwoLensesSeparated}).

But this is a microscope. We want to know how much bigger our final image
is. And of course, we really want to know the magnification of the system.
It's not just the eyepiece that will magnify. We know what the eyepiece will
do because it is being used as just a simple magnifier. Recall that
magnifications are a factors. Think from our basic equation 
\[
m=\frac{h_{i}}{h_{o}}
\]%
tells us that 
\[
h_{i}=mh_{o}
\]%
or in words, $m$ is the factor by which $h_{i}$ is bigger (or smaller) than $%
h_{o}.$ So magnifications are factors. For example, $d_{i}$ could be $10$
time bigger than $d_{o}$. And we have two lenses. Each lens will have a
magnification of it's object. But how do we combine the magnifications of
the two lenses? We just multiply the two magnifications together! The
objective lens makes an image that is a factor of $m_{o}$ times bigger (or
smaller) than the object. Then the eyepiece makes an image of the first
image that is $m_{e}$ times bigger than the first image. So the final image
is $m_{o}m_{e}$ times bigger than the object.

We should ask if this works for angular magnification. To find the overall
system angular magnification we start with the magnification of the eyepiece 
$M_{e},$ say, $20$ times bigger than it would look with just our eye. So
whatever the object, the eyepiece gives us a virtual image of that object
that is $M_{e}$ times bigger than it would look without the eyepiece. But
the object in our case is the real image formed by the objective lens. And
this real image may be bigger than the actual object. If you think about it,
it makes sense that if the objective makes the object look $10$ times
bigger, and the eyepiece makes the image look $20$ times bigger than if you
looked at it with your eye, then the whole system makes it look $200$ times
bigger. The combined magnification for the two-lens system is 
\begin{equation}
M_{system}=m_{o}M_{e}
\end{equation}%
The minimum magnification of the eye piece will be roughly 
\[
M_{e_{\min }}\approx \frac{25\unit{cm}}{f_{e}}
\]%
and the maximum will be 
\[
M_{e_{\max }}=1+\frac{25\unit{cm}}{f_{e}}
\]%
But remember the object for the eyepiece is the image from the first lens.
And that image is larger than the object by an amount \ 
\[
m_{o}=\frac{-d_{io}}{d_{oo}}
\]%
where $d_{io}$ is the image distance for the objective lens and $d_{oo}$ is
the object distance for the objective lens. And let's estimate how big the
magnification due to the first lens will be. Because for a microscope $%
d_{oo}\approx f_{o}.$ and $d_{io}\approx L$ (very roughly) 
\[
m_{o}=\frac{-d_{io}}{d_{oo}}\approx -\frac{L}{f_{o}}
\]

The combined magnification for the two-lens system is about 
\begin{equation}
M_{system}=m_{o}M_{e}\approx -\frac{L}{f_{o}}\frac{25\unit{cm}}{f_{e}}
\end{equation}%
this is the minimum magnification (because we used the minimum magnification
formula for the eyepiece).

This gives us a rough idea of how a microscope works, let's try telescopes
next.

\subsection{Telescopes}

There are two main types of telescopes \emph{refracting} and \emph{reflecting%
}. We will study refracting telescopes first.

\textbf{Refracting Telescopes}

Like the microscope, we combine two lenses and call one the objective and
the other the eyepiece. The eyepiece again plays the role of a simple
magnifier, magnifying the image produced by the objective.\FRAME{dtbpF}{%
5.1084in}{2.9516in}{0pt}{}{}{Figure}{\special{language "Scientific
Word";type "GRAPHIC";maintain-aspect-ratio TRUE;display "USEDEF";valid_file
"T";width 5.1084in;height 2.9516in;depth 0pt;original-width
4.4996in;original-height 2.5884in;cropleft "0";croptop "1";cropright
"1";cropbottom "0";tempfilename 'TN9QSZ7G.wmf';tempfile-properties "XPR";}}%
We again form a real, inverted image with the objective. We are now looking
at distant objects. This means $d_{oo}$ is very large. and using%
\[
d_{oi}=\frac{d_{oo}f_{o}}{d_{oo}-f_{o}}\approx \frac{d_{oo}f_{o}}{d_{oo}}%
=f_{o}
\]%
we can see that the image distance $d_{io}\approx f_{o}.$ Once again, we use
the image from the objective as the object for the eyepiece. The eye piece
forms an upright virtual image (that looks inverted because the object for
the eyepiece is the image from the objective, and the real image from the
objective is inverted). Once again we could find $d_{ie}$ using equation (%
\ref{TwoLensesSeparated}).

Once again we want the magnification. The largest magnification is when the
rays exit the eyepiece parallel to the principal axis. Then the image from
the eyepiece is formed at near infinity (but it is very big, so it is easy
to see). This gives a lens separation of $f_{o}+f_{e}$ which will be roughly
the length of the telescope tube.

The angular magnification will be 
\begin{equation}
M=\frac{\theta _{i}}{\theta _{o}}
\end{equation}%
where $\theta _{o}$ is the angle subtended by the object at the objective
(see figure above) it's what you see if you stand in front of the telescope
and just use your eye. And $\theta _{i}$ is subtended by the final image at
the viewer's eye. Consider the height of the final image, $h_{o},$ is very
large. And let's call the height of the image made by the objective lens $%
h_{io}.$ We see from the figure that 
\begin{equation}
\tan \theta _{o}\approx -\frac{h_{io}}{f_{o}}
\end{equation}%
and with the distance from the objective lens to the object, $d_{oo},$ very
large we can use small angles. 
\begin{equation}
\theta _{o}\approx -\frac{h_{io}}{f_{o}}
\end{equation}

The angle $\theta _{i}$ will be the angle formed by rays bent by the lens of
the eyepiece. This angle can be found by following a ray that leaves the tip
of the first image and goes through the eyepiece and into the person's eye.
The easiest ray to chose is the one that goes through the middle of the
eyepiece. Since we placed the eyepiece so $d_{eo}\approx f_{e}$ 
\begin{equation}
\tan \theta _{i}\approx \frac{h_{io}}{f_{e}}\approx \theta _{i}
\end{equation}%
so

The magnification is nearly%
\begin{equation}
M=\frac{\theta _{i}}{\theta _{o}}\approx \frac{\frac{h_{io}}{f_{e}}}{-\frac{%
h_{io}}{f_{o}}}\approx -\frac{f_{o}}{f_{e}}
\end{equation}

\textbf{Reflecting Telescopes}

Reflecting telescopes use a series of mirrors to replace the objective lens.
Usually, there is an eyepiece that is refractive (although there need not
be, radio frequency telescopes rarely have refractive pieces).\FRAME{dhF}{%
3.6348in}{2.7354in}{0pt}{}{}{Figure}{\special{language "Scientific
Word";type "GRAPHIC";maintain-aspect-ratio TRUE;display "USEDEF";valid_file
"T";width 3.6348in;height 2.7354in;depth 0pt;original-width
3.5864in;original-height 2.693in;cropleft "0";croptop "1";cropright
"1";cropbottom "0";tempfilename 'TN9QSZ7H.wmf';tempfile-properties "XPR";}}

There are two reasons to build reflective telescopes. The first is that
reflective optics do not suffer from chromatic aberration. The second is
that large mirrors are much easier to make and mount than refractive optics.
The Keck Observatory in Hawaii has a $10\unit{m}$ reflective system. The
largest refractive system is a $1\unit{m}$ system. Mechanical engineering
students will appreciate that putting giant $1\unit{m}$ diameter hunk of
glass on the end of a hollow tube and then turning the tube to point in
different directions would cause the tube and lens mass to oscillate. This
would make the imagery from the telescope very bad unless you wait for the
\textquotedblleft ringing\textquotedblright\ oscillations to die out due to
damping. Putting the heavy mass on the bottom of the tube is much more
stable.

The Hubble telescope has a $2.5\unit{m}$ aperture and James Webb has a $6%
\unit{m}$. Reflective telescopes are just easier to build.

The telescope pictured in the figure is a Newtonian, named after Newton, who
designed this type of telescope. Many other designs exist. Popular designs
for space applications include the Cassegrain telescope. The rough design of
a reflective telescope can be worked out using refractive pieces, then the
rough details of the reflective optics can be formed.

\end{document}
